\documentclass[11pt,oneside]{book}
\usepackage[margin=1.2in]{geometry}
\usepackage[toc,page]{appendix}
\usepackage{graphicx}
\usepackage{natbib}
\usepackage{lipsum}
\usepackage{caption}
\usepackage[T1]{fontenc}
\usepackage{titlesec, blindtext, color}
\usepackage{xcolor,tikz}
\usetikzlibrary{patterns}
\usetikzlibrary{decorations.markings}
\usepackage{amsmath,amssymb,amsthm,mathrsfs,amsfonts,xfrac,pifont,bbold,physics}
\usepackage[utf8]{inputenc}
\usepackage{amsthm}
\usepackage{slashed}
\usepackage[breakable, theorems, skins]{tcolorbox}
\usepackage[colorlinks = true,
            linkcolor = red,
            urlcolor  = blue,
            citecolor = red,
            anchorcolor = red]{hyperref}
\usepackage{cleveref}

\usepackage{soul}
\usepackage{frcursive}
\usepackage{booktabs}
\usepackage{enumitem}
\usepackage{cancel}

\usepackage{array}
\newcolumntype{C}[1]{>{\centering\arraybackslash}m{#1}}

% -------------------------------------------------------------------
% Theorem Styles
% -------------------------------------------------------------------

\theoremstyle{definition} % Define theorem styles here based on the definition style (used for definitions and examples)
\newtheorem*{definition}{Definition}

\theoremstyle{plain} % Define theorem styles here based on the plain style (used for theorems, lemmas, propositions)
\newtheorem{theorem}{Theorem}[section]
\newtheorem{axiom}{Axiom}
\newtheorem{corollary}[theorem]{Corollary}
\newtheorem{lemma}[theorem]{Lemma}
\newtheorem{proposition}[theorem]{Proposition}

\theoremstyle{remark} % Define theorem styles here based on the remark style (used for remarks and notes)
\newtheorem*{notation}{Notation}
\newtheorem*{solution}{Solution}
\newtheorem*{terminology}{Terminology}

\newtheoremstyle{underline}% name
{}        % Space above, empty = `usual value'
{}              % Space below
{}              % Body font
{}    % Indent amount (empty = no indent, \parindent = para indent)
{}              % Thm head font
{.}             % Punctuation after thm head
{1.5mm}         % Space after thm head: \newline = linebreak
{{\underline{\textit{\thmname{#1}\thmnumber{ #2}}~\thmnote{(#3)}\unskip}}}  % Thm head spec

\theoremstyle{underline}

\newtheorem{remark}[theorem]{Remark}
\newtheorem{example}[theorem]{Example}
\newtheorem{claim}[theorem]{Claim}
% -------------------------------------------------------------------
% Chapter Headings
% -------------------------------------------------------------------

\setcounter{chapter}{-1}

\makeatletter
%\renewcommand{\@chapapp}{Lecture}
\makeatother
\definecolor{lightergray}{rgb}{0.9,0.9,0.9}

\usepackage{titlesec}
\titleformat{\section}{\large\bfseries\raggedright}{}{0em}{\colorsection}[\titlerule]
\titleformat{name=\section,numberless}{\large\scshape\bfseries\raggedright}{}{0em}{\colorsectionnonumber}[\titlerule]

\titleformat{\subsection}{\bfseries\raggedright}{}{0em}{\colorsubsection}
\titleformat{name=\subsection,numberless}{\bfseries\raggedright}{}{0em}{\colorsubsectionnonumber}

\newcommand{\colorsection}[1]{%
    \colorbox{lightergray}{\parbox{\dimexpr\textwidth-2\fboxsep}{\thesection\ \ #1}}}
\newcommand{\colorsectionnonumber}[1]{%
    \colorbox{lightergray}{\parbox{\dimexpr\textwidth-2\fboxsep}{#1}}}
    
\newcommand{\colorsubsection}[1]{%
    \colorbox{lightergray}{\parbox{\dimexpr\textwidth-2\fboxsep}{\thesubsection\ #1}}}
\newcommand{\colorsubsectionnonumber}[1]{%
    \colorbox{lightergray}{\parbox{\dimexpr\textwidth-2\fboxsep}{#1}}}
    
\definecolor{gray75}{gray}{0.75}
\newcommand{\hsp}{\hspace{20pt}}
\titleformat{\chapter}[hang]{\Huge\bfseries}{\thechapter\hsp\textcolor{gray75}{|}\hsp}{0pt}{\Huge\bfseries}

% Begin/end Shortcuts
\def\bc{\begin{corollary}}
\def\ec{\end{corollary}}
\def\bcl{\begin{claim}}
\def\ecl{\end{claim}}
\def\bd{\begin{definition}}
\def\ed{\end{definition}}
\def\ben{\begin{enumerate}}
\def\een{\end{enumerate}}
\def\be{\begin{equation}}
\def\ee{\end{equation}}
\def\bse{\begin{equation*}}
\def\ese{\end{equation*}}
\def\bex{\begin{example}}
\def\eex{\end{example}}
\def\bit{\begin{itemize}}
\def\eit{\end{itemize}}
\def\bl{\begin{lemma}}
\def\el{\end{lemma}}
\def\bnn{\begin{notation}}
\def\enn{\end{notation}}
\def\bn{\begin{note}}
\def\en{\end{note}}
\def\bp{\begin{proposition}}
\def\ep{\end{proposition}}
\def\bq{\begin{proof}}
\def\eq{\end{proof}}
\def\br{\begin{remark}}
\def\er{\end{remark}}
\def\bs{\begin{solution}}
\def\es{\end{solution}}
\def\btab{\begin{table}}
\def\etab{\end{table}}
\def\btb{\begin{tabular}}
\def\etb{\end{tabular}}
\def\bt{\begin{theorem}}
\def\et{\end{theorem}}
\def\btik{\begin{tikzpicture}}
\def\etik{\end{tikzpicture}}
\def\bter{\begin{terminology}}
\def\eter{\end{terminology}}

% Letter Shortcuts
% Greek letters
\def\a{\alpha}
\def\b{\beta}
\def\del{\delta}
\def\g{\gamma}
\def\l{\lambda}
\def\n{\nabla}
\def\sig{\sigma}

% Mathbb
\def\b1{\mathbb{1}} % identity
\def\C{\mathbb{C}} % complex numbers
\def\F{\mathbb{F}}
\def\N{\mathbb{N}} % Natural
\def\Q{\mathbb{Q}} % Rational
\def\R{\mathbb{R}} % Reals
\def\Z{\mathbb{Z}} % Integers

% Mathcal
\def\cA{\mathcal{A}}
\def\cB{\mathcal{B}}
\def\cD{\mathcal{D}}
\def\cF{\mathcal{F}}
\def\cH{\mathcal{H}}
\def\cJ{\mathcal{J}}
\def\cK{\mathcal{K}}
\def\cL{\mathcal{L}}
\def\cM{\mathcal{M}}
\def\cN{\mathcal{N}}
\def\cO{\mathcal{O}}
\def\cP{\mathcal{P}}
\def\cT{\mathcal{T}}
\def\cQ{\mathcal{Q}}

% Mathfrak 
\DeclareMathOperator{\GL}{GL}
\def\gl{\mathfrak{gl}}
\DeclareMathOperator{\Ort}{O}
\def\ort{\mathfrak{o}}
\DeclareMathOperator{\SL}{SL}
\def\sl{\mathfrak{sl}}
\DeclareMathOperator{\SO}{SO}
\def\so{\mathfrak{so}}
\DeclareMathOperator{\SU}{SU}
\def\su{\mathfrak{su}}
\def\fF{\mathfrak{F}}

% Misc. 
\def\diag{\text{diag}}
\def\p{\partial}
\def\ss{\subset}
\def\se{\subseteq}
\def\sm{\setminus}
\def\ve{\varepsilon}
\def\vn{\varnothing}
\def\cl{\colon}
\def\la{\langle}
\def\ra{\rangle}


% Boxes 
\DeclareRobustCommand{\mybox}[2][colframe=blue!10!black,before skip=10pt]{%
\begin{tcolorbox}[
        breakable,
        width=\dimexpr\textwidth\relax, 
        boxsep=5pt,
        arc=5pt,outer arc=5pt,
        ]
        #2
\end{tcolorbox}
}
\def\bbox{\begin{tcolorbox}[colback=white!5,colframe=blue!75!black,title=Exercise]}
\def\ebox{\end{tcolorbox}}

% Feynman Diagram Stuff 
\def\midarrow{\draw[thick, decoration={markings, mark=at position 0.5 with {\arrow{>}}}, postaction={decorate}]}
\def\midarrowblue{\draw[thick, blue, decoration={markings, mark=at position 0.5 with {\arrow{>}}}, postaction={decorate}]}
\def\midarrowred{\draw[thick, red, decoration={markings, mark=at position 0.5 with {\arrow{>}}}, postaction={decorate}]}
\def\aftermidarrow{\draw[thick, decoration={markings, mark=at position 0.75 with {\arrow{>}}}, postaction={decorate}]}
\def\aftermidarrowblue{\draw[thick, blue, decoration={markings, mark=at position 0.7 with {\arrow{>}}}, postaction={decorate}]}
\def\aftermidarrowred{\draw[thick, red, decoration={markings, mark=at position 0.7 with {\arrow{>}}}, postaction={decorate}]}
\def\beforemidarrow{\draw[thick, decoration={markings, mark=at position 0.25 with {\arrow{>}}}, postaction={decorate}]}
\def\wavey{\draw[thick, decorate, decoration={snake, segment length=2.5mm, amplitude=0.75mm}]}
\newcommand{\AxisRotatorL}[1][rotate=0]{%
    \tikz [x=0.25cm,y=0.60cm,line width=.2ex,-stealth,#1] \draw (0,0) arc (-150:150:1 and 1);%
}
\newcommand{\AxisRotatorR}[1][rotate=0]{%
    \tikz [x=0.25cm,y=0.60cm,line width=.2ex,-stealth,#1] \draw (0,0) arc (150:-150:1 and 1);%
}

\begin{document}

\captionsetup[figure]{margin=1.5cm,font=small,labelfont={bf},name={Figure},labelsep=colon,textfont={it}}
\captionsetup[table]{margin=1.5cm,font=small,labelfont={bf},name={Table},labelsep=colon,textfont={it}}
\setlipsumdefault{1}

\frontmatter

\begin{titlepage}
	\centering
	    \scshape % Use small caps for all text on the title page
        \vspace*{\baselineskip} % White space at the top of the page
        
	    \rule{\textwidth}{1.6pt}\vspace*{-\baselineskip}\vspace*{2pt} % Thick horizontal rule
	    \rule{\textwidth}{0.4pt} % Thin horizontal rule
	    
	    \vspace{0.75\baselineskip} % Whitespace above the title
	    
	    {\LARGE Supersymmetry} % Title
	    
	    \vspace{0.75\baselineskip} % Whitespace below the title
	    
	    \rule{\textwidth}{0.4pt}\vspace*{-\baselineskip}\vspace{3.2pt} % Thin horizontal rule
	    \rule{\textwidth}{1.6pt} % Thick horizontal rule
	    
        \vspace{5\baselineskip} % Whitespace after the title block

	    Course delivered in 2020 by 
	
	    \vspace{0.5\baselineskip} % Whitespace before the editors
	
	    {\scshape\Large Dr. Stefano Cremonesi} % Lecturer Name
	
	    \vspace{0.5\baselineskip} % Whitespace below the editor list
	
	    \textit{Durham University} % Lecturer Institution
	    
	    \vspace{2\baselineskip} % Whitespace after the title block
	    
	    \includegraphics[width=8cm]{images/DurhamLogo.png}\\[1cm] % Logo 

	    Notes taken by 
	
	    \vspace{0.5\baselineskip} % Whitespace before my name
	
	    {\scshape\Large Richie Dadhley} % Lecturer Name
	   
	    \vspace{0.5\baselineskip} % Whitespace below my name
	    \textit{richie@dadhley.com} % Email
	
	    \vfill % Whitespace between editor names and publisher logo
\end{titlepage}



% -------------------------------------------------------------------
% Acknowledgements
% -------------------------------------------------------------------

\newpage
\section*{Acknowledgements}

These are my notes on the 2020 lecture course "Supersymmetry" taught by Dr. Stefano Cremonesi at Durham University as part of the Particles, Strings and Cosmology Msc. For reference, the course lasted 16 hours and was taught over 4 weeks. \\

I have tried to correct any typos and/or mistakes I think I have noticed over the course. I have also tried to include additional information that I think supports the taught material well, which sometimes has resulted in modifying the order the material was taught. Obviously, any mistakes made because of either of these points are entirely mine and should not reflect on the taught material in any way. \\

% For any other notable sources used:
% One of the main sources I have used to obtain the additional information is Dr. David Tong’s String Theory notes, available via the following link  

%\begin{center}
%    \textcolor{red}{Link}
%\end{center}

I would like to extend a message of thanks to Dr. Cremonesi for teaching this course. I would also like to thank Thimo Preis for helping getting these notes started.\\

If you have any comments and/or questions please feel free to contact me via the email provided on the title page. \\

These notes are currently a work in progress, so for updated notes, as well as a list of other notes/works I have available, visit my blog site

\begin{center}
    \href{https://richie291.wixsite.com/theoreticalphysics}{https://richie291.wixsite.com/theoreticalphysics}
\end{center}

These notes are not endorsed by Dr. Cremonesi or Durham University.

\vspace{1cm}

\begin{flushright}
    \Huge{{\cursive\setul{0.1ex}{}\ul{~~Richie Dadhley~~}}}
\end{flushright}

% -------------------------------------------------------------------
% Contents
% -------------------------------------------------------------------

\tableofcontents

% -------------------------------------------------------------------
% Main sections (as required)
% -------------------------------------------------------------------

\mainmatter

\chapter{Introduction}

It is always useful to have an introduction and motivation for what it to come. This is especially true when the material can become rather abstract, as otherwise it's very easy to get lost in the world of equations and algebra. We therefore start with just this. As with most introductions\footnote{At least in my experience.} it is likely that some of the stuff written here will mean nothing to the reader, this should not deter from reading on. Instead the introduction is meant to introduce us to what we're going to study and, perhaps more importantly, why we care. This chapter, then, can be viewed more as a `grounding' point to revisit when questions about what on Earth we're doing arise. 

So without further ado, let's go. 

\section{SUSY: What is it?}
\label{sec:WhatIsSUSY}

The first question we should ask it "What \textit{is} SUSY?" Well "SUSY" itself stands for supersymmetry, but then we just ask "what is supersymmetry?" Well as the name suggests, it is a \textit{spacetime symmetry}, much like the Poincar\'{e} symmetry group, but what is "super" about it? Well, as we will see later, it turns out to map Bosons to Fermions and vice versa
\bse
    \text{Boson, integer spin } \ket{B} \qquad \Longleftrightarrow \qquad \ket{F} \text{ Fermions, half-odd spin}.
\ese

As with the Poincar\'{e} symmetry of `normal' QFT, there will be conserved charges associated to our SUSY, which we creatively call \textit{supercharges}.\footnote{Get ready to start sticking the word "super" in front of everything...} There will be two such charges, and it is standard to denote them by $Q$ and $\bar{Q}$. Our symmetry map above can then be written as 
\bse 
    Q\ket{B} = \ket{F} \qquad Q \ket{F} = \ket{B} \qquad \bar{Q}\ket{B} = \ket{F} \qand \bar{Q}\ket{F} = \ket{B},
\ese 
where each $\ket{F}/\ket{B}$ are meant to just mean \textit{some} Fermion/Boson state, i.e. they're not all the same states. 

\bnn 
    As with the above, we will often write formulas that hold for both the barred expressions and unbarred expressions. In order to save essentially writing everything twice, we will adopt the notation of putting a tilde\footnote{It might seem more reasonable to use $\overset{(-)}{Q}$ etc, but personally I don't think this typesetting looks nice, so tilde it is.} to cover both cases. That we we can write the above simply as
    \be 
    \label{eqn:QOnBAndF}
        \widetilde{Q}\ket{B} = \ket{F} \qand \widetilde{Q}\ket{F} = \ket{B}. 
    \ee 
\enn 

Let's make some comments/introduce some terminology. 
\ben[label=(\roman*)]
    \item We denote the index structure on our supercharges using the early part of the Greek alphabet, e.g. $\a,\beta$ etc. We also adopt the standard notation that barred things have a dotted index. In other words we have $Q_{\a}$ and $\bar{Q}_{\dot{\a}}$. 
    \item It follows from the fact that Bosons carry integer spin and Fermions half-odd spin and \Cref{eqn:QOnBAndF} that we require $Q_\alpha$, $\bar{Q}_{\dot{\a}}$ to carry spin $1/2$.
    \item These charges will form an algebra, just like the generators of the Poincar\'{e} group. This algebra will, of course, have representations, and we call the irreducible representations (henceforth just \textit{irreps}) \textit{supermultiplets}. Particles/fields in same supermultiplet are called \emph{superpartners}.\footnote{I told you, get ready to stick "super" in front of everything...}
    \item Our supercharges are not completely blind to other symmetries of the QFT, and indeed relations arise. An obvious example is the Poincar\'{e} transformations $\{P, M\}$, where $P$ are the spacetime translations and $M$ our Lorentz transformations. Perhaps less obvious examples are so-called internal symmetries,\footnote{Sometimes we also call these flavour symmetries, to distinguish from $R$-symmetries} which we denote by $B$, and so-called $R$-symmetries which transform the different supercharges into each other.\footnote{A bit more technically, an $R$-symmetry is the largest subgroup of the automorphism group of the SUSY algebra that commutes with the Lorentz group. That is, it is the largest group that commutes with the Lorentz group, rotates the supercharges between each other \textit{and} leaves the anticommutator \Cref{eqn:QQbarAnticommutator} invariant.} The commutation relations turn out to schematically be 
    \bse 
    [P,\widetilde{Q}] = 0, \qquad [M, \widetilde{Q}] \propto \widetilde{Q}, \qquad [B,\tilde{Q}]=0, \qand [R,\widetilde{Q}]\propto \widetilde{Q}.
    \ese 
    It follows from these relations that superpartners have
    \ben
         \item same mass (if SUSY is preserved by vacuum),
         \item different spin (raise or lower spin by one-half by applying SUSY),
         \item same quantum numbers under internal global symmetries,
         \item different $R$-charge.
    \een
    \item As the supercharges map Fermions into Bosons, they have to be anti-commuting. This is just because Bosons commute but Fermions anticommute, so we need to account for this. Thus we consider
    \be 
    \label{eqn:QQbarAnticommutator}
        \{Q,\bar{Q}\} \propto P 
    \ee 
    thus applying an infinitesimal SUSY transformation squared gives you a translation. Recall that in GR we get diffeomorphism invariance by gauging out by translations. The above tells us that if gauge by SUSY we also gauge by translations. This gives us a theory which is supersymmetric and diffeomorphism-invariant, which is \emph{supergravity}.
    \item As we have written it so far, we only have one copy of SUSY. We can  extend SUSY by having several copies of $\{Q_{\a}, \bar{Q}_{\dot{\a}}\}$. We denote each copy with a capital Latin index, e.g. $Q^I_{\a},\bar{Q}^J_{\dot{\a}}$ with $I,J=1,\dots,\cN$. We generate a supermultiplet by starting from lowest-weight state or highest-weight state and then by acting with $Q,\bar{Q}$. The number of states will increase exponentially if you increase amount of SUSY, i.e. increase $\cN$. The maximum spin in a supermultiplet will grow linearly with $\cN$, however we require that we have \textit{no} interacting degrees of freedom with spin 
    \ben 
        \item $>1$ when gravity is absent. In this case we have $\cN \leq 4$, which is seen by the fact that $4$ supercharges allows us to go from spin $-1$ to spin $+1$. This is super-quantum field theory (SQFT).
        \item $>2$ when we have gravity. In this case we have $\cN \leq 8$. This corresponds to (4D) supergravity. 
    \een 
\een

As we just mentioned, SQFT is a thing, the question is "what is it?" Well, as it's a QFT so we will have 
\ben[label=(\roman*)]
    \item Some field content,
    \item Some Lagrangian/Action
\een
both of which are constrained by SUSY. The task of this course is to answer the question of "what are these constraints?" 

\br 
    Nowadays you look at theories which do not have a Lagrangian formulation (like for CFT), but you will still have constraints by SUSY.
\er 

\section{Why SUSY?}

Ok now that we have an introductory knowledge of \textit{what} SUSY is, we now want to answer the question of "\textit{why} do we care about SUSY?"

\subsection{Theoretical reasons}

\subsubsection{Most General Symmetry Of Interacting QFTs}

SUSY is the most general symmetry for interacting theories. This is essentially given by the following theorems.\footnote{This is one case where you shouldn't panic if this part of the introduction seems scary. This will become more clear as we go on.}

\bt[Coleman-Mandula theorem (1967)]
    Consider a unitary, relativistic QFT with finitely many d.o.f. below a mass scale, i.e. with mass $<M$ (for any $M$) and assume there to be an analytic, non-trivial S-matrix. Then the Lie group of symmetries of such a theory (or of the S-matrix) is 
    \begin{center}
        (Poincar\'{e}) $\times$ (compact internal symmetry).
    \end{center} 
\et

\br 
    If you relax particle finiteness, then there is a similar statement where the Poincare group is replaced by the conformal group
    \begin{center} 
    (Conformal) $\times$ (Compact Internal Symmetry).
    \end{center}
\er 

\bt[Haag-Lopuszanski-Sohnius theorem (1975)]
\label{thrm:HLS}
    Extend the symmetry algebra to a \textbf{superalgebra}, or \textbf{graded Lie algebra}, that includes anticommutators. Then\footnote{The symbol $\ltimes$ is a \textit{semidirect product}, with the SuperPoincar\'{e} group being the so-called \textit{normal subgroup}. It doens't matter too much what this means apart from that it implies that the commutator of an element of the R symmetry group and the SuperPoincar\'{e} group is an element of the SuperPoincar\'{e} group. In other words $[R,\widetilde{Q}] \propto \widetilde{Q}$, which we have already seen above.}
    \begin{center}
        (R-symmetry) $\ltimes$ (SuperPoincar\'{e}) $\times$ (Compact Internal Symmetry)
    \end{center}
\et 

\br 
     Note that Superpoincar\'{e} group is a subset of Superconformal group, and if you relax particles finiteness you get
    \begin{center}
        (R-symmetry) $\ltimes$ (SuperConformal) $\times$ (Compact Internal Symmetry)
    \end{center}
\er 

\subsubsection{Generic Prediction Of String Theory}

It turns out that SUSY is crucial for string theory, in particular it is crucial for the stability of the vacuum. In string theory it is needed in order to solve the problems of so-called \textit{tachyons}, which are particles that have negative mass-squared. On top of this, the non-SUSY string theory (sometimes known as Bosonic string theory) has no (clear) way to introduce Fermions into the theory. However all the symmetry of the system are used to constrain the Bosonic theory, and so if we want to introduce Fermions, we need some more symmetry, i.e. SUSY, so that we can constrain our Fermion fields correctly. 

\subsubsection{SQFTs As Theoretical Laboratories}

You can use SUSY QFT as theoretical laboratory, and it gives you improved quantum behaviour, which in turn allows us to control theory better (the more SUSY introduced the more control). This allows us to obtain exact results (often), at least for subset of SUSY theories, even at strong coupling.

\subsection{Phenomenological reasons for SUSY}

Besides the theoretical reasons given above, there are also several phenomenological reasons for studying SUSY. Let's now give a few. 

\subsubsection{Naturalness vs. Fine Tuning: The Hierarchy Problem}

It's an experimental fact that electroweak symmetry breaking occurs at 
\bse 
    m_{\cancel{EW}} \sim 250GeV << m_{\text{Planck}} \sim 10^{19}GeV.
\ese
Why is this a problem? Well consider the Higgs $2$-point function 
\begin{center}
    \btik 
        \draw[thick, dashed] (-2,0) -- (-0.75,0);
        \node at (-2.2,0) {$H$};
        \draw[thick, dashed] (0.75,0) -- (2,0);
        \node at (2.2,0) {$H$};
        \beforemidarrow (0,0) circle [radius=0.75cm];
        \aftermidarrow (0,0) circle [radius=0.75cm];
        \node at (0,1) {$\psi$};
        \node at (0,-1) {$\psi$};
        \draw[fill=black] (-0.75,0) circle [radius=0.05cm] node [right] {$\l_F$};
        \draw[fill=black] (0.75,0) circle [radius=0.05cm] node [left] {$\l_F$};
        \node[right] at (2.5,0) {$\sim \, -|\l_f|^2\Lambda_{UV}^2$};
    \etik 
\end{center}
This is quadratically divergent\footnote{Perhaps more technically,  quadratically sensitive to new UV physics.} this leads to a renormalisation of $m_H^2$ and we would expect $m_H \sim \Lambda_{UV}$, which tells us that \textit{theoretically} we expect
\bse 
    m_{\cancel{EW}} \sim \Lambda_{UV}.
\ese 
However this doesn't agree with our experimental result above. This constitutes what is known as a \textit{hierarchy problem}, i.e. you would have to fine tune parameters by many orders of magnitude to get this. This would be unnatural and actually quantum corrections would spoil the fine tuning.

So how do we fix this? Well suppose there existed a scalar $S$ with $\lambda_S H^2 |S|^2$ where the two-point function has one self-interaction contribution
\begin{center}
    \btik 
        \draw[thick, dashed] (-2,0) -- (2,0);
        \node at (-2.2,0) {$H$};
        \node at (-2.2,0) {$H$};
        \draw[thick, dashed] (0,0.75) circle [radius=0.75cm];
        \node at (0,1.7) {$S$};
        \draw[fill=black] (0,0) circle [radius=0.05cm] node [below] {$\l_S$};
        \node[right] at (2.5,0) {$\sim \, +\l_S\Lambda_{UV}^2$};
    \etik 
\end{center}
where we notice the difference in sign compared to the $\l_F$ case above. Therefore \textit{if} $\lambda_S=|\lambda_f|^2$, then the quadratic divergences cancel and hierarchy problem is solved.\footnote{Also turns out that the log divergence cancels too.} In SUSY, the scalar $S$ would would be the superpartner to the above Fermion. This allows them to cancel as they sit in the same supermultiplet. This argument works perturbatively to all orders and also non-perturbatively.

This all seems great, but as of yet, we have not observed any superpartners in our collider experiments. This tells us that if SUSY is relevant at all in describing nature, it \textit{must} be broken in nature. The scale of the SUSY breaking would occur in the range
\bse 
    10^3 GeV < m_{\cancel{SUSY}} \leq m_{\text{Planck}}
\ese 
where the latter inequality stems from the fact that you want SUSY for a quantum gravity theory, and such it should be unbroken at Planck scale.

It turns out that if SUSY is broken spontaneously, the quadratic divergences still cancel and so the hierarchy problem is still fixed. However it turns out that the log divergences comeback. This seems bad, and we need to introduce new corrections to account for this. However this still reduces the amount of fine tuning needed to a reasonable level. In a minimal SUSY extension of the standard model (MSSM) people argued that a reasonable fine tuning estimate 
\bse 
    m_{\cancel{SUSY}} \sim 1 TeV,
\ese 
which is known as \textit{low energy SUSY}. This scale is already in a struggle with the LHC, but the tension can be reduced by increasing the fine tuning slightly, or by making modifications to the MSSM. 

\subsubsection{Gague Coupling Unification}

One of the major goals of high energy particle physics is grand unification theories (GUTs). These stem from the fact that in the SM the gauge couplings seem to tend towards a common point, but are off ever so slightly. 
\begin{center}
    \btik 
        \draw[->] (-0.5,0) -- (5.5,0);
        \draw[->] (0,-0.5) -- (0,3);
        \node[right] at (5.2,-0.3) {Energy};
        \node at (-0.6,2.8) {$\frac{1}{g_i^2}$};
        \draw[thick, blue] (0.5,0.5) -- (4.5,1.5);
        \draw[thick, red] (0.5,1) -- (4.5,1.4);
        \draw[thick] (0.5,2.5) -- (4.5,0.5);
        \draw (3.7,0.2) --(3.7,-0.2) node [below] {$10^{15}$GeV};
        \node at (6.5,2.5) {U(1)};
        \node at (6.5,1.5) {\textcolor{red}{SU(2)}};
        \node at (6.5,0.5) {\textcolor{blue}{SU(2)}};
    \etik 
\end{center}

It would obviously be much nicer of nature if they did indeed meet perfectly and combine into one mother-of-all couplings. This would occur if we actually had a bigger Lie group at higher energy scales, which was broken by the vacuum expectation value (vev) of some field to the $U(1)\times SU(2)\times SU(3)$ of the SM. 
\bse 
    G_{GUT} \stackrel{\expval{\phi}\neq 0}{\longrightarrow} G_{SM} \stackrel{\expval{H}}{\longrightarrow } SU(3) \times U(1)_{em}
\ese 
with $m^2_{GUT} \sim 10^{15,16} GeV$ and $m_{\cancel{EM}} \sim 10^2$.

\begin{center}
    \btik 
        \draw[->] (-0.5,0) -- (5.5,0);
        \draw[->] (0,-0.5) -- (0,3);
        \node[right] at (5.2,-0.3) {Energy};
        \node at (-0.6,2.8) {$\frac{1}{g_i^2}$};
        \draw[thick] (0.5,0.5) -- (3.5,1.2);
        \draw[thick] (0.5,1) -- (3.5,1.2);
        \draw[thick] (0.5,2.5) -- (3.5,1.2);
        \draw[thick] (3.5,1.2) -- (5,1.3);
        \draw[fill=black] (3.5,1.2) circle [radius=0.07cm];
        \draw (3.5,0.2) -- (3.5,-0.2) node [below] {$m_{GUT}$};
    \etik 
\end{center}

The problems with non-SUSY GUT theories are
\ben[label=(\roman*)]
    \item The couplings don't quite meet up.
    \item GUT Yakawa couplings would induce a 
    \bse 
        p^+ \to e^+ + \pi^0
    \ese 
    decay channel. This is ruled out by the proton lifetime. 
    \item We get a new hierarchy problem for $m_{\cancel{EM}} << m_{GUT}$.
\een 

The claim is that low energy SUSY gives us that 
\ben[label=(\roman*)]
    \item The couplings meet (within error bars) at $m_{GUT} \sim 10^{16}GeV$.
    \item Proton decay not a problem. 
    \item Hierarchy is maintained. 
\een 

\subsubsection{Dark Matter}

Another big current area for phenomenological physics is dark matter. 

\textcolor{red}{Fill in later.}
\chapter{Poincar\'{e}, Lorentz \& Spinors}

If we are going to construct our superPoincar\'{e} group and its irreps, of course it's important that we understand the `regular' Poincar\'{e} group properly first. We therefore start the course with such a discussion. 

\section{Poincar\'{e} group}

The Poincar\'{e} group is both the Lorentz group, $SO(1,3)$, and the set of spacetime translations. We denote the Poincar\'{e} group by $ISO(1,3)$ and it acts as follows
\be 
    ISO(1,3) : \; x^{\mu} \mapsto {\Lambda^{\mu}}_{\nu} x^{\nu} + a^{\mu} \qquad s.t. \qquad  \Lambda^T \eta \Lambda = \eta 
\ee 
where ${\Lambda^{\mu}}_{\nu}$ are the generates of the Lorentz group and $a^{\mu}$ is a spacetime translation. As we know the Lorentz group has $6$ independent components (the boosts and spatial rotations) and we collect these into an object denoted $M_{\mu\nu} = - M_{\nu\mu}$. Similarly we have $4$ spacetime translations, which we collect into $P_{\mu}$. These are the generators of the group and so live in the Lie algebra. They then obey a set of Lie bracket (which here is just the commutator) relations. These turn out to be 
\mybox{
    \be 
    \label{eqn:PoincareCommutator}
        \begin{split}
            [P_{\mu}, P_{\nu} ] & = 0 \\
            [M_{\mu \nu}, P_{\nu}] & = - i \eta_{\mu \rho} P_{\nu} + i\eta_{\nu \rho} P_{\mu} \\
            [M_{\mu \nu}, M_{\rho \sigma}] & = -i \eta_{\mu \rho} M_{\nu \sigma} + i \eta_{\nu\rho}M_{\mu\sig} - i \eta_{\nu\sig} M_{\mu\rho} + i \eta_{\mu\sig} M_{\nu\rho},
        \end{split}
    \ee 
}
\noindent where we use the "mostly minus" convention $\eta_{\mu\nu} = \text{diag}(+,-,-,-)$. 

In order to label our representations, we want to find the Casimir operators of this algebra. That is we find matrices which commute with every element of the algebra, and then Schur's Lemma tells us that we can label our irreps via these. 

\bcl 
    We have two Casimir invariants for our Poincare algebra given by 
    \be
    \label{eqn:PoincareCasimirs}
        P^2 := P^{\mu}P_{\mu} \qand W^2 := W^{\mu}W_{\mu},
    \ee 
    where 
    \bse 
        W^\mu = \frac{1}{2} \epsilon^{\mu \nu \rho \sigma} P_\nu M_{\rho \sigma}
    \ese
    is the so-called \textit{Pauli-Lebanski} vector.
\ecl 

\bbox 
    Prove the above claim. That is show that \Cref{eqn:PoincareCasimirs} commute with our generators $P_{\mu}$ and $M_{\mu\nu}$. 
\ebox 

The irreps of the Poincar\'{e} group are then our particles. These split into two general cases
\begin{enumerate}
    \item \textbf{Massive particles}: Go to rest frame $P_\mu = (m,\Vec{0})$ such that
    \bse 
    P^2=m^2, \qand W^2 = -m^2 s(s+1),
    \ese 
    with $s$ being the spin. Thus, by Schur's Lemma, the irreps can be labelled by mass and spin. Given a certain rep, there will be certain weight state labelled by its eigenvalue $m_s$ of the $z$-direction spin operator, $S_z$.
    \item \textbf{Massless particles}: Here we can't go to rest frame, but we can go to light cone frame $P_\mu = (E,0,0,E)$. Then we get
    \bse 
    P^2=W^2=0, \qand W_{\mu} = M_{12} P_{\mu},
    \ese 
    such that we can not use mass or spin, but we can use angular momentum in the plane orthogonal to direction of motion, $M_{12}P_{\mu}$. The necessary eigenvalues are the helicity $\pm s$. Irreps are then labelled by the absolute value of their helicity, $s$. Thus it is the same for all states in a multiplet. The weight states will be distinguished by the sign of the helicity.
\end{enumerate}

\section{Lorentz group \& $SL(2,\C)$}

We now want to look more closely at the Lorentz subgroup, $SO(1,3)$. As we said above, this acts as
\be 
SO(1,3) :  x^{\mu} \mapsto {\Lambda^{\mu}}_{\nu} x^{\nu},
\ee 
with generators $M_{\mu\nu} = -M_{\nu\mu}$. These generators can be split into rotations, $M_{ij}=\epsilon_{ijk} J_k$, and boosts, $M_{0i}=K_i$. The group of rotations is compact, which tells us that the the generators are Hermitian $J=J^{\dagger}$, whilst the group of boosts is non-compact, which tells us that the generators are \textit{anti}-hermitian $K^\dagger=-K$.\footnote{The non-compactness of the boosts is easily understood by the fact that there is a limit to how much we can boost something. Contrasting that to the fact that we can rotate by whatever angle we like.}

These generators satisfy the commutation relations
\bse 
    [J_i,J_k] = i \epsilon_{ijk} J_k, \qquad [J_i,K_k] = i \epsilon_{ijk} K_k, \qand [K_i,K_j] = - i \epsilon_{ijk} J_k,
\ese 
which can easily be checked. Now, the rotations look nice because they close under the Lie bracket, and so are isomorphic to $\mathfrak{su}(2)$. However the boosts are not playing so nicely and so we really want to do something to fix this. This is a standard problem and the answer is to complexify the Lie algebra by defining
\bse 
    \Vec{J}^{\pm} = \frac{1}{2}(\Vec{J}\pm i \Vec{K}),
\ese 
which are Hermitian. If we then compute the commutators of our $J^{\pm}$s, we see that 
\bse 
    [J^+_i, J^+_j] = i \epsilon_{ijk} J^+_k, \qquad [J^-_i, J^-_j] = i \epsilon_{ijk} J^-_k, \qand [J^+_i,J^-_j]=0.
\ese 
This tells us we have two independent copies of $\mathfrak{su}(2)$. In other words, we can trade the Lorentz algebra, $\mathfrak{so}(1,3)$, for two copies of the $\mathfrak{su}(2)$ algebra at the price of complixification. We can now label the irreps of the Lorentz group via this decomposition by two half-integers $(s_+,s_-)$, where
\bse 
    \big( \Vec{J}^{\pm} \big)^2 = s_{\pm} (s_{\pm} + 1), \qquad s_{\pm} \in \{0, 1/2, 1, ... \}.
\ese 

Note that by complex conjugaton we swap the copies
\bse 
    (s_+,s_-) \Longleftrightarrow (s_-,s_+),
\ese 
which is often indicated by writing
\bse 
    SO(1,3) = \frac{SU(2)\times SU(2)^*}{\Z_2}.
\ese 
The $\Z_2$ quotient imposes that $SO(1,3)$ irreps have spin $s_+ + s_- \in \Z$. 

The object in the `numerator' (i.e. the group which is quotiened), is known as the \textit{spin group}
\bse 
    \text{Spin}(1,3) \equiv SL(2,\C) = SU(2) \times SU(2)^*,
\ese 
and it is the \textit{double cover} of the Lorentz group, as seen by the $\Z_2$ quotient. We label the irrps by
\bse 
(s_+,s_-) \quad s_+ + s_-\in \frac{\Z}{2},
\ese 

\subsection{$SL(2,\C)$ vs. $SO(1,3)$}

Let's flush out the double cover business mention above a bit more. First let's give a definition of the $SL(2,\C)$ group in terms of matrices.
\be 
    SL(2,\C) = \left\{ M = \begin{pmatrix}
        a&b \\
        c&d\\
    \end{pmatrix} \Big| a,b,c,d \in \C, \, \det M = ad-bc=1 \right\}.
\ee 

The double cover idea above is equivalent to the following claim. 

\bcl 
    There is a $2$ to $1$ homomorphism 
    \bse 
        \begin{split}
            \Lambda : SL(2,\C) & \to SO(1,3) \\
            M & \mapsto \Lambda(M).
        \end{split}
    \ese 
    such that
    \be 
    \label{eqn:DoubleCoverConditions}
        \Lambda(M_1M_2) = \Lambda(M_1) \Lambda(M_2) , \qand \Lambda(M)= \Lambda(-M).
    \ee 
\ecl 

\bq 
    First let's look at the conditions \Cref{eqn:DoubleCoverConditions}. The first one just tell us that $\Lambda$ is a group homomorphism,\footnote{A homomorphism is a `structure preserving map', and the structure here is the group multiplication, so what we require is $\Lambda( M_1 \bullet M_2)= \Lambda(M_1)\circ \Lambda(M_2)$, where $\bullet$ is the group multiplication in $SL(2,\C)$ and $\circ$ the one in $S)(1,3)$. Of course we have just suppressed the notation above.} and the second one is what gives us the "2 to 1" bit as both $M$ and $-M$ are mapped to the same element in $SO(1,3)$.
    
    Ok so now we want to try and construct such a $\Lambda$. First we introduce the definition\footnote{Note the difference in index placement. However the two are \textit{not} simply related by raising an index as then we would have $\bar{\sig}^{\mu} = (\b1, - \sig_i)$. We will see shortly how the two are related.} 
    \be
    \label{eqn:SigmaMu}
        \sig_{\mu} := (\b1 ,\sig_i) \qand \bar{\sig}^{\mu} = (\b1,\sigma_i),
    \ee 
    where 
    \bse 
        \sig_1 = \begin{pmatrix}
            0&1 \\
            1 & 0 \\
        \end{pmatrix}, \qquad \sig_2=
        \begin{pmatrix}
            0 & -i \\
            i & 0 \\ 
        \end{pmatrix}, \qquad  \sig_3 = 
        \begin{pmatrix}
            1 & 0 \\
            0& -1 \\
        \end{pmatrix},
    \ese
    are the Pauli matrices. As we know the Pauli matrices with $\b1_{2\times2}$ form a basis for the space of $2\times 2$ complex matrices. We can show that these obey 
    \be 
    \label{eqn:TraceSigBarSig}
        \tr[\sig_{\mu}\bar{\sig}^{\nu}] = 2\del^{\nu}_{\mu}.
    \ee 
    Next, given a $x^{\mu}\in \R^{1,3}$, we define 
    \bse 
        X:= x^\mu \sigma_\mu = 
        \begin{pmatrix}
            x^0 +x^3 & x^1-ix^2 \\
            x^1+i x^2 & x^0-x^3 \\ 
        \end{pmatrix},
    \ese 
    which we can easily check obeys 
    \bse 
        X^\dagger = X, \qand \det X = x^\mu x_\mu. 
    \ese 
    We can invert this using \Cref{eqn:TraceSigBarSig} to get 
    \bse 
        x^{\mu} = \frac{1}{2}\tr[X \bar{\sig}^{\mu}].
    \ese 
    Finally consider taking a homogeneous, linear transformation on $X$ as
    \bse 
        X \to X^{\prime} = M X M^{\dagger}, \qquad \text{with} \qquad M \in SL(2,\C).
    \ese 
    It follows from the definition of $X$ that this transformation also holds on $x^{\mu}$ as 
    \be
    \label{eqn:xPrime}
        x^{\prime\nu}\sig_{\nu} = M \sig_{\nu} M^{\dagger} x^{\nu},
    \ee 
    where we have used the fact that $x^{\nu}\in\R$ for a given $\nu$ so can just move it about freely. Now it's easy to check that the above transformation preserves Hermiticity of $X$ (which is the same as preserving reality of $x^{\nu}$) and it also preserves the determinant condition. In particular we have $x^{\prime\mu}x_{\mu}^{\prime} = x^{\mu}x_{\mu}$, which is a Lorentz transformations! So we can extract the transformation 
    \bse 
        x^{\prime \mu} = {\Lambda^{\mu}}_{\nu}(M) x^{\nu} 
    \ese 
    and simply read off ${\Lambda^{\mu}}_{\nu}$: simply multiply both sides of \Cref{eqn:xPrime} by $\frac{1}{2}\bar{\sig}^{\mu}$ and take a trace to get
    \bse 
        x^{\prime\mu} = \frac{1}{2}\tr[M\sig_{\nu}M^{\dagger}\bar{\sig}^{\mu}] x^{\nu},
    \ese 
    which let's us read off 
    \be 
    \label{eqn:Lambda(M)}
        {\Lambda^{\mu}}_{\nu}(M) = \frac{1}{2}\tr[M\sig_{\nu}M^{\dagger}\bar{\sig}^{\mu}].
    \ee 
    It is trivial to see that this obeys our $\Lambda(M) = \Lambda(-M)$ condition, and so we have our 2 to 1 homomorphism. 
\eq 

\section{Spinors}

Now that we know what the spin group is, we want to introduce two $\C$ component spinors. We will focus mostly on Chiral/Weyl spinors, but will also briefly mention Dirac and Majoranna spinors.

\subsection{Weyl Spinors}

Let's focus on Weyl\footnote{We shall flip flop between saying "Weyl" and "Chiral" in these notes.} spinors. These transform in the two fundamental irreps of $SL(2,\C)$, which we denote by 
\bse
    \begin{split}
        \text{Left-Handed} \qquad & (1/2,0) \\
        \text{Right-Handed} \qquad & (0,1/2).
    \end{split}
\ese 

\bnn 
    We use a notation where lower indices indicate the \textit{row} and upper indices the \textit{column}. For example 
    \bse 
        ({M_{\a}}^{\beta}) = \begin{pmatrix}
            {M_1}^1 & {M_1}^2 & ... & {M_1}^N \\
            {M_2}^1 & {M_2}^2 & ... & {M_2}^N \\
            \vdots & \vdots & \ddots & \vdots \\
            {M_N}^1 & {M_N}^2 & ... & {M_N}^N
        \end{pmatrix}.
    \ese 
\enn 

Left-handed Weyl spinors transform in the fundamental representation, i.e. 
\bse 
    \psi \mapsto M \psi \qquad M \in SL(2,\C),
\ese
which we can write in components as
\bse 
    \psi_{\a} \mapsto \psi'_{\a} = {M_{\a}}^{\beta} \psi_{\beta}.
\ese 
Similarly for right-handed Weyl spinors transform in the antifundamental:
\bse 
    \bar{\psi}_{\dot{\a}} \mapsto  {(M^*)_{\dot{\a}}}^{\dot{\beta}} \bar{\psi}_{\dot{\beta}} = \bar{\psi}_{\dot{\beta}} (M^{\dagger}{)^{\dot{\beta}}}_{\dot{\a}},
\ese 
where we have used the standard notation of putting dots on antifundamenal indices. This will be used throughout the course and can simply be remembered via "barred objects come with dotted indices." This is equivalent to saying
\bse 
    (1/2,0)^* = (0,1/2),
\ese 
and so we can identify $\bar{\psi}_{\dot{\a}} := (\psi_{\a})^*$ or $\bar{\psi}_{\dot{\a}} := (\psi_{\a})^{\dagger}$, depending on whether we are looking at a number or an operator, respectively. 

Now our experience with GR tells us that its often very useful to raise and lower indices in order to make contractions etc. The question if "how do we do this here?" It is tempting to say "just use the metric $\eta_{\mu\nu}$", but then we realise we can't do this. Why? Well the easiest way to see this is because it has the wrong index structure, i.e. $\mu,\nu$ are spacetime indices but our $\a,\beta,\dot{\a},\dot{\beta}$ are $SU(2)$ indices. What do we do then? Well we recall that a Lie group is, in particular, a manifold and so we can define the following 2-forms on $SU(2)$ and $SU(2)^*$. These are given in matrix form as\footnote{Really we shouldn't put an equal sign between the dotted and undotted $\epsilon$s, as they live in different spaces. Basically all we're saying is that their matrix forms are the same.}
\bse 
    (\epsilon^{\a\beta}) = (\epsilon^{\dot{\a}\dot{\beta}}) = \begin{pmatrix}
        0 & 1 \\
        -1 & 0
    \end{pmatrix} \qand (\epsilon_{\a\beta}) = (\epsilon_{\dot{\a}\dot{\beta}}) = \begin{pmatrix}
        0 & -1 \\
        1 & 0
    \end{pmatrix}
\ese 
Where the two matrices are inverses of each other, i.e. $(\epsilon_{\a\beta}) = (\epsilon^{\a\beta})^{-1}$. We can see this as a way to "raise and lower indices" as 
\bse 
    \psi_{\a} = \epsilon_{\a\beta}\psi^{\beta}
\ese    
for left-handed spinors, and similarly for right-handed ones. 

\br 
    It is worth emphasising that our $\epsilon$s are 2-forms, \textit{not} metrics. In particular they are antisymmetric, while a metric is symmetric. Therefore we need to be careful about signs when raising and lowering indices. 
\er 

\bbox 
    By raising the spinor indices using the $\epsilon$s, show the following transformation behaviours
    \bse 
        \psi^{\a} \mapsto \psi^{\beta} {(M^{-1})_{\beta}}^{\a}
    \ese 
    and 
    \bse 
        \bar{\psi}^{\dot{\a}} \mapsto \bar{\psi}^{\dot{\beta}} {(M^{*-1})_{\dot{\beta}}}^{\dot{\a}}.
    \ese
\ebox  

\subsection{Scalar Product (anticommuting spinors)}

Now that we know how to raise/lower spinor indices, we can talk about taking scalar products of spinors. We emphasise before going forward, that the index placement is crucial to the definitions that follow. This is because we raise indices with a 2-form, and so contracting `top left to bottom right' is \textit{not} the same as contracting `bottom left top right'. This is easiest seen from the fact that that Weyl spinors are \textit{anti}commuting, and so their components are Grassman odd numbers, so 
\be 
\label{eqn:PsiChiMinusChiPsi}
    \psi^{\a}\chi_{\a} = - \chi_{\a}\psi^{\a}. 
\ee 

Ok first we construct a scalar product on left-handed spinors simply as `top right to bottom left' contraction
\mybox{
    \be  
    \label{eqn:LeftHandedInnerProduct}
        \psi \chi := \psi^{\a} \chi_{\a} = \epsilon^{\a\beta} \psi_{\beta}\chi_{\a}.
    \ee 
}

If we want this to be an inner product, we want it to be symmetric, i.e.
\bse 
    \psi \chi = \chi \psi.
\ese 
At first this seems in contrast to \Cref{eqn:PsiChiMinusChiPsi}, however this is where our convention of how to the contraction becomes important. We note that 
\bse 
    \psi\chi := \psi^{\a}\chi_{\a} = \epsilon^{\a\beta}\psi_{\beta}\chi_{\a}= - \epsilon^{\a\beta} \chi_{\a}\psi_{\beta} = + \epsilon^{\beta\a} \chi_{\a}\psi_{\beta} = \chi^{\beta}\psi_{\beta} =: \chi\psi,
\ese 
where we have explicitly used the 2-form nature, $\epsilon^{\a\beta}=-\epsilon^{\beta\a}$. Again we stress that the inner product is only symmetric because of the way we define our contractions. 

Ok so what about right-handed spinors? Here we impose the \textit{opposite} contraction convention. That is we contract `bottom left to top right':
\mybox{
    \be
    \label{eqn:RightHandedInnerProduct}
        \bar{\psi}\bar{\chi} := \bar{\psi}_{\dot{\a}} \bar{\chi}^{\dot{\a}} = \bar{\psi}_{\dot{\a}}\bar{\chi}_{\dot{\beta}}\epsilon^{\dot{\beta}\dot{\a}}.
    \ee
}

\noindent This might seem like a strange thing to do, but this convention has the following nice property:
\bse 
    (\psi\chi)^{\dagger} = \big(\psi^{\a}\chi_{\a}\big)^{\dagger} = (\chi_{\a})^{\dagger}(\psi^{\a})^{\dagger} = \bar{\chi}_{\dot{\a}} \bar{\psi}^{\dot{\a}} = \bar{\chi}\bar{\psi} = \bar{\psi}\bar{\chi},
\ese 
so the two inner products are related by Hermition conjugation. 

The $\sig_{\mu}, \bar{\sig}^{\mu}$ matrices we introduced before carry one dotted and one un-dotted index. This can be seen by the putting the spinor indices on our formula for ${\Lambda^{\mu}}_{\nu}(M)$, \Cref{eqn:Lambda(M)}, 
\bse 
    {\Lambda^{\mu}}_{\nu}(M) = \frac{1}{2}\tr[ M \sig_{\nu} M^{\dagger} \overline{\sig}^{\mu}] = \frac{1}{2} {M_{\a}}^{\beta} (\sig_{\nu})_{\beta\dot{\beta}} {(M^{\dagger})^{\dot{\beta}}}_{\dot{\g}} (\bar{\sig}^{\mu})^{\dot{\g}\a}. 
\ese 
Note the spinor index placement: the un-barred $\sig_{\mu}$ has lower indices with the un-dotted index appearing first, while the barred $\bar{\sig}^{\mu}$ has upper indices with the dotted index first. This is forced upon us as it is the only way to contract the indices in the above formula. We therefore have
\mybox{
    \be 
    \label{eqn:SigmaMuWithSpinorIndices}
        (\sig^{\mu})_{\a\dot{\a}} = (\b1, -\sig_i)_{\a\dot{\a}}, \qand (\overline{\sig}^{\mu})^{\dot{\a}\a} = (\b1, +\sig^i)^{\dot{\a}\a}
    \ee
}
\noindent where we note that the $-\sig_i$ in the first expression comes from raising the $\mu$ index with $\eta^{\mu\nu}=\text{diag}(+,-,-,-)$. 

\br 
\label{rem:SigmaMapping}
    There is a nice way to remember how the spinor index structure appears on our $\sig^{\mu}$ and $\bar{\sig}^{\mu}$. We simply note that they essentially act as maps\footnote{Note that the $(\sig^{\mu})_{\a\dot{\a}}$ acts on a right-handed spinor \textit{with} an $\epsilon^{\dot{\a}\dot{\beta}}$ as it needs a raised index. The map idea is still the same, though.}  
    \bse 
        (\sig^{\mu})_{\a\dot{\a}} : (0,1/2) \to (1/2,0) \qand (\bar{\sig}^{\mu})^{\dot{\a}\a} : (1/2,0) \to (0,1/2),
    \ese 
    and so can remember that the un-barred one has the undotted index first and similarly the barred one has the dotted index first. As for `up/down' placement, we just remember that the unbarred one is defined with a lower index $\sig_{\mu}$ and tell ourselves the spinor indices match. 
\er 

This placement of indices seems a bit of nuisance, however they are actually quite nice because it allows us to relate the two $\sig^{\mu}$ and $\bar{\sig}^{\mu}$ as follows
\be 
\label{eqn:BarSigmaRelationSigma}
    (\bar{\sig}^{\mu})^{\dot{\a}\a} = \epsilon^{\dot{\a}\dot{\beta}}\epsilon^{\a\beta} (\sig^{\mu})_{\beta\dot{\beta}}. 
\ee 
There is another really nice thing about the spinor index placement on our $\sig/\bar{\sig}$: it can be placed into our inner products easily. In particular 
\bse 
    \psi \sig^{\mu} \bar{\chi} = \psi^{\a} (\sig^{\mu})_{\a\dot{\a}} \bar{\chi}^{\dot{\a}}  \qand \bar{\psi} \bar{\sig}^{\mu} \chi = \bar{\psi}_{\dot{\a}} (\bar{\sig}^{\mu})^{\dot{\a}\a} \chi_{\a}.
\ese 
We therefore often reduce the notation and simply write the above as 
\bse 
    \psi \sig^{\mu}\bar{\chi} = \psi^{\a}\big(\sig^{\mu}\bar{\chi}\big)_{\a} \qand \bar{\psi} \bar{\sig}^{\mu}\chi = \bar{\psi}_{\dot{\a}}\big(\sig^{\mu} \chi\big)^{\dot{\a}},
\ese
which is just using the mapping idea from \Cref{rem:SigmaMapping}. Now if we consider the transformation behaviour of our $\psi\sig^{\mu}\bar{\chi}$ and $\bar{\psi}\bar{\sig}^{\mu}\chi$, we see that they transform as 4-vector, hence the $\mu$ index we've been using all along. 

\subsection{Dirac Spinors}

Above we have talked specifically about Weyl/Chiral spinors, but these are of course not the only kind of spinor. Another important representation is the
\bse
    \text{Dirac Spinor} \qquad (1/2,0) \oplus (0,1/2).
\ese 
These correspond to $4$-component Dirac spinors, which we conventionally denote 
\bse 
    \Psi_D = \begin{pmatrix}
        \psi_{\a} \\
        \bar{\chi}^{\dot{\a}}
    \end{pmatrix}.
\ese 
This matrix form follows simply from the definition of the direct sum $\oplus$, and we see that 
\bse 
    \begin{pmatrix}
        \psi_{\a} \\
        0
    \end{pmatrix} \qand \begin{pmatrix}
        0 \\
        \bar{\chi}^{\dot{\a}}
    \end{pmatrix}
\ese
are left-handed and right-handed Weyl spinors. 

We see from the above that the Dirac representation is \textit{not} an irrep; it is given by the direct sum of two irreps. Why is it interesting, then? Well we introduce the famous gamma matrices
\bse 
    \g^{\mu} = \begin{pmatrix}
        0 & \sig^{\mu} \\
        \bar{\sig}^{\mu} & 0
    \end{pmatrix}
\ese 
and then define the \textit{Chirality gamma matrix}
\bse 
    \g^5 := i \g^0 \g^1 \g^2 \g^3 = \begin{pmatrix} 
        \b1_2 & 0 \\
        0 & -\b1_2
    \end{pmatrix}. 
\ese
Therefore the Dirac spinor with only the left-handed components (i.e. $\bar{\chi}^{\dot{\a}}=0$) have Chirality\footnote{Eigenvalue of $\g^5$.} $+1$. Similarly the ones with only right-handed have Chirality $-1$. A Dirac spinor let's us package these into one object that we can manipulate at once. This finds tremendous use in the SM.

\bbox 
    Prove that the gamma matrices satisfy a \textit{Clifford algebra}
    \be 
    \label{eqn:CliffordAlgebra}
        \{\g^{\mu},\g^{\nu}\} = 2\eta^{\mu\nu}\b1_4.
    \ee 
\ebox 

\subsection{Majoranna Spinor}

There is one more important type of spinor worth mentioning, known as \textit{Majoranna} spinors. These are Dirac spinors that satisfy $\bar{\chi}= \bar{\psi} = \psi^{\dagger}$. We often summarise this as "Majoranna spinors are their own antiparticle".

\section{Lorentz Generators} 

We finish off this review by recalling how the Lorentz generators can be expressed in terms of our gamma matrices. They are simply given by\footnote{The motivation for the notation $\Sigma^{\mu\nu}$ is hopefully reasonably clear.}
\bse 
    M^{\mu\nu} \equiv \Sigma^{\mu\nu} = \frac{i}{4}[\g^{\mu}, \g^{\nu}] = i \begin{pmatrix}
        \sig^{\mu\nu} & 0 \\
        0 & \bar{\sig}^{\mu\nu} 
    \end{pmatrix}
\ese
where 
\mybox{
    \be 
        ({\sig^{\mu\nu})_{\a}}^{\beta} = \frac{1}{4}(\sig^{\mu}\bar{\sig}^{\nu} - \sig^{\nu}\bar{\sig}^{\mu}{\big)_{\a}}^{\beta} \qand ({\bar{\sig}^{\mu\nu})^{\dot{\a}}}_{\dot{\beta}} = \frac{1}{4}(\bar{\sig}^{\mu}\sig^{\nu} - \bar{\sig}^{\nu}\sig^{\mu}{\big)^{\dot{\a}}}_{\dot{\beta}}
    \ee 
}

Using the results 
\bse 
    (\g^0)^{\dagger} = \g^0 \qand (\g^i)^{\dagger} = -\g^i,
\ese 
we get 
\bse 
    \begin{split}
        \text{Boosts} \qquad  \big(\Sigma^{0i}\big)^{\dagger} & = - \Sigma^{0i}, \\
        \text{Rotations} \qquad \big(\Sigma^{ij}\big)^{\dagger} & = + \Sigma^{ij}.
    \end{split}
\ese 
We also have 
\bse
    i \sig^{12} = i\bar{\sig}^{12} = \frac{1}{2}\sig^3  \qquad \iff \qquad J_3 = S_3,
\ese 
and so 
\bse 
    J_3 \begin{pmatrix}
        \psi_1 \\
        \psi_2
    \end{pmatrix} = \begin{pmatrix}
        \frac{1}{2}\psi_1 \\
        -\frac{1}{2}\psi_2
    \end{pmatrix}, \qand J_3 \begin{pmatrix}
        \bar{\chi}^{\dot{1}} \\
        \bar{\chi}^{\dot{2}}
    \end{pmatrix} = \begin{pmatrix}
        \frac{1}{2}\bar{\chi}^{\dot{1}}\\
        -\frac{1}{2}\bar{\chi}^{\dot{2}}
    \end{pmatrix}.
\ese 
\chapter{Supersymmetry Algebra \& Supermultiplets}

Ok so we have reviewed the Lorentz group and the spin group, we now want to go on to add all the "super" prefixes to things. The first thing we need to do is generalise the idea of a Lie algebra and then talk about the representations of such a superalgebra. 

\br 
    We will focus on $4$-dimensional SUSY, but the logic will apply elsewhere. The only difference is that the representations (i.e. the spinors) will change, but the logic is exactly the same. 
\er

\section{Lie Superalgebra (Graded Lie Algebra Of Degree 1)}

As we just said, one of the most important things for us to generalise is the notion of a Lie algebra to what is known as a \textit{graded} Lie algebra. Such objects are their own beasts and extend beyond SUSY, but the simplest kind --- a graded Lie algebra \textit{of degree 1} --- finds massive application in SUSY and is known as a \textit{Lie Superalgebra}.

So how do we make a Lie superalgebra? Well the first thing we have to do is extend the notion of a vector space (which a Lie algebra is) to a \textit{graded vector space}. For SUSY in particular, we want what is known as a \textit{$\Z_2$ graded vector space} or more simply a \textit{super vector space}, which we define as follows.
\bd[$\Z_2$ Graded Vector Space/Super Vector Space]
    A \textit{super vector space} is a vector space $V$, which can be written as a decomposition 
    \be 
        V = V_0 \oplus V_1
    \ee 
    where we call any vector that is purely an element of $V_0$ or purely an element of $V_1$ \textit{homogeneous}. We associate a \textit{parity} to the homogeneous vectors as follows 
    \bse 
        |v| = \begin{cases}
            0 & \text{if } v\in V_0 \\
            1 & \text{if } v\in V_1.
        \end{cases}
    \ese 
    A homogeneous vector of parity $0$ is called \textit{even/Bosonic}, while one with parity $1$ is called \textit{odd/Fermionic}. The addition and $\C$-multiplication\footnote{Or more generally $\F$ for some underlying field $\F$.} are simply inherited component wise.\footnote{That is for $v,w\in V_0$ and $p,q\in V_2$, we have $(v,p)+(w,q) = (v+w,p+q)\in V_0\oplus V_1$ etc.}
\ed 

We can show how operators such as direct sums and direct products carry over to super vector spaces however these details are omitted here.\footnote{For a nice discussion, see Section 2.1 of David Skinner's notes on SUSY.}

Ok now that we have a super vector space, we can now try equip it with some kind of Lie bracket structure to give us a Lie superalgebra. 

\bd[Lie Superalgebra]
    Let $L=L_0\oplus L_1$ be a super vector space. We can make this into a \textit{Lie superalgebra} by equipping it with a bilinear bracket 
    \bse 
        [\cdot, \cdot\} : L \times L \to L,
    \ese 
    obeying: for all $x_i\in L_i$ and $x_j \in L_j$
    \ben[label=(\roman*)] 
        \item (Grading Consistency): $[ x_i , x_j \} \subseteq L_{i+j}$  
        \item ((Anti)-Symmetry): $[x_i , x_j\} = -(-1)^{ij} [x_i , x_j\}$.
        \item (Jacobi) :
            \bse 
                (-1)^{ik} \big[x_i , [x_j , x_k\} \big\} + (-1)^{kj} \big[x_k , [x_i , x_j\} \big\} + (-1)^{ji} \big[x_j , [x_k , x_i\} \big\} = 0. 
            \ese 
    \een 
\ed 

Let's break this down a little bit. From the above properties, we can see that
\ben[label=(\roman*)]
    \item the even part is closed under the bracket, and it is indeed itself a Lie algebra as 
    \bse 
        [L_0, L_0\} = [L_0,L_0],
    \ese 
    where the right hand side is understood to be a Lie bracket, specifically the commutator (as we will consider matrix groups).
    \item The odd part is \textit{not} closed under the bracket as 
    \bse 
        [L_1, L_1\} \subseteq L_0
    \ese 
    and the bracket becomes the \textit{anti}commutator.
    \item Finally we have 
    \bse 
        [L_0, L_1\} \subseteq L_1
    \ese
    with the bracket being the commutator. 
\een 

\br 
    Condition (iii) above actually tells us that the odd part is a representation space for the even part, with the representation given by an adjoint-type action.
\er 

\br 
    Again we emphasise that we have focused specifically on the case relevant to SUSY. We can extend the above definitions to more general graded Lie algebras \textit{of degree $n$}, where the vector space is given by 
    \bse 
        V = \bigoplus_{i=0}^n L_i.
    \ese
    We then get our Lie superalgebra by imposing a $\Z_2$ quotient $L_i = L_{i+2}$, which is where the names above came from. 
\er 

\section{4D SUSY (Or Superpoincar\'{e}) Algebra}

As the naming in our definition of a super vector space suggested, the idea is that the even objects are Bosonic while the odd objects are Fermionic. We have 
\ben[label=(\roman*)]
    \item $L_0$ is the Poincar\'{e} algebra (plus any potential R symmetry and central $U(1)$ symmetries)
    \item $L_1$ is the supercharges $Q^I_{\a}$ and $\bar{Q}_{I\dot{\a}} = (Q_{\a}^I)^{\dagger}$, where $I= i,..., \cN$ takes care of potentially having multiple different SUSYs.
\een
In addition to the commutators of the Poincar\'{e} algebra, \Cref{eqn:PoincareCommutator}, we have:
\mybox{
\be 
    \label{eqn:SuperPoincareCommutators}
        \begin{split}
            [P_{\mu}, Q_{\a}^I] = 0 = [P_{\mu}, \bar{Q}_{I\dot{\a}}] \\
            [M_{\mu\nu}, Q^I_{\a}] = i {(\sig^{\mu\nu})_{\a}}^{\beta} Q_{\beta}^I \\ 
            [M_{\mu\nu}, \bar{Q}_U^{\dot{\a}}] = i {(\bar{\sig}^{\mu\nu})^{\dot{\a}}}_{\dot{\beta}} \bar{Q}^{\dot{\beta}}_I \\
            \{ Q^I_{\a}, \bar{Q}_{J\dot{\beta}}\} = 2(\sig^{\mu})_{\a\dot{\beta}} P_{\mu} \del^I_J \\
            \{ Q_{\a}^I, Q_{\beta}^J \} = \epsilon_{\a\beta} Z^{IJ} \\
            \{ \bar{Q}_{I\dot{\a}}, \bar{Q}_{J\dot{\beta}}\} = \epsilon_{\dot{\a}\dot{\beta}} \bar{Z}_{IJ}
        \end{split}
    \ee 
    where 
    \be 
    \label{eqn:CentralCharges}
        Z^{IJ}=-Z^{JI} \qand \bar{Z}_{IJ} = (Z^{IJ})^{\dagger}
    \ee 
}
\noindent We call the $Z/\bar{Z}$ the \textit{central charges}, as commute with all the generators.\footnote{Hence they are in the centre of the group.} We refer to this set of commutators/anticommutators as the \textit{superPoincar\'{e} algebra}.

In addition, there \textit{might} be a $R$-symmetry $\subseteq U(\cN)$, which acts on indices $I,J$ of supercharges such that $Q_{\a}^I$ are in the fundamental rep of $U(\cN)$ and $\bar{Q}_{I\dot{\a}}$ in the antifundamental.\footnote{Note this is consistent with the convention of upper/lower indices for fundamental/antifundamental.} If the central charges vanish, the $R$-symmetry is an automorphism of the above commutators.\footnote{This is essentially why we took the semi-direct product in \Cref{thrm:HLS} with the SuperPoincar\'{e} group being the normal subgroup.} Whether this $R$-symmetry is realised as an actual symmetry of the theory or not depends on the theory you are considering (in particular on the central charges and on the interactions). 

\br 
    Note that for $\cN=1$ we could have the R-symmetry as then, $I,J=1$ only and so \Cref{eqn:CentralCharges} can only be satisfied if the central charges vanish. However we stress that we don't \textit{have} to have an $R$-symmetry even in this case, as some interactions could break this symmetry. In other words 
    \begin{center}
        having $R$-symmetry $\implies$ central charges vanish,
    \end{center}
    but the reverse is \textit{not} true. 
\er 

\bnn 
    From now on we shall often drop the parentheses around $(\sig^{\mu})_{\a\dot{\a}}$ and $(\bar{\sig}^{\mu})^{\dot{\a}\a}$ in order to lighten notation. That is we will just write $\sig^{\mu}_{\a\dot{\a}}$ and $\bar{\sig}^{\mu\dot{\a}\a}$.
\enn 

We are now in a little better place to understand the Haag-Lopuszanski-Sohnius theorem, \Cref{thrm:HLS}. The idea is that the superPoincar\'{e} algebra, \Cref{eqn:SuperPoincareCommutators}, is completely fixed by requiring
\ben[label=(\roman*)] 
    \item Poincar\'{e} symmetry,
    \item Extension to a graded Lie algebra (of degree 1), and 
    \item Coleman-Mandula theorem (rules out conserved charges of spin >1). 
\een

\subsection{Basic Consequences of SUSY Algebra}

Let's now look at some of the basic consequences of our SUSY algebra. 

\ben 
    \item If we have unbroken SUSY then it follows from $[P^2, Q]=0=[P^2, \bar{Q}]$ along with Schur's Lemma that superpartners have the same mass.
    \item Using $\a=\dot{\a}$, in the sense that $\a=1 \iff \dot{\a} =1$\footnote{We need this for when we take the Hermitian conjugates below.} etc, we have (the indices $I$ are \textit{not} summed over)
    \bse
        0 \leq \| Q_{\a}^I \ket{\phi}\|^2 + \| \bar{Q}_{I\dot{\a}}\ket{\phi}\|^2 = \bra{\phi} \bar{Q}_{I\dot{\a}} Q^I_{\a} \ket{\phi} + \bra{\phi} Q^I_{\a} \bar{Q}_{I\dot{\a}}\ket{\phi} = 2 (\sig^{\mu})_{\a\dot{\a}} \bra{\phi}P_{\mu}\ket{\phi},
    \ese 
    where we have used the anticommutator relation above. Now sum over $\a=\dot{\a}$, i.e. take the trace, to obtain
    \bse 
        0 \leq  2 \tr[\sig^{\mu}] \bra{\phi} P_{\mu} \ket{\phi} = 4\bra{\phi} P_0 \ket{\phi} = 4E_{\phi},
    \ese 
    so we see the energy is non-negative and that $E_{\phi} = 0$ iff $Q^I_{\a}\ket{\phi}= 0 = \bar{Q}_{I\dot{\a}}\ket{\phi}$ for all $I,\a,\dot{\a}$. This allows us to conclude 
    \ben 
        \item Energy $E \geq 0$ for \textit{all} physical states\footnote{Note this means we can't just shift the energy around as we normally do.} 
        \item The equality is only satisfied for the \textit{supersymmetric ground state}. 
    \een 
    These are robust statements that are protected against quantum corrections, as they follows from the algebra \textit{alone}. Spontaneous \cancel{SUSY} occurs if and only if the vacuum energy is greater than zero. This is a brilliant way to check for SUSY, namely if you measure the lowest energy state of the system and it's non-zero you \textit{must} have spontaneous \cancel{SUSY}. Similarly if the lowest energy state is zero you must have unbroken SUSY.\footnote{The important thing to note is that if must be the actual lowest state of the system. That is our measured non-vanishing of the Universe vacuum energy only tells us that we have spontaneously broken SUSY \textit{if} we assume that the Universe's vacuum is the lowest state of the system.} 
    \item A supermulitplet (not vacuum) contains an equal number of Bosonic and Fermionic states, $n_B = n_F$. This is a non-trivial statement and we shall now prove it.
    \bq 
        We start by defining the \textit{Fermion number operator} $F$ which acts as
        \bse 
            \begin{split}
                \bra{b}F\ket{b} \qquad  \text{even} \\
                \bra{f}F\ket{f} \qquad \text{odd}
            \end{split}
        \ese 
        where $\ket{b}/\ket{f}$ are Boson/Fermion states, respectively. In particular, we could write it as $F=2s$ where $s$ is the spin. Therefore we have 
        \bse 
            (-1)^F \ket{b} = + \ket{b}, \qand (-1)^F \ket{f} = -\ket{f}.
        \ese
        Since $Q/\bar{Q}$ shift spin by $1/2$, it changes the statistics, thus
        \be 
        \label{eqn:FOnQ}
            (-1)^F \widetilde{Q} = - \widetilde{Q}(-1)^F,
        \ee
        which is equivalent to saying that the supercharges are Fermionic.
        
        Ok now let's consider a supermultiplet.\footnote{That is a finite-dimensional irrep of SUSY algebra} Now by the cylictity of the trace, we have\footnote{Note we use a capital $\Tr$ here to differentiate the trace over the supermultiplet from the trace over spinor indices we had before which we denoted as a lower case $\tr$.}
        \bse 
            \Tr[(-1)^F\bar{Q}_{J\dot{\beta}}Q^I_{\a}] = \Tr[Q^I_{\a}(-1)^F\bar{Q}_{J\dot{\beta}}],
        \ese 
        and so using \Cref{eqn:FOnQ}, we have
        \bse 
            \begin{split}
                0 & = \Tr[-Q_{\a}^I (-1)^F \bar{Q}_{J\dot{\beta}} + (-1)^F \bar{Q}_{J\dot{\beta}} Q^I_{\a}\} ]  \\
                & = \Tr[ (-1)^F \{Q^I_{\a},\bar{Q}_{J\dot{\beta}}\} ] \\
                & = 2 (\sig^{\mu})_{\a\dot{\beta}} \del^{I}_J \Tr[(-1)^F P_{\mu}] \\
                & = 2 (\sig^{\mu})_{\a\dot{\beta}} \del^I_J p_{\mu} Tr[(-1)^F],
            \end{split}
        \ese 
        where we have made use of the superPoincar\'{e} algebra, and where the little $p_{\mu}\in \R$ is the common eigenvalue of $P_{\mu}$ on the supermultiplet. Now if we choose $p_{\mu}\neq 0$ (that is $E\neq 0$, so \textit{not} the vacuum) we obtain
        \bse 
            0 = \Tr[(-1)^F] = n_B - n_F.
        \ese
    \eq 
\een 

\section{Supermultiplets}

As we have said a few times, supermultiplets are irreps of the superPoincar\'{e} algebra. As the Poincar\'{e} algebra is a subalgebra of the superPoincar\'{e} algebra, we see that supermultiplets are (generally reducible) representations of the Poincar\'{e}. This tells us that they contain particles, and if the Poincar\'{e} rep is reducible, the supermultiplet actually contains multiple different types of particle. They are, by definition, \textit{on-shell}. We will get an \textit{off-shell} generalisation later when we introduce superfields. Let's now discuss these supermultiplets further.

As the supermulitplets contain particles, we can split them into massless and massive reps, which we will now discuss.

\subsection{Massless Supermultiplets}

Here we can go to a light-cone frame so that 
\bse 
    P_{\mu} = E(1,0,0,1)
\ese 
from which a quick calculation gives
\bse 
    \sig^{\mu}P_{\mu} = \begin{pmatrix}
        0 & 0 \\
        0 & 2E
    \end{pmatrix}.
\ese
Then using our superPoincar\'{e} algebra relation 
\bse 
    \{Q_{\a}^I, \bar{Q}_{J\dot{\beta}}\} = 2 \sig^{\mu}_{\a\beta} P_{\mu} \del^I_J,
\ese 
we have
\bse 
    \{Q_{\a}^I, \Bar{Q}_{J\dot{\a}} \} = \begin{pmatrix}
        0 & 0 \\
        0 & 4E
    \end{pmatrix} \del^I_J \qquad \implies \qquad \{Q_1^I , \bar{Q}_{J\dot{1}}\} = 0 \quad \forall I,J \in 1,...,\cN.
\ese
From this we have
\bse 
    0 = \bra{\phi} \{Q_1^I , \bar{Q}_{J\dot{1}}\} \ket{\phi} = \|\bar{Q}_{I\dot{1}}\ket{\phi}\|^2 + \|Q_1^I\ket{\phi} \|^2,
\ese 
and so if we assume unitarity (or non-negative definitness of our Hilbert space of states) we conclude 
\bse 
    \bar{Q}_{I\dot{i}}\ket{\phi} = 0 = Q_1^I\ket{\phi},
\ese 
for \textit{all} physical states in our supermultiplet. The only way we can satisfy this is if we have 
\bse 
    Q_1^I = 0 = \bar{Q}_{I\dot{1}}.
\ese 
In particular this gives us 
\bse 
    Z^{IJ}\ket{\phi} = \bar{Z}_{IJ}\ket{\phi} = 0 
\ese 
on this supermultiplet.

This tells us that $\cN$ of our total $2\cN$ (2 from $\a$ and $\cN$ from $I$) supercharges act trivially on the supermultiplet. What about the remaining $\cN$? Well, we define\footnote{Note the denominators are included to counteract the $4E$ above.} 
\mybox{
    \be 
    \label{eqn:CreationAnnihilationMassless}
        a^I := \frac{1}{2\sqrt{E}} Q_2^I, \qand a_I^{\dagger} := \frac{1}{2\sqrt{E}} \bar{Q}_{I\dot{2}},
    \ee 
}
\noindent which satisfy
\bse 
    \{ a^I, a_J^{\dagger}\} = \del^I_J, \qquad \{ a^I,a^J \} = 0 = \{ a_I^{\dagger},a_J^{\dagger} \},
\ese 
but these are just the anticommutation relations for the creation/annihilation operators for Fermions! We can see how they affect the helicity by computing
\bse 
    [M_{12} , a^I] = -\frac{1}{2}a^I \qand [M_{12}, a_J^{\dagger}] = \frac{1}{2}a_J^{\dagger},
\ese 
which tells us that $a^I$ lowers the helicity by $1/2$ while $a_J^{\dagger}$ raises it by $1/2$. 

Ok so we have our raising/lowering operators, so we can now try to build up our representation. As always we do this by acting on the vacuum. We define the \textit{Clifford vacuum} as $\ket{\l_0}$ which has helicity $\l_0$ and is annihilated by all $a^I$,
\bse 
    a^I \ket{\l_0} = 0.
\ese 
We now act with the $\cN$ Fermionic creation operators $a_I^{\dagger}$ and produce the states
\bse
    \begin{split}
        \ket{\l_0} & \\
        a_I^{\dagger} \ket{\l_0} & = \ket{\l_0+\frac{1}{2}}_I \\
        a_I^{\dagger} a_J^{\dagger} \ket{\l_0} & = \ket{\l_0+1}_{[IJ]}\\
        & \vdots \\
        a_{I_1}^{\dagger} ... a_{I_{\cN}}^{\dagger}\ket{\l_0} & = \ket{\l_0 +\frac{\cN}{2}}_{[I_1...I_{\cN}]}
    \end{split}
\ese 
where the subscripts on the states remind us that we have antisymmetric operators, i.e. if we act with the same operator twice the state vanishes. The most general state is given by 
\bse 
    a_{I_1}^{\dagger} ... a_{I_k}^{\dagger}\ket{\l_0} = \ket{\l_0 +\frac{k}{2}}_{[I_1...I_k]}.
\ese
We can use this to work out how many different states there are of a given helicity. With the comment above about not being able to apply the same operator twice, it's clear that there are\footnote{Read "$\cN$ choose $k$".}
\bse 
    \begin{pmatrix}
        \cN \\
        k
    \end{pmatrix} := \frac{\cN!}{k!(\cN-k)!},
\ese
with helicity $\l_0 + \frac{k}{2}$. We can then use this to work out the \textit{total} number of states in the supermultiplet. Using the binomial formula 
\bse 
    (x+y)^n = \sum_{k=1}^n \begin{pmatrix}
        n \\
        k
    \end{pmatrix} x^k y^{n-k},
\ese 
we have 
\bse 
    (\text{total \# states}) = \sum_{k=0}^{\cN} \begin{pmatrix}
        \cN \\
        k
    \end{pmatrix} = \sum_{k=0}^{\cN} \begin{pmatrix}
        \cN \\
        k
    \end{pmatrix} 1^{k}1^{\cN-k} = (1+1)^{\cN} = 2^{\cN},
\ese 
where the second equality follows by trivially multiplying by $1$. So our massless supermultiplet has $2^{\cN}$ states. This is obviously a lot of states, but it is still considerably less then we had before ariving at \Cref{eqn:CreationAnnihilationMassless}. On top of this, we wont take $\cN$ to be huge, so we don't need to be too scared of this result.

\br 
    We can also check that we still have $n_F=n_B$ by computing
    \bse 
        \begin{split}
            \Tr(-1)^F & = \sum_{k=0}^{\cN} \begin{pmatrix}
                \cN \\
                k
            \end{pmatrix} (-1)^{2\l_0+k} \\
            & = (-1)^{2\l_0} \sum_{k=0}^{\cN} \begin{pmatrix}
                \cN \\
                k
            \end{pmatrix} (-1)^k 1^{\cN-k} \\
            & = (-1)^{2\l_0} (-1+1)^{\cN} \\
            & = 0,
        \end{split}
    \ese 
    where we have used $F = 2|s| = 2(\l_0 + \frac{k}{2})$, and again cleverly multiplied by $1$ on the third line. 
\er

Ok great so we know how to build up our supermultiplet, but there's an important point we haven't covered yet, which is the content of the next proposition. 

\bp 
    Any unitary, locally Lorentz invariant QFT must be CPT invariant. We call such theories \textit{self-conjugate}. 
\ep 

Why do we care about this? Well because a CPT operation flips the helicity of a particle, i.e. 
\bse 
    CPT : \l \mapsto -\l,
\ese 
and there is no reason why our supermultiplet constructed above should respect this symmetry. Indeed such a constructed supermultiplet typically will \textit{not} be self-conjugate, and in particular we will only get a self-conjugate theory if 
\be 
\label{eqn:MasslessSelfConjugateCondition}
    \l_0 = -\frac{\cN}{4}.
\ee
We can see this reasonably easily by considering "pairing off" the states into self conjugate pairs: 
\bse 
    \begin{split}
        \ket{\l_0} & \longleftrightarrow \ket{\l_0 + \frac{\cN}{2}} \\
        \ket{\l_0+ \frac{1}{2}} & \longleftrightarrow \ket{\l_0 + \frac{\cN-1}{2}} \\
        & \quad  \vdots \\
        \ket{\l_0 + \frac{\cN}{4} + \frac{1}{2}} & \longleftrightarrow \ket{\l_0 + \frac{\cN}{4} - \frac{1}{2}} \\
        \ket{\l_0 + \frac{\cN}{4}} & \longleftrightarrow \ket{\l_0 + \frac{\cN}{4}},
    \end{split}
\ese 
we can only satisfy the last condition if 
\bse 
    \l_0 + \frac{\cN}{4} = 0,
\ese
which is exactly \Cref{eqn:MasslessSelfConjugateCondition}.

\br 
    We can also see \Cref{eqn:MasslessSelfConjugateCondition} by noting that under a CPT transformation our raising/lowering operators essentially flip roles, and so out highest weight state $\ket{\l_0 + \cN/2}$ and lowest weight state $\ket{\l_0}$ also flip. If these are going to be invariant, it's clear we need \Cref{eqn:MasslessSelfConjugateCondition}.
\er 

Of course in general we are not going to satisfy such a condition, and when we don't we will have to restore CPT invariance by adding the CPT-conjugate states. This gives us a supermultiplet with $2^{\cN+1}$ states in total. The particles that are identified under a CPT transformation have opposite helicity and charge, and so are particle-antiparticle pairs. 

Before considering specific examples, let's introduce the notation we are going to use and what it means. We will denote the content of the supermultiplet by listing the allowed helicity values. This is clearly a \textit{set} and so we should use the notation 
\bse 
    \{ \l_0, \l_0 + 1/2 , ... , \l_0 + \cN/2\},
\ese    
however we will simply use a bracket notation 
\bse 
    (\l_0, \l_0 + 1/2, ... , \l_0 + \cN/2). 
\ese 
Now in the cases when we have to add the CPT-conjugate states, really we are taking the \textit{disjoint union}
\bse 
    (\l_0, ... , \l_0 + \cN/2 ) \, \dot{\cup} ( -\l_0 , ... , -\l_0 -\cN/2).
\ese 
It's a union simply because we're putting two sets together, and it is disjoint because, for example, if we start from $\l_0=0$ but don't have a self-conjugate system, then our CPT-conjugate state will also have a $\l=0$ entry. We need to keep track of both of these $0$s as they correspond to different particles.\footnote{Well more correctly, they correspond to different degrees of freedom. This will hopefully become clear shortly.} As well as this, if we have $\cN>1$, it's possible that we can produce the same helicity value in multiple ways \textit{within} the original set. Again these are separate degrees of freedom, and so we need to keep track of them all.

For the reasons above, here we will simply use a $+$ to denote the union and will adopt the, somewhat strange, notation of using a $\times$ and inner brackets when we have multiple particles with the same helicity. This will hopefully become clear with the examples that follow. 


\subsubsection{$\cN=1$ Massless Supermultiplets}

When we have $\cN=1$ we only have two helicity values, namely
\bse 
    (\l_0, \l_0 + 1/2).  
\ese 
Now, since the helicity is always half integer there is no way for this to be self-conjugate, and so we will always have to add the CPT conjugate states:
\bse 
    (\l_0, \l_0+1/2) + (-\l_0, -\l_0-1/2).
\ese 
We then categorise the different supermultiplets via the $\l_0$ value. We summarise all the possible combinations in the table below, which we shall explain below. 

\begin{center}
	\begin{tabular}{@{} C{1.5cm} C{4cm} C{4cm} C{5cm} @{}}
		\toprule
		$\l_0$ & Multiplet Name & Helicity Content & Particle Content \\
		\midrule 
		$0$ & (Chiral) $\chi$-plet & $(-1/2, \, 2\times (0), \,  1/2)$ & Complex Scalar \& Weyl Fermion \\ \\
		$1/2$ & (Vector) V-plet & $(-1, \, -1/2, \, 1/2, \, 1)$ & Gauge Boson \& Weyl Fermion \\ \\
		$1$ & Gravitino Multiplet & $(-3/2, \, -1, \, 1, \, 3/2)$ & Gauge Boson \& Gravitino \\ \\
		$3/2$ & Gravity Multiplet & $(-2,\, -3/2, \, 3/2, 2)$ & Graviton \& Gravitino \\ 
		\bottomrule
	\end{tabular}
\end{center}

So how did we construct this table? Well the first column is obviously just the $\l_0$ values, which we restrict to below $s\leq 2$ in accordance with condition (vi)(b) of \Cref{sec:WhatIsSUSY}. The next column is just the names of the given supermultiplets. Note that the last two are only valid when we have gravity (in accordance with (vi)(a) from \Cref{sec:WhatIsSUSY}), which is where their names comes from. The third column can be obtained give the procedure outlined above, for example for the $\chi$-plet we have 
\bse 
    (0,1/2) + (0,-1/2) = (-1/2, \, 2\times (0), \, 1/2). 
\ese 
The last column comes from us knowing how the helicity content given particles. We get these from considering the degrees of freedom of the particles and their spin. That is:
\ben[label=(\roman*)]
    \item Complex scalar is $2\times (0)$, 
    \item Weyl Fermion is $(-1/2,\, 1/2)$,
    \item Gauge Boson (which is a massless vector Boson) is $(-1, \, 1)$, where we don't have the $0$ part because it is massless,
    \item Graviton is $(-2,\, 2)$, i.e. a spin-$2$ massless Boson,
    \item We get the gravitino from the fact that it is the superpartner to the graviton, as per the last column. It is therefore $(-3/2, \, 3/2)$. 
\een 

\bter 
    It is common terminology to name the superpartner of a Boson by adding "ino" to end of the name, just like we did for the gravitino above. On the other hand for Fermions, the superpartners are named by putting an "s" in front of it, e.g. the superpartner of a top quark is a "stop squark". 
\eter 

\subsubsection{$\cN=2$ Massless Supermultiplets}

Let's now consider the $\cN=2$ case. Here we have two types of creation operators, so we have 
\bse 
    (\l_0, \, 2\times (\l_0+1/2), \, \l_0 + 1)
\ese 
usually this is \textit{not} self conjugate and so we need to add CPT-conjugate states. The helicity content here is longer and so we present them as a list rather then a table. 
\ben[label=(\roman*)] 
    \item $\cN=2$ V-plet: $\l_0=0$:
    \bse 
        (-1, \, 2\times (-1/2), \, 2\times(0), \, 2\times(1/2), \, 1).
    \ese 
    We have 2 Weyl spinors, a gauge boson and a complex scalar. However we note these are all part of the $\cN=1$ above so we can decompose it into one $\cN=1$ V-plet and one $\cN=1$ $\chi$-plet.
    \item $\cN=2$ half-hyper multiplet, $\frac{1}{2}$H-plet: $\l_0=-1/2$ we start with
    \bse 
        (-1/2, 2\times(0), 1/2),
    \ese
    which we note is already potentially self-conjugate, because it is a $\cN=1$ $\chi$-plet. Technically this is only possible if the representation of this multiplet is so-called \textit{pseudo-real}.
    \item $\cN=2$ H-plet: This is the above but with CPT conjugate added, so we have 
    \bse 
        (2\times(-1/2), 4 \times(0), 2\times (1/2))
    \ese 
    which is two lots of $\cN=1$ $\chi$-plet.
\een 

\bbox 
    Construct the $\cN=2$ gravitino ($\l_0=-3/2$) and graviton ($\l=-2$) multiplets. By looking at the particle content, find the decompositions in terms of $\cN=1$ multiplets. 
    \br 
        The result of this plus another calculation should explain why we have `skipped' the $\l_0=-1$ multiplet. 
    \er 
\ebox 

\subsubsection{$\cN=4$ Massless Supermultiplets}

Here we have $4$ super charges, and so we have 
\bse 
    (\l_0, \, 4\times (\l_0+1/2), \, 6\times(\l_0+1), \, 4 \times (\l_0+3/2), \, \l_0+2).
\ese 
It follows from this that if we don't have gravity, i.e. $s \leq 1$, that there is only one $\cN=4$ supermultiplet. It is a V-plet, $\l_0=-1$ which is self conjugate and given by
\bse 
    (-1, 4\times(-1/2), 6\times(0), 4\times(1/2), 1)
\ese
which we can decompose as a $\cN=2$ V-plet and a $\cN=2$ H-plet. 

This theory is known as \textit{$\cN=4$ super-Yang Mills} and corresponds to the most symmetric non-Abelian QFT in $4D$. It is the most symmetric theory we can have without studying supergravity (i.e. getting $\l_0>1$). 

\br 
    The $R$-symmetry here turns out to be $SU(4)$ rather then $U(4)$. This comes from the fact that the V-plet is self conjugate.
\er 

\bbox 
    Find the $\cN=4$ graviton multiplet, $\l_0=-2$. Decompose the result in terms of $\cN=2$ and $\cN=1$ multiplets.
\ebox 

\subsubsection{$\cN=3$ Massless Supermultiplets}

As you might have noticed, we skipped past $\cN=3$ massless supermultiplets. The reason we did it is the content of the next exercise.

\bbox 
    Using fermionic oscillators, construct explicitly the most general massless supermultiplet of particles in a supersymmetric QFT with $\cN = 3$ SUSY.
    \ben 
        \item What is the particle content of this multiplet?
        \item Show that this multiplet coincides with an $\cN = 4$ vector multiplet.
    \een 
\ebox 

\br 
    Apparently people have recently discovered $\cN=3$ super-CFTs which are \textit{not} $\cN=4$. However these turn out to have no Lagrangian description, so we wont discuss them here. 
\er 

\subsection{Massive Supermultiplets \& BPS States}

The other category for our particles, and therefore the supermultiplets, is obviously massive ones. Here we go to the rest frame
\bse 
    P_{\mu} = (m,0,0,0).
\ese
We can plug this into our superPoincar\'{e} algebra relation and obtain
\bse 
    \{Q_{\a}^I, \bar{Q}_{J\dot{\beta}}\} = 2m\del_{\a\dot{\beta}} \del^I_J.
\ese 
as well as 
\bse 
    \{ Q_{\a}^I , Q_{\beta}^J \} = \epsilon_{\a\beta}Z^{IJ} \qand  \{ \bar{Q}_{I\dot{\a}} , \bar{Q}_{J\dot{\beta}} \} = \epsilon_{\dot{\a}\dot{\beta}}\bar{Z}_{IJ}.
\ese 
By a $U(\cN)$ rotation, i.e. a $R$-symmetry transformation, we can skew-diagonalise $Z^{IJ}$ as
\bse 
    Z^{IJ} = \begin{pmatrix}
        0 & z_1 && && & \\
        - z_1 & 0  && && &\\ 
        & & 0 & z_2 && &\\
        & & -z_2 & 0 && &\\
        & & & & \ddots & &\\
        & & & & & 0 & z_{\cN/2} \\
        & & & & &  -z_{\cN/2} & 0
    \end{pmatrix}
\ese 
where $z_i \in \R$. This clearly only works if $\cN \in 2\Z$, however we can easily adapt it to odd $\cN$ values by treating the last one a $\cN=1$. As we will see in a moment, this corresponds to just putting another row/column of $0$s, i.e.
\bse 
    Z^{IJ} = \begin{pmatrix}
        0 & z_1 && && & & 0 \\
        - z_1 & 0  && && & & 0 \\ 
        & & 0 & z_2 && & & 0 \\
        & & -z_2 & 0 && & & 0 \\
        & & & & \ddots & & & \vdots \\
        & & & & & 0 & z_{(\cN-1)/2} & 0\\
        & & & & &  -z_{(\cN-1)/2} & 0 & 0 \\
        0 & 0 & 0 & 0 & \dots & 0 & 0 & 0
    \end{pmatrix}
\ese 

\br 
    In the rest frame $SO(1,3)$ is broken to $SO(3)$ and similarly $SL(2,\C)$ to $SU(2)$. This is why we have the $\a$ and $\dot{\beta}$ "talking to each other" in $\del_{\a\dot{\beta}}$ above. In other words, dotting an index no longer means anything. However, we will still use the dotted notation so we can remember where the indices came from.
\er 

We can't make any $\cN$ independent comment about the available states here as the central charges don't vanish. That is, for the massless case we could define the creation/annihilation operators \textit{for any} $\cN$ value and then proceeded from there. Here we have to go case by case.  

\subsubsection{$\cN=1$ Massive Supermultiplets}

Of course first we consider $\cN=1$ case. Now from the antisymmetry condition, $Z^{IJ}=-Z^{JI}$, with $I,J=1$ only, we see $Z=0$. We can therefore define 
\mybox{
    \be 
    \label{eqn:CreationAnnihilationMassive}
        a_{\a} = \frac{1}{2m} Q_{\a}, \qquad a_{\dot{\a}}^{\dagger} = \frac{1}{\sqrt{2m}} \bar{Q}_{\dot{\a}} \qquad \a,\dot{\a} = 1,2
    \ee 
}
\noindent where the important thing to note is that $\a,\dot{\a}$ now take two values. This is different to the massless case where we only had $\a,\dot{\a}=2$. We therefore have twice as many Fermionic oscillators compared to the massless case. 

Now how do \Cref{eqn:CreationAnnihilationMassive} effect the spin?\footnote{Note we talk about \textit{spin} not helicity, as the particles are massive and so do not have well defined helicity.} Well if we plug these into the superPoincar\'{e} algebra relations, we can see that 
\be
    \begin{split}
        a_1, a_2^{\dagger} : m_s & \to m_s+\frac{1}{2} \\
        a_1^{\dagger}, a_2 : m_s & \to m_s-\frac{1}{2},
    \end{split}
\ee  
and so the former are our raising operators and the latter the lowering operators. 

\br 
    Note that creation/annihilation and raising/lowering do \textit{not} agree here. That is our creation operators are the daggered $a_1^{\dagger}, a_2^{\dagger}$, but the raising operators are $a_1, a_2^{\dagger}$. Therefore when we build the states up below we will both increase and decrease the spin value. The way we can remember which does which is that in the massless case raising = creation and there we only had the $2$ index. That is, for the massive case, $a_2^{\dagger}$ is both a creation operator and a raising operator, but $a_1^{\dagger}$ is a creation operator \textit{but} it is a \textit{lowering} operator. 
\er

Great, now we start again from our Clifford vacuum which label the by the mass (which is omitted here) and spin $s$. Now it's important to note that we actually  have \textit{vacua}, plural. This comes from the spin degeneracy
\bse 
    m_s \in \{ -s, -s+1 , ... ,s-1,s\}
\ese 
So we have $2s+1$ different vacua. These vacua are all annihilated by the \textit{annihilation} operators $a_1$ and $a_2$, and we build our states by acting with the \textit{creation} operators $a_1^{\dagger}$ and $a_2^{\dagger}$. Again we emphasise that creation does \textit{not} mean raising here. 

As we now have multiple raising operators and multiple vacua, we tend to indicate this building of the rep diagrammatically, as we now demonstrate.\footnote{We only give the multiplets in the absence of gravity, i.e. $m_s\leq 1$. Of course there are also massive supergravity multiplets too. } 

\ben[label=(\roman*)]
    \item $\cN=1$ massive $\chi$-plet, $s=0$: 
    \begin{center}
        \btik 
            \draw[thick, ->] (-2,0) -- (2,0);
            \node at (1.8,-0.2) {$m_s$};
            \draw[thick] (0,-0.1) -- (0,0.1) node [above] {$0$};
            \node at (0,-0.5) {$0$};
            \draw[->] (-0.2,-0.75) -- (-1.25,-1.5);
            \node[left] at (-0.5,-0.75) {$a_1^{\dagger}$};
            \draw[->] (0.2,-0.75) -- (1.25,-1.5);
            \node[right] at (0.5,-0.75) {$a_2^{\dagger}$};
            \node at (-1.5,-1.8) {$-\frac{1}{2}$};
            \node at (1.5,-1.8) {$\frac{1}{2}$};
            \draw[->] (-1.25,-2.1) -- (-0.2,-2.85);
            \node[left] at (-0.7,-2.7) {$a_2^{\dagger}$};
            \draw[->] (1.25,-2.1) -- (0.2,-2.85);
            \node[right] at (0.7,-2.7) {$a_1^{\dagger}$};
            \node at (0,-3) {$0$};
        \etik 
    \end{center}
    So our states are 
    \bse 
        (-1/2, \, 2\times(0) , \, 1/2)
    \ese 
    and we have a massive complex scalar, $2\times(0)$, and a Majoranna Fermion, $(-1/2,\,1/2)$. 
    \item $\cN=1$ V-plet, $s=1/2$: 
    \begin{center}
        \btik 
            \draw[thick, ->] (-4,0) -- (4,0);
            \node at (3.8,-0.2) {$m_s$};
            \draw[thick] (0,-0.1) -- (0,0.1) node [above] {$0$};
            \begin{scope}[xshift=-1.5cm]
                \node at (-0.1,-0.4) {$-\frac{1}{2}$};
                \draw[->] (-0.2,-0.75) -- (-1.25,-1.5);
                \node[left] at (-0.5,-0.75) {$a_1^{\dagger}$};
                \draw[->] (0.2,-0.75) -- (1.25,-1.5);
                \node[right] at (0.5,-0.75) {$a_2^{\dagger}$};
                \node at (-1.5,-1.8) {$-1$};
                \node at (1.5,-1.8) {$0$};
                \draw[->] (-1.25,-2.1) -- (-0.2,-2.85);
                \node[left] at (-0.7,-2.7) {$a_2^{\dagger}$};
                \draw[->] (1.25,-2.1) -- (0.2,-2.85);
                \node[right] at (0.7,-2.7) {$a_1^{\dagger}$};
                \node at (-0.1,-3.1) {$-\frac{1}{2}$};
            \end{scope}
            \begin{scope}[xshift=1.5cm]
                \node at (0,-0.5) {$\frac{1}{2}$};
                \draw[->] (-0.2,-0.75) -- (-1.25,-1.5);
                \node[left] at (-0.5,-0.75) {$a_1^{\dagger}$};
                \draw[->] (0.2,-0.75) -- (1.25,-1.5);
                \node[right] at (0.5,-0.75) {$a_2^{\dagger}$};
                \node at (1.5,-1.8) {$1$};
                \draw[->] (-1.25,-2.1) -- (-0.2,-2.85);
                \node[left] at (-0.7,-2.7) {$a_2^{\dagger}$};
                \draw[->] (1.25,-2.1) -- (0.2,-2.85);
                \node[right] at (0.7,-2.7) {$a_1^{\dagger}$};
                \node at (0,-3) {$\frac{1}{2}$};
            \end{scope}
        \etik 
    \end{center}
    So our states are 
    \bse 
        (-1, \, 2\times(-1/2), \, 2\times(0), \, 2\times(1/2), \, 1)
    \ese 
    which corresponds to a massive vector, $(-1,\,0,\,1)$, a massive Dirac Fermion, $(2\times(-1/2),\, 2\times(1/2))$, and a massive real scalar, $(0)$. 
\een 

As these two examples illustrate, we don't have to worry about CPT here as it is accounted for in our vacuum degeneracy and the fact that one of our creation operators is a raising operator and the other is a lowering operator. 

\br 
    We can compare the $\cN=1$ massive V-plet to the massless multiplets. How? Well recall that the Higgs gives mass to gauge Bosons and Fermions. Well what is happening here is the $\cN=1$ massless V-plet \textit{eats} a $\cN=1$ massless $\chi$-plet and gives us the $\cN=1$ massive V-plet. This can easily be checked by going back to the table above and checking that the degrees of freedom all add up (i.e. the numbers inside the brackets are the same). This is basically the content of the superHiggs mechanism. 
\er 


\subsubsection{Extended SUSY $\cN\geq2$}

Ok what if we have extended SUSY, i.e. $\cN\geq 2$? Well now we can't conclude that $Z^{IJ}=0=\bar{Z}^{IJ}$ and so things are more complicated. The first thing we note is what we said above: if we have odd $\cN$ we treat the last one as the $\cN=1$ case, which is why we put $0$s everywhere in the matrix above. So we only need to worry about the even $\cN$ case. 

As we did before, we skew diagonalise $Z^{IJ}$ to get $z_1, ... , z_{\cN/2}$. We then define 
\mybox{
    \be 
    \label{eqn:AnnihilationMassiveN>1}
        \begin{split}
            a^r_{\a} & := \frac{1}{\sqrt{2}}( Q_{\a}^{2r-1} + \epsilon_{\a\beta} (Q^{2r}_{\beta})^{\dagger}) \\
            b^r_{\a} & := \frac{1}{\sqrt{2}}( Q_{\a}^{2r-1} - \epsilon_{\a\beta} (Q^{2r}_{\beta})^{\dagger}) 
        \end{split}
    \ee
}
\noindent with $r=1,...,\cN/2$ labelling the $2\times 2$ blocks. We get the creation operators by taking Hermitian conjugate of these. This seems like a very unintuitive definiton, but the reason we do it is because they disentangle the anticommutation relations, and we have
\be 
\label{eqn:abAnticommutators}
    \{ a_{\a}^r , (a^s_{\beta})^{\dagger}\} = (2m + z_r) \del^{rs}\del_{\a\beta} \qand \{ b_{\a}^r , (b^s_{\beta})^{\dagger}\} = (2m - z_r) \del^{rs}\del_{\a\beta}
\ee
and all others vanishing. 

\bbox 
    Check that \Cref{eqn:abAnticommutators} hold.
\ebox 

Let's note how much more complicated this is to the above case. We have 
\bse 
    \underbrace{2}_{a/b} \times \underbrace{\frac{\cN}{2}}_{r} \times \underbrace{2}_{\a} = 2\cN 
\ese 
\textit{different} creation operators. We obviously have to keep track of all of them and see how they act on the states. We can check that  $(a_1^r)^{\dagger}, (b_1^r)^{\dagger}$ lower $m_s$ while  $(a_2^r)^{\dagger}, (b_2^r)^{\dagger}$ raise it. Again we remember this in the same way as for the $\cN=1$ case. 

This all seems very complicated, but we can make things a bit simpler in the following way. Unitarity (i.e. non-negative definitness of the Hilbert space) tells us that we require the anticommutation relations to be non-negative. The first expression in \Cref{eqn:abAnticommutators} tells us $z_r \geq -2m$ and the second expression tells us $z_r \leq 2m$. In total this give us
\mybox{
    \be
    \label{eqn:BPSBound}
        2m \geq |z_r| \qquad \forall r = 1,... , \cN/2.
    \ee 
}
\noindent which is known as the \textit{BPS bound}.\footnote{Bogomolny, Prasad, Sommerfield found a related bound for solitons without SUSY. Then  Witten and Olive came along and showed it works for SUSY as above.} We call a state that saturates the equality above a \textit{BPS particle/state}. 

Why does this help us? Well note that if we saturate the BPS bound for some $z_r$ then either the associated supercharges $a^r_{\a}, (a^r_{\a})^{\dagger}$ or $b^r_{\a}, (b^r_{\a})^{\dagger}$ annihilate the multiplet, which is \textit{shorter} as a result. The emphasise on shorter is a technical term, which we now expand on.

Let $0 \leq k \leq \cN/2$  be the number of central charges $z_r$ that saturate the bound, then 
\ben[label=(\roman*)] 
    \item $k=0$: (BPS not saturated) this is known as a \textit{long} multiplet. Here we have $2^{2\cN}$ d.o.f. from oscillators.
    \item $0<k<\cN/2$: we have a \textit{short} multiplet. Here we have $2^{2(\cN-k)}$ d.o.f. from oscillators. We say the multiplet is "$\frac{k}{\cN}$-BPS".\footnote{This is the fraction of supercharges \textit{preserved} by the multiplet.}
    \item $k=\cN/2$: then we get \textit{ultra-short} multiplet (shortest possible massive multiplet). Here we have $2^{\cN}$ d.o.f. from oscillators. In accordance with the above, this is $1/2$-BPS.
\een 

Let's now list some properties of BPS-saturated particles/states:
\ben 
    \item The defining relation, $2m = |z|$, is \textit{protected against} (i.e. not effected by) continuous deformations (changes of coupling constant, or of $\hbar$ or of some vev), because the number of d.o.f. cannot jump continuously. Of course if we have continuous parameters, we could change some expectation value continuously, and so $m$ and $z$ can continuously, but the relation $2m=|z|$ will always hold.\footnote{There is a caveat that we could have two short multiplets which recombine to make a longer multiplet.}
    \item A BPS state can only decay into BPS products with aligned central charges.\footnote{This comes from the idea that the $z$ are really vectors in some vector space and so by aligned we mean they are not only parallel but aligned. That is they are not antiparallel.} Indeed, consider the decay into $2$ decay products: 
    \ben
        \item Central charge conservation: $z = z_{(1)} + z_{(2)}$ where the brackets tell us which decay product we are talking about, it is \textit{not} a value of $r$. It follows from this that 
        \bse 
            |z| = |z_{(1)} + z_{(2)}| \leq |z_{(1)}| + |z_{(2)}| 
        \ese 
        \begin{center}
            \btik 
                \draw[thick, ->] (0,0) -- (0,1.5);
                \node at (-0.2,1.3) {$z$};
                \draw[thick, ->] (0,0) -- (-1,0.3);
                \node at (-1.2,0.5) {$z_{(1)}$};
                \draw[thick, ->] (0,0) -- (0.5,1);
                \node at (1,0.8) {$z_{(2)}$};
            \etik 
        \end{center}
        So if our initial state is BPS we have $|z|=2m$. Then if our decay products are also BPS we also have $|z_{(i)}|=2m_i$, so in total we have 
        \bse 
            2m \leq 2(m_1+m_2).
        \ese 
        \item Kinematics: Go to rest frame of decaying particle, then the decay can only happen if 
        \bse 
            m \geq m_1+m_2.
        \ese 
    \een 
    Combining both of these inequalities we are forced to conclude 
    \bse 
        m = m_1+m_2
    \ese 
    which tells us that the central charges must align as in the following diagram
    \begin{center}
        \btik 
            \draw[thick, ->] (0,0) -- (0,1.5);
            \node at (-0.2,1.3) {$z$};
            \draw[thick, ->] (-0.1,0) -- (-0.1,0.5);
            \node[left] at (-0.1,0.2) {$z_{(1)}$};
            \draw[thick, ->] (0.1,0) -- (0.1,1);
            \node[right] at (0.1,0.7) {$z_{(2)}$};
        \etik 
    \end{center}
\een

\subsubsection{$\cN=2$ Long Massive Multiplets}

Restricting to $s\leq 1$ (i.e. no gravity) the only $\cN=2$ long massive multipelet is the vector multiplet $s=0$:\footnote{We only label the $a_i^{\dagger}/b_i^{\dagger}$s on the first line. The rest are hopefully easily understood.} 
\begin{center}
    \btik 
        \draw[thick, ->] (-4,0) -- (4,0);
        \node at (1.8,-0.2) {$m_s$};
        \draw[thick] (0,-0.1) -- (0,0.1) node [above] {$0$};
        \node at (0,-0.5) {$0$};
        \draw[->] (-0.2,-0.75) -- (-1.25,-1.5);
        \node[left] at (-0.5,-0.75) {$a_1^{\dagger},b_1^{\dagger}$};
        \draw[->] (0.2,-0.75) -- (1.25,-1.5);
        \node[right] at (0.5,-0.75) {$a_2^{\dagger},b_2^{\dagger}$};
        \node at (-1.4,-1.8) {$2\times\big(-\frac{1}{2}\big)$};
        \node at (1.5,-1.8) {$2\times\big(\frac{1}{2}\big)$};
        \draw[->] (-1.25,-2.1) -- (-0.2,-2.85);
        \draw[->] (1.25,-2.1) -- (0.2,-2.85);
        \node at (0,-3.1) {$4\times(0)$};
        %
        \draw[->] (-1.5,-2.1) -- (-2.7,-2.85);
        \draw[->] (1.5,-2.1) -- (2.7,-2.85);
        \node at (2.9,-3.1) {$1$};
        \node at (-3,-3.1) {$-1$};
        %
        \draw[->] (-2.7,-3.45) -- (-1.65,-4.2);
        \draw[->] (2.7,-3.45) -- (1.65,-4.2);
        \draw[->] (-0.2,-3.45) -- (-1.25,-4.2);
        \draw[->] (0.2,-3.45) -- (1.25,-4.2);
        \node at (-1.4,-4.6) {$2\times\big(-\frac{1}{2}\big)$};
        \node at (1.5,-4.6) {$2\times\big(\frac{1}{2}\big)$};
        \draw[->] (-1.25,-4.8) -- (-0.2,-5.55);
        \draw[->] (1.25,-4.8) -- (0.2,-5.55);
        \node at (0,-5.7) {$0$};
    \etik 
\end{center}
which we can list as 
\bse 
    ( -1, 4\times(-1/2), 6\times(0), 4\times (1/2), 1),
\ese 
which has field content: one massive vector $(-1,0,1)$, two massive Dirac Fermions $(4\times(-1/2),4\times(1/2))$, and five real scalars $5\times(0)$.

\subsubsection{$\cN=2$ Short Multiplets}

Here we only have one type of creation operator, as the other saturate the BPS bound.\footnote{Note that for $\cN=2$ short and ultra-short coincide.} We don't label the creation operators as could be either $\{a_1^{\dagger},a_2^{\dagger}\}$ \textit{or} $\{b_1^{\dagger},b_2^{\dagger}\}$.

Again restricting to $s\leq 1$, we have three possible $\cN=2$ short massive multiplets 
\ben[label=(\roman*)]
    \item $\cN=2$ Massive Half-Hyperplet, $s=0$:
    \begin{center}
        \btik 
            \draw[thick, ->] (-2,0) -- (2,0);
            \node at (1.8,-0.2) {$m_s$};
            \draw[thick] (0,-0.1) -- (0,0.1) node [above] {$0$};
            \node at (0,-0.5) {$0$};
            \draw[->] (-0.2,-0.75) -- (-1.25,-1.5);
            \draw[->] (0.2,-0.75) -- (1.25,-1.5);
            \node at (-1.5,-1.8) {$-\frac{1}{2}$};
            \node at (1.5,-1.8) {$\frac{1}{2}$};
            \draw[->] (-1.25,-2.1) -- (-0.2,-2.85);
            \draw[->] (1.25,-2.1) -- (0.2,-2.85);
            \node at (0,-3) {$0$};
        \etik 
    \end{center}
    which has 
    \bse 
        (-1,2\times(0),1)
    \ese 
    which is a massive complex scalar $(2\times(0))$ and a massive (symplectic) Majorana Fermion $(-1/2,1/2)$. This is only possible for pseudoreal reps of the gauge group. 
    \item $\cN=2$ Massive Hyperplet $s=0$: This is the same as (i) but we add the CPT-conjugate states, so in total we have
    \bse 
        (2\times(-1),4\times(0),2\times(1))
    \ese 
    which is two massive complex scalars and one massive Dirac Fermion. 
    
    Sending $z\to 0$, (i) and (ii) become the massless half-hyperplet and hyperplet, respectively. 
    \item $\cN=2$ Short Massive Vector Multiplet $s=1/2$: 
    \begin{center}
        \btik 
            \draw[thick, ->] (-4,0) -- (4,0);
            \node at (3.8,-0.2) {$m_s$};
            \draw[thick] (0,-0.1) -- (0,0.1) node [above] {$0$};
            \begin{scope}[xshift=-1.5cm]
                \node at (-0.1,-0.4) {$-\frac{1}{2}$};
                \draw[->] (-0.2,-0.75) -- (-1.25,-1.5);
                \draw[->] (0.2,-0.75) -- (1.25,-1.5);
                \node at (-1.5,-1.8) {$-1$};
                \node at (1.5,-1.8) {$0$};
                \draw[->] (-1.25,-2.1) -- (-0.2,-2.85);
                \draw[->] (1.25,-2.1) -- (0.2,-2.85);
                \node at (-0.1,-3.1) {$-\frac{1}{2}$};
            \end{scope}
            \begin{scope}[xshift=1.5cm]
                \node at (0,-0.5) {$\frac{1}{2}$};
                \draw[->] (-0.2,-0.75) -- (-1.25,-1.5);
                \draw[->] (0.2,-0.75) -- (1.25,-1.5);
                \node at (1.5,-1.8) {$1$};
                \draw[->] (-1.25,-2.1) -- (-0.2,-2.85);
                \draw[->] (1.25,-2.1) -- (0.2,-2.85);
                \node at (0,-3) {$\frac{1}{2}$};
            \end{scope}
        \etik 
    \end{center}
    which is 
    \bse 
        (-1, 2\times(-1/2), 2\times(0), 2\times(1/2), 1)
    \ese
    with field content: one massive vector $(-1,0,1)$, one massive Dirac Fermion $(2\times(-1/2),2\times(1/2))$ and one massive real scalar $(0)$. This has the same degrees of freedom as the $\cN=1$ massive vector multiplet. 
    
    This corresponds to a superHiggs mechanism where the $\cN=2$ massless vector eats some of its own degrees of freedom to become a $\cN=2$ short massive vector.  In the $\cN=1$ language, we say the $\cN=1$ V-plet eats the $\cN=1$ $\chi$-plet in the adjoint rep that is part of the same $\cN=2$ V-plet. 
    
    Sending $z\to 0$, this reduces to an $\cN=2$ massless vector multiplet. 
\een 

\subsubsection{$\cN=4$ Ultrashort Multiplets}

Again restricting to $s\leq 1$, we have a single ultrashort multiplet. Here we have $4/2=2$ sets of creation operators, which again we don't label but obviously account for in terms of degrees of freedom. The diagram is simply 
\begin{center}
    \btik 
        \draw[thick, ->] (-4,0) -- (4,0);
        \node at (1.8,-0.2) {$m_s$};
        \draw[thick] (0,-0.1) -- (0,0.1) node [above] {$0$};
        \node at (0,-0.5) {$0$};
        \draw[->] (-0.2,-0.75) -- (-1.25,-1.5);
        \draw[->] (0.2,-0.75) -- (1.25,-1.5);
        \node at (-1.4,-1.8) {$2\times\big(-\frac{1}{2}\big)$};
        \node at (1.5,-1.8) {$2\times\big(\frac{1}{2}\big)$};
        \draw[->] (-1.25,-2.1) -- (-0.2,-2.85);
        \draw[->] (1.25,-2.1) -- (0.2,-2.85);
        \node at (0,-3.1) {$4\times(0)$};
        %
        \draw[->] (-1.5,-2.1) -- (-2.7,-2.85);
        \draw[->] (1.5,-2.1) -- (2.7,-2.85);
        \node at (2.9,-3.1) {$1$};
        \node at (-3,-3.1) {$-1$};
        %
        \draw[->] (-2.7,-3.45) -- (-1.65,-4.2);
        \draw[->] (2.7,-3.45) -- (1.65,-4.2);
        \draw[->] (-0.2,-3.45) -- (-1.25,-4.2);
        \draw[->] (0.2,-3.45) -- (1.25,-4.2);
        \node at (-1.4,-4.6) {$2\times\big(-\frac{1}{2}\big)$};
        \node at (1.5,-4.6) {$2\times\big(\frac{1}{2}\big)$};
        \draw[->] (-1.25,-4.8) -- (-0.2,-5.55);
        \draw[->] (1.25,-4.8) -- (0.2,-5.55);
        \node at (0,-5.7) {$0$};
    \etik 
\end{center}
which is the same as the $\cN=2$ long massive V-plet:
\bse 
    ( -1, 4\times(-1/2), 6\times(0), 4\times (1/2), 1),
\ese 
corresponding to: one massive vector $(-1,0,1)$, two massive Dirac Fermions $(4\times(-1/2),4\times(1/2))$, and five real scalars $5\times(0)$.

Again this is a superHiggs mechanism where the $\cN=4$ massless vector eats some of its own degrees of freedom to become a $\cN=2$ short massive V-plet. In the $\cN=2$ language: $\cN=2$ V-plet eats $\cN=2$ adjoint hyperplet in the same $\cN=4$ V-plet. 

Sending $z_r\to 0$\footnote{Note for $\cN=2$ we just wrote $z\to0$, as there there is only one $z_r$. Here we have two.} reduces to an $\cN=4$ massless V-plet.  
\chapter{Superspace, Superfileds \& Supersymmetric Actions}

So we have discussed the irreps of the SUSY algebra. The irreps correspond to our physical particles (i.e. they tell us the mass etc) and so they are, by definition, \textit{on-shell}. Of course if we want to do QFT we also want to know the \textit{off-shell} stuff. That is we want to be able to talk about off-shell propagators etc. We therefore need to introduce the notion of a superfield and superspace so that we can build up our SUSY actions. That is what we now do.

There are two formalisms for constructing our linear off-shell stuff 
\ben 
    \item Component Formalism: Here you realise SUSY algebra as transformations of Bosonic and Fermionic fields, known as \textit{components}. Again the SUSY will map Bosonic to Fermionic and vice versa. This is always possible but the problem is that the SUSY is not manifest. 
    \item Superfield Formalism: We introduce the notion of superspace on which the super charges $Q/\Bar{Q}$ act as differential operators. Superfields will be fields on the superspace, and we package all the components into a single object, our superfield. Here we will be able to make SUSY manifest, which makes writing down SUSY actions much easier.
\een 

Before moving on, let's make a few comments.
\ben[label=(\roman*)] 
    \item As it allows for the construction of SUSY actions easier, we will adopt the superfield formalism. We will focus on $\cN=1$ 4-dimensional as the superfield formalism is always available here. However we should point out that this formalism is \textit{not} always available as we increase $\cN$. 
    \item The component formalism can also be realised \textit{on-shell}, using exactly the same d.o.f. that appeared in the construction of supermultiplets. 
    \item It's a known fact that in a QFT going \textit{off-shell} results in Fermions gaining extra degrees of freedom. This is just because $\slashed{p}\psi=0$, for example, is a vector valued expression, and so it allows us to relate the different degrees of freedom, thereby reducing them. If we want to maintain SUSY, these extra d.o.f. must be matched by an equal number of Bosonic d.o.f, which become trivial when we go back on-shell.\footnote{We can also obtain this result by considering how the fields transform under the action of the supercharges. If you work through this you will find that a Boson transforms as $\del_{\epsilon}\phi = \epsilon\psi$ and a Fermion transforms as $\del_{\epsilon}\psi_{\a} = -i(\sig^{\mu}\bar{\epsilon})_{\a} \p_{\mu}\phi$. From here you can show that the commutator $[\del_{\epsilon_1},\del_{\epsilon_2}]$ acts on $\phi$ to give something of the form $\p_{\mu}\phi$, which is the action of $P_{\mu}$ on $\phi$. So the algebra closes \textit{off-shell} for the Bosons. However if you do the same calculation for the Fermion $\psi$ you will get exactly the same $\p_{\mu}\psi$ term, but then you get two additional terms which are (proportional to) the equation of motion for $\psi$. So off-shell the algebra doesn't close (the EoM terms stop it from), but \textit{on-shell} these terms vanish and so the algebra closes. The idea is that we can introduce another Bosonic field $F$ such that $\del_{\epsilon}\psi_{\a} = -i(\sig^{\mu}\bar{\epsilon})_{\a} \p_{\mu}\phi + \epsilon_{\a} F$ and then this $F$ term varys in \textit{exactly} the correct way to remove the EoM terms above. In doing this we get that the algebra even closes \textit{off-shell}. This is a very rough explanation of this argument and a more detailed one can be found in Section 1.3.1 of "Perspective On Supersymmetry" by Kane.} These new Bosonic fields, denoted $F$, are called \textit{auxiliary fields}, and they appear quadratically but without derivatives in the action. This means when we go on-shell we simple get $F=0$. Another way to see that they are irrelevant for on-shell physics is to consider the path integral approach, where these auxiliary fields can just be integrated out as they appear in Gaussian form (i.e. quadratically without derivatives).
\een 

\section{Superspace \& Superfields}

As we said above, we will focus on $\cN=1$ $D=4$ SUSY, and so we start by introducing coordinates in superspace $(x^{\mu}, \theta_{\a}, \Bar{\theta}_{\dot{\a}})$\footnote{Some people use notation like $\R^{1,3|4}$ to denote the spacetime + Grassman coordinates for the full superspace.} where $\theta/\bar{\theta}$ are Grassman-odd. We then use the fact that $P_{\mu}$ generates translations in $x^{\mu}$, to motivate us defining stuff such that $Q_{\a}$ does same for $\theta$ and $\bar{Q}_{\dot{\a}}$ for $\bar{\theta}_{\dot{\a}}$, all in such a way that $\{Q_{\a}, \bar{Q}_{\dot{\a}}\} = 2\sig^{\mu}_{\a\dot{\a}} P_{\mu}$. 

\subsection{Superspace Translations}

We start by introducing  $\epsilon_1/\bar{\epsilon}_2$, which are Grassman-odd, spinor, SUSY parameters, so that
\bse 
    [\epsilon_1 Q , \bar{\epsilon}_2 \bar{Q}] = 2 (\epsilon_1 \sig^{\mu}\bar{\epsilon}_2)P_{\mu},
\ese
where the contraction of $\a,\dot{\a}$ indices are implied. Next we note that if we are going to have the anticommutator of $Q$ and $\bar{Q}$ to give us spacetime translations, they can't just translate $\theta$ and $\bar{\theta}$. That is, we need $Q/\bar{Q}$ to also generate translations of $x^{\mu}$.

Ok so now that we have an idea of what we want our $Q/\bar{Q}$ to do, we need to work out how to obtain a form for them. In order to do this, let's think about standard (even) translations. The coordinates translate as
\bse 
    x^{\mu} \mapsto x^{\mu} + a^{\mu},
\ese 
and fields as
\bse 
    \phi(x) \mapsto \phi(x+a) = e^{ia^{\mu}\cP_{\mu}} \phi(x) e^{-ia^{\mu}\cP_{\mu}}, 
\ese 
where $\cP_{\mu}$ is an abstract generator. If we consider an infinitesimal translation, we can Taylor expand to get\footnote{Bonus exercise: check this.} 
\bse 
    \phi(x+a) = \phi(x) + ia^{\mu} [\cP_{\mu}, \phi(x)] + \cO(a^2).
\ese 
This holds for \textit{any} transformation, but we know how translations act so we can Taylor expand left hand side as 
\bse 
    \phi(x+a) = \phi(x) + a^{\mu}\p_{\mu}\phi(x) + \cO(a^2),
\ese
which allows us to conclude 
\bse 
    \del_a\phi(x) = ia^{\mu}[\cP_{\mu},\phi(x)]  = a^{\mu}\p_{\mu}\phi(x) \equiv i a^{\mu} P_{\mu}\phi(x)
\ese 
where $P_{\mu}$ is a differential operator given by $P_{\mu} = -i\p_{\mu}$. 

We now want to try mimic this for our super (odd) translations. 
\bse 
    \begin{split}
        \theta_{\a} & \mapsto \theta_{\a} + \epsilon_{\a} \\
        \bar{\theta}_{\dot{\a}} & \mapsto \bar{\theta}_{\dot{\a}} + \bar{\epsilon}_{\dot{\a}} \\
        x^{\mu} & \mapsto x^{\mu} + i \theta \sig^{\mu} \bar{\epsilon} - i \epsilon \sig^{\mu} \bar{\theta},
    \end{split}
\ese 
where the prefactors in the last result comes from the fact that we need to contract $\a/\dot{\a}$ indices and have a Bosonic result. We also need the parameter to be real and so the prefactors of the other terms have to be related. In other words, we should should have a $c$ and $\bar{c}$ on the two terms, but we have already set them to be 1, as it turns out this will give us the correct commutator relation.\footnote{Note that it follows from this relation between the two terms that we always have to consider the action of \textit{both} $Q$ and $\bar{Q}$ together as otherwise our spacetime translation won't be real.} 

\br 
    Note that we say that $\epsilon/\bar{\epsilon}$ transformations are infinitesimal because they are Grassman odd, and so the Taylor expansion will be truncated to the first order term. 
\er 

Now just as we had a field $\phi$ which depended on our spacetime coordiantes $\phi=\phi(x)$, we now want a \textit{superfield}, which we denote $Y$, that depends on our full superspace, i.e. $Y=Y(x,\theta,\bar{\theta})$. We then plug in our coordinate transformations above to obtain
\be
\label{eqn:SuperfieldTransformation}
    Y(x,\theta, \bar{\theta}) \mapsto Y(x^{\mu} + i \theta \sig^{\mu} \bar{\epsilon} - i \epsilon \sig^{\mu} \bar{\theta}, \theta+\epsilon , \bar{\theta} + \bar{\epsilon}).
\ee 
This is essentially our \textit{definition} of what a superfield is. That is a superfield is defined such that it transforms in this way. We then do what we did for the even case above and say that a general transformation is given by 
\bse 
    Y(x,\theta, \bar{\theta}) \mapsto  e^{i(\epsilon \cQ + \bar{\epsilon}\bar{\cQ})} Y(x,\theta,\bar{\theta}) e^{-i(\epsilon\cQ + \bar{\epsilon}\bar{\cQ})},
\ese 
with $\cQ$ being the abstract generators. Next, again we can Taylor expand to obtain 
\bse 
    Y(x,\theta,\bar{\theta}) + i[\epsilon\cQ + \bar{\epsilon}\bar{\cQ} , Y(x,\theta,\bar{\theta})] + \cO(\epsilon^2).
\ese 
Again this is true for any transformation rule, but we now use our definition of the transformation of the superfield, \Cref{eqn:SuperfieldTransformation}, which can then Taylor expand to be\footnote{Using $\p_{\a} := \frac{\p}{\p \theta^{\a}}$ and $\bar{\p}_{\dot{\a}}$ similar.}
\bse 
    Y(x,\theta,\bar{\theta}) + \big[i(\theta\sig^{\mu}\bar{\epsilon} - \epsilon\sig^{\mu}\bar{\theta})\p_{\mu} + \epsilon^{\a}\p_{\a} + \bar{\epsilon}^{\dot{\a}}\bar{\p}_{\dot{\a}}\big] Y(x,\theta,\bar{\theta}).
\ese
Note that the indices are \textit{not} contracted in our convention here. That is we defined our inner product for dotted indices to be `bottom left to top right' but here we have $\epsilon^{\dot{\a}}\bar{\p}_{\dot{\a}}$. This is just because of the way the Taylor expansion works, and will result in us getting the correct sign in the end. Putting this together, we have 
\bse 
    \begin{split}
        \del_{\epsilon,\bar{\epsilon}} Y(x,\theta,\bar{\theta}) & = i [\epsilon\cQ + \bar{\epsilon}\bar{\cQ} , Y(x,\theta,\bar{\theta})] \\
        & = \big[i(\theta\sig^{\mu}\bar{\epsilon} - \epsilon\sig^{\mu}\bar{\theta})\p_{\mu} + \epsilon^{\a}\p_{\a} + \bar{\epsilon}^{\dot{\a}}\bar{\p}_{\dot{\a}}\big] Y(x,\theta,\bar{\theta}) \\
        & \equiv i\big( \epsilon^{\a} Q_{\a} + \bar{\epsilon}_{\dot{\a}} \bar{Q}^{\dot{\a}}\big)Y(x,\theta,\bar{\theta})
    \end{split}
\ese 
where on the last line the indices are contracted in the "usual" way. We then finally conclude 
\mybox{
    \be 
    \label{eqn:QbarQDiffOperators}
        Q_{\a} = -i(\p_{\a} - i(\sig^{\mu}\bar{\theta})_{\a} \p_{\mu}) \qand \bar{Q}_{\dot{\a}} = i(\bar{\p}_{\dot{\a}} - i(\theta\sig^{\mu})_{\dot{\a}} \p_{\mu}).
    \ee 
}

\bp 
    The product of two superfields is a superfield.
\ep 

\bq 
    This just follows from the Leibniz rule for differentiation and the commutator result 
    \bse 
        [A,BC] = [A,B]C + B[A,C].
    \ese
    That is, we consider the transformation of the product $Y_1Y_2$ giving us 
    \bse 
        \begin{split}
            \del_{\epsilon,\bar{\epsilon}} (Y_1Y_2) & = [\epsilon\cQ + \bar{\epsilon}\bar{\cQ} , Y_1Y_2] \\ 
            & = [\epsilon\cQ + \bar{\epsilon}\bar{\cQ} , Y_1]Y_2 + Y_1[\epsilon\cQ + \bar{\epsilon}\bar{\cQ} , Y_2] \\
            & = (\del_{\epsilon,\bar{\epsilon}}Y_1)Y_2 + Y_1 (\del_{\epsilon,\bar{\epsilon}}Y_2) \\
            & = (i(\epsilon Q + \bar{\epsilon}\bar{Q})Y_1)Y_2 + Y_1 (i(\epsilon Q + \bar{\epsilon}\bar{Q})Y_2) \\
            & = i(\epsilon Q + \bar{\epsilon}\bar{Q}) (Y_1Y_2).
        \end{split}
    \ese 
\eq

Let's just clarify the notation we have used a bit better. Our derivatives are 
\bse 
    \p_{\mu} = \frac{\p}{\p x^{\mu}}, \qquad \p_{\a} = \frac{\p}{\p \theta^{\a}} \qand \bar{\p}_{\dot{\a}} = \frac{\p}{\p \bar{\theta}^{\dot{\a}}}.
\ese 
Now recall that 
\bse 
    [\p_{\mu}, x^{\nu}] = (\p_{\mu}x^{\nu})\b1 = \del_{\mu}^{\nu} 
\ese
where the terms in the commutator should be treated as \textit{operators}, in the sense that $x^{\nu}$ is the operator that says "multiply by the \textit{number} $x^{\nu}$", while on the right-hand side they are simply derivatives/numbers. In other words, really we should consider the action of the commutator on some field $f(x)$ and obtain 
\bse 
    \begin{split}
        [\p_{\mu}, x^{\nu}]f(x) & = \p_{\mu}\big(x^{\nu}f(x)\big) - x^{\nu}\p_{\mu}f(x) \\
        & = \big(\p_{\mu}x^{\nu}\big)f(x) + x^{\nu}\p_{\mu}f(x) - x^{\nu}\p_{\mu}f(x) \\
        & = \del^{\nu}_{\mu} f(x),
    \end{split}
\ese 
and so we simply "strip off" the $f(x)$ as it was arbitrary to obtain the expression above. In a completely analogous way, we also have
\bse 
    [\p_{\mu},\p_{\nu}] = 0.
\ese

\br 
    Note really we should have a $\b1$ on the right-hand side for our equal sign to make sense. That is the commutator is an operator that is defined by its action on a field, so the right-hand side should also be an operator that says "act with the identity, weighted by the number $\del^{\nu}_{\mu}$". Of course it is very standard notation to drop the $\b1$, and so we have done so here. However this remark is included as it might help clear up confusion with the argument made above. 
\er 

We now want to translate these relations into ones for our Grassman-odd derivatives $\p_{\a}$ and $\bar{\p}_{\dot{\a}}$. As with all even to odd relations, commutators are replaced with anticommutators, and so we obtain
\bse
    \begin{split}
        \{\p_{\a}, \theta^{\beta}\} & = (\p_{\a}\theta^{\beta}) = \del_{\a}^{\beta} \\
        \{\bar{\p}_{\dot{\a}}, \bar{\theta}^{\dot{\beta}}\} & = (\p_{\a}\theta^{\beta}) = \del_{\dot{\a}}^{\dot{\beta}} \\
        \{\p_{\a}, \bar{\theta}^{\dot{\beta}}\} & = 0 = \{\bar{\p}_{\dot{\a}}, \theta^{\beta}\}
    \end{split}
\ese
and 
\be 
\label{eqn:GrassmanDerivativesAnticommutators}
    \{\p_{\a}, \p_{\beta}\} = \{\bar{\p}_{\dot{\a}}, \bar{\p}_{\dot{\beta}}\} = \{\p_{\a}, \bar{\p}_{\dot{\beta}}\} = 0.
\ee
Finally note that it follows from our conventions that 
\bse 
    (\p_{\a})^{\dagger} = \bar{\p}_{\dot{\a}}. 
\ese 

\br 
    There's another easy way to understand \Cref{eqn:GrassmanDerivativesAnticommutators}. Firstly we note that dotted and undotted indices live in different spaces and don't talk to each other, so we expect the cross anticommutator to vanish. Next if we take $\p_{\a}\p_{\beta}$ with $\a\neq \beta$ we are differentiating w.r.t. two different $\theta^{\gamma}$s and so the result will vanish (the minus sign coming from having to swap the $\theta^{\a}\theta^{\beta} = -\theta^{\beta}\theta^{\a}$ in the function it acts on. Then we simply note that because $\theta^{\a}/\bar{\theta}^{\dot{\a}}$ are Grassman-odd we never have $\theta^{\a}\theta^{\a}$ or $\bar{\theta}^{\dot{\a}} \bar{\theta}^{\dot{\a}}$ as these vanish. Therefore the action of the same derivative twice will always vanish, as our function can have at most 1 power of the variable so the first derivative removes this and then the second has nothing to act on. This gives us the results above straight away. Note that essentially what we've shown is that we can treat the derivatives themselves as Grassman-odd expresions and so their anticommutators obviously vanish. 
\er 

\bbox 
    \ben 
        \item Show that $\{Q_{\a},\bar{Q}_{\dot{\a}}\} = 2\sig^{\mu}P_{\mu}$ and $\{Q_{\a},Q_{\beta}\}=0=\{\bar{Q}_{\dot{\a}},\bar{Q}_{\dot{\beta}} \}$. 
        \item Using $(\p_{\a})^{\dagger} = \bar{\p}_{\dot{\a}}$ and $(\p_{\mu})^{\dagger} = -\bar{\p}_{\mu}$, show that $(Q_{\a})^{\dagger}= \bar{Q}_{\dot{\a}}$ and that $(\epsilon Q + \bar{\epsilon}\bar{Q})$ is Hermitian if $\bar{\epsilon}_{\dot{\a}} = (\epsilon_{\a})^*$. 
        \item Show that $[\epsilon_1 Q, \bar{\epsilon}_2\bar{Q}] = \epsilon_1^{\a}\bar{\epsilon}_2^{\dot{\a}} \{ Q_{\a} ,\bar{Q}_{\dot{\a}}\} = 2(\epsilon_1 \sig^{\mu} \bar{\epsilon}_2) P_{\mu}$. 
        \item Show that 
        \be 
        \label{eqn:RaisingIndexOnAlphaDerivative}
            \p^{\a} := \frac{\p}{\p \theta_{\a}} = - \epsilon^{\a\beta}\p_{\beta} \qand \bar{\p}^{\dot{\a}} := \frac{\p}{\p \bar{\theta}_{\a}} = - \epsilon^{\bar{\a}\bar{\beta}}\bar{\p}_{\dot{\beta}}.
        \ee 
        \textit{Hint: Show that $\p^{\a}F(\theta) = - \epsilon^{\a\beta}\p_{\beta}F(\theta)$ for any function $F$ of $\theta_1$ and $\theta_2$.}
    \een 
\ebox 

\br 
    Note that \Cref{eqn:RaisingIndexOnAlphaDerivative} is a somewhat funny fact, given our experience of GR because of the minus sign. This again is all related to the fact that we are raising the indices using a 2-form instead of a metric. 
\er 

\subsection{Components Of A Superfield}

We now extend the argument we in the above remark about only getting linear terms in $\theta^{\a}$ etc. If we have a general superfield $Y(x,\theta,\bar{\theta})$ this can be \textit{at most} quadratic in $\theta$ or $\bar{\theta}$, separetly.\footnote{I.e. we really mean $\theta\theta$ and $\bar{\theta}\bar{\theta}$. Of course we can also have the product $\theta\theta\bar{\theta}\bar{\theta}$, as we will see shortly.} Its quadratic because $\a=1,2$ so we can have cross terms like
\bse 
    \theta_1\theta_2.
\ese
This means that in what follows we will have $\theta\theta$ and $\bar{\theta}\bar{\theta}$ terms and we might at first be tempted to say "that's zero", but should remember that it just means the cross terms like above. In other words, our expansions will vanish at cubic order because 
\bse 
    \theta_{\a}\theta_{\beta}\theta_{\g} = 0 = \bar{\theta}_{\dot{\a}}\bar{\theta}_{\dot{\beta}}\bar{\theta}_{\dot{\g}} 
\ese 
as two of the indices must be the same.

\bnn 
    In what follows we shall use our inner product notations, namely 
    \bse 
        \theta\theta := \theta^{\a}\theta_{\a} \qand \bar{\theta}\bar{\theta} = \bar{\theta}_{\dot{\a}}\bar{\theta}^{\dot{\a}}.
    \ese 
    Any time two $\theta/\bar{\theta}$s appear next to each other this inner product is assumed. We could be abit more careful and put brackets around everything like 
    \bse 
        (\theta\theta)(\bar{\theta}\bar{\theta})
    \ese    
    so that we know the $\theta$s are associated and separately the $\bar{\theta}$s. However we don't have an inner product defined between our $\theta^{\a}$ and $\bar{\theta}^{\dot{\a}}$ and so no hopefully confusion should arise. However if we have 4 objects of the same index structure next to each other, we shall try use brackets, e.g. in the relation
    \bse 
        (\theta\psi)(\theta\chi) = - \frac{1}{2}(\theta\theta)(\psi\chi).
    \ese 
    However should we forget to do this, again hopefully no confusion should arise because there is actually only one way to read $\theta\psi\theta\chi$, and it is the left-hand side above. 
\enn 

Ok with that cleared up, let's consider expanding a general superfield around $(x,0,0)$ and terminate the result at $\theta\theta \bar{\theta}\bar{\theta}$. The resulting $x$ dependent functions are known as \textit{components}. We have 
\mybox{
    \be 
    \label{eqn:SuperfieldInComponents}
        \begin{split}
            Y(x,\theta,\bar{\theta}) = & y(x) + \theta \psi(x) + \bar{\theta}\bar{\chi}(x)  + \theta\theta m(x) + \bar{\theta}\bar{\theta} \bar{n}(x) + \theta \sig^{\mu} \bar{\theta} v_{\mu}(x) \\
            & + \theta \theta \bar{\theta} \bar{\l}(x) + \bar{\theta}\bar{\theta}  \theta\rho(x) + \theta\theta  \bar{\theta}\bar{\theta}D(x).
        \end{split}
    \ee 
}
\noindent We call $y(x) = Y(x,0,0)$ the \textit{bottom component}, and $D(x)$ the the \textit{top component}. 

We can now ask the question of "how do the components transform under a SUSY translation?" The answer is to recall that a superfield transforms as
\bse 
    \del_{\epsilon\bar{\epsilon}} Y(x,\theta,\bar{\theta}) = i\big(\epsilon Q + \bar{\epsilon}\bar{Q}\big) Y(x,\theta,\bar{\theta}). 
\ese 
We can then use our definitions of $Q/\bar{Q}$ in terms of derivatives, \Cref{eqn:QbarQDiffOperators}, and act on our $Y(x,\theta,\bar{\theta})$ and extract the transformation of the components by comparing the result to 
\bse 
    \del_{\epsilon\bar{\epsilon}} Y(x,\theta,\bar{\theta}) = \del_{\epsilon,\bar{\epsilon}}y(x) + \theta \del_{\epsilon,\bar{\epsilon}} \psi(x) + \bar{\theta} \del_{\epsilon,\bar{\epsilon}}\bar{\chi} + ... + \theta\theta \bar{\theta}\bar{\theta} \del_{\epsilon,\bar{\epsilon}} D(x).
\ese
For a general superfield this is obviously a very long and tedious calculation and so we don't do it here. We will, however, do this calculation for a so-called \textit{Chiral superfield} soon. The important point to note about this calculation is that the $Q/\bar{Q}$ are differential operators in superspace, so they act on the \textit{total} expansion of $Y$. That is they act on the $\theta/\bar{\theta}$, not just the components.  

We now raise an important point: we saw earlier that if we have $\cN=1$ without gravity\footnote{Which we want if we are going to do QFT.} then all the massless irreps had $4$ d.o.f. (i.e. there were $4$ numbers in our (-1, ... , 1) stuff), and the most we could have for a massive irrep was $8$, which the $\cN=1$ massive V-plet. However if we leave our superfield, \Cref{eqn:SuperfieldInComponents}, completely general then there is \textit{no way} for us to form an irrep. That is if the components are all independent and complex we would have\footnote{We can work out what are Bosons and what are Fermions using the Grassman nature. That is $Y(x,\theta,\bar{\theta})$ is Grassman even (as $y(x) = Y(x,0,0)$ is) and so anything that appears with an odd number of $\theta/\bar{\theta}$s is a Fermion and anything else is a Boson.} 
\ben[label=(\roman*)]
    \item Bosons: 
        \begin{itemize}
            \item $y,m, \bar{n}$ and $D$: 2 $\R$ d.o.f. each.
            \item $v_{\mu}$: 8 $\R$ d.o.f.
        \end{itemize}
    \item Fermions: 
        \begin{itemize}
            \item $\psi, \bar{\chi}, \bar{\l},\rho$: 4 $\R$ d.o.f. each.
        \end{itemize} 
\een 
So in total we have $32$ real d.o.f. If we want to obtain an irrep from our superfield we are, therefore, going to have to impose some relations between the components.\footnote{Note that we can still get a representation for a superfield it will just be \textit{reducible}.} Of course these constraints will need to be consistent with SUSY otherwise everything we have done is gone. Essentially there are two constraints we can impose
\ben[label=(\roman*)]
    \item Impose reality condition, e.g. $Y=\bar{Y} =: Y^{\dagger}$.
    \item Differential constraint in superspace. 
\een 

The first thing we note is that the reality condition itself wont be good enough for us to get an irrep as it will simply reduce our $32$ d.o.f. to $16$, which is still greater than $8$.\footnote{Big boy maths right there...} We therefore focus now on (ii) and try impose some differential constraints and will return to the reality condition later. 

\subsection{Supercovariant Derivatives}

As we just said, our constraint must be SUSY consistent and so we can't just use any old derivative. In other words we need a derivative that is \textit{covariant} w.r.t to SUSY. We therefore introduce the \textit{supercovariant derivatives}:
\mybox{
    \be 
    \label{eqn:SupercovariantDerivatives}
        D_{\a}  = \p_{\a} + i(\sig^{\mu}\bar{\theta})_{\a} \p_{\mu} \qand \bar{D}_{\dot{\a}}  = \bar{\p}_{\dot{\a}} + i(\theta\sig^{\mu})_{\dot{\a}} \p_{\mu}.
    \ee 
}

These have essentially be pulled out of thin air, but we can justify their form by explaining how they were constructed. 
\ben[label=(\roman*)]
    \item As with basically everything so far, we get the barred version by taking the Hermitian conjugate 
    \bse 
        \big(D_{\a}\big)^{\dagger} = \bar{D}_{\dot{\a}}. 
    \ese 
    \item They anticommute to give
    \be 
    \label{eqn:DbarDAnticommutator}
        \{D_{\a}, \bar{D}_{\dot{\a}} \} = 2i \sig^{\mu}_{\a\dot{\a}} \p_{\mu},
    \ee 
    and all others vanishing. 
    \item They anticommute with the supercharges 
    \be 
    \label{eqn:DQCommutators}
        \{D_{\a}, Q_{\beta}\} = \{D_{\a}, \bar{Q}_{\dot{\beta}}\} = \{\bar{D}_{\dot{\a}} , \bar{Q}_{\dot{\beta}} \} = \{\bar{D}_{\dot{\a}}, Q_{\beta}\} = 0.
    \ee 
\een

\bbox 
    Prove that \Cref{eqn:DbarDAnticommutator,eqn:DQCommutators} hold.
\ebox 

Why do we want these properties? Well (i) is just because that's how everything we have constructed so far has worked. Condition (ii) is because again it gives us this "squaring superderivative gives spacetime derivative" relation which we have used several times so far. In other words it is the differential operator version of $\{Q_{\a},\bar{Q}_{\dot{\a}}\} = 2\sig^{\mu}_{\a\dot{\a}}P_{\mu}$.\footnote{Note there appears to be a sign off as $P_{\mu}=-i\p_{\mu}$. However there is a non-trivial relation between going between commutators of \textit{operators} and \textit{derivatives}. Essentially you pick up a minus sign, $[Q,\bar{Q}]= - [D,\bar{D}]$. The reason this is the case can be found in Section 3.3.1 of my CFT notes.} Condition (iii) is actually the most useful in terms of fixing our d.o.f. problem above. In this sense, condition (iii) is really our defining property for the supercovariant derivatives. 

The idea is we essentially defined our superfileds by their transformation property, which are generated by our supercharges $Q,\bar{Q}$. Now because our supervcovariant derivative anticommute with the $Q,\bar{Q}$ we can "move them though" the transformation and impose constraints on the components of your superfield \textit{without effecting the SUSY transformation property}. That is, we still have a superfield, which this quick calculation proves
\bse 
    D_{\a}(\del_{\epsilon,\bar{\epsilon}}Y) = D_{\a}i(\epsilon Q + \bar{\epsilon} \bar{Q}) Y = i(\epsilon Q + \bar{\epsilon} \bar{Q}) D_{\a}Y = \del_{\epsilon,\bar{\epsilon}}(D_{\a}Y)
\ese 
and similarly for the barred stuff. We can therefore use $D_{\a}$ to impose conditions on our $Y$ such as $D_{\a}Y =0$ which obviously constrains the components. 

\section{Chiral Superfield}

Armed with our supercovariant derivatives, we can now try constrain our superfields and obtain irreps. The first thing we do is to define \textit{Chiral superfields}

\bd[(Anti)Chiral Superfield]
    A \textit{Chiral superfield}, $\Phi$, is a superfield that obeys 
    \be 
    \label{eqn:ChiralSuperfield}
        \bar{D}_{\dot{\a}}\Phi = 0.
    \ee 
    Similarly we define an \textit{antichiral superfield} $\bar{\Phi} := \Phi^{\dagger}$ which satisfies 
    \be
    \label{eqn:AntichiralSuperfield}
        D_{\a}\bar{\Phi} = 0.
    \ee 
\ed 


\bp 
    Products of Chiral superfields are Chiral superfields.
\ep 

\bq 
    Again this just follows from Leibniz so we don't write it out again.
\eq 

\br 
    Note if we make the superfield both chiral \textit{and} antichiral, i.e. it is annihilated by both $D$ and $\bar{D}$, then, from the fact that the commutator of the two is proportional to $\p_{\mu}$, we have that the superfield is constant. That is it doesnt depend on $\theta,\bar{\theta}$ or $x$. This is boring so we wont consider this.
\er 

To solve the (anti)chiral constraint, introduce new (anti)chiral superspace coordinates
\mybox{
    \be 
    \label{eqn:yCoordinates}
        y^{\mu} := x^{\mu} + i \theta \sig^{\mu} \bar{\theta} \qand \bar{y}^{\mu} := x^{\mu} - i \theta \sig^{\mu} \bar{\theta}
    \ee 
}

Why are these coordinates useful? Well first we write our supercovariant derivative in $y$ coordinates as
\bse 
    \bar{D}_{\dot{\a}} = a \bar{\p}_{\dot{\a}} + b \big(\theta\bar{\sig}^{\mu}\big)_{\dot{\a}}\p_{\mu}^y,
\ese 
where $a$ and $b$ are constants we find considering the action on $\bar{\theta}^{\dot{\beta}}$ and $y^{\nu}$. Now note
\bse 
    \begin{split}
        \bar{D}_{\dot{\a}} y^{\mu} & = (\bar{\p}_{\dot{\a}} + i(\theta\sig^{\nu})_{\dot{\a}}\p_{\nu}) (x^{\mu} + i\theta\sig^{\mu}\bar{\theta}) \\
        & = -i(\theta\sig^{\mu})_{\dot{\a}} + i (\theta\sig^{\nu})_{\dot{\a}}\del^{\mu}_{\nu} \\
        & = 0,
    \end{split}
\ese 
which allows us to read off $b=0$. We can also show that $a=1$, and so in total
\bse 
    \bar{D}_{\dot{\a}} = \bar{\p}_{\dot{\a}}.
\ese 
This is why the $y$ coordinates are useful; our chiral superfield condition now simply becomes 
\bse 
    \bar{\p}_{\dot{\a}}\Phi(y,\theta,\bar{\theta}) = 0 \qquad \iff \qquad \Phi(y,\theta,\bar{\theta}) = \Phi(y,\theta). 
\ese 
Similarly if we work with $\bar{y}$ then we get 
\bse 
    D_{\a} = \p_{\a} 
\ese 
and so our antichiral condition becomes 
\bse 
    \bar{\Phi}(\bar{y},\theta,\bar{\theta}) = \bar{\Phi}(\bar{y},\bar{\theta}). 
\ese 

Ok great, so if we switch from $(x,\theta,\bar{\theta})$ to $(y,\theta,\bar{\theta})$, then the chiral constraint is solved by
\mybox{
    \be
    \label{eqn:ChiralSuperfieldComponentsIny}
        \Phi(y,\theta) = \phi(y) + \sqrt{2} \theta \psi(y) - \theta \theta F(y)
    \ee 
}
\noindent where the second line is just an expansion in $\theta$, the prefactors are just conventions that will be useful later. This is the most general solution to the constraint. 

\br 
    Note, by a similar logic to before when we were working out the degrees of freedom for a general superfield, we see that if $\Phi$ was a $\C$ scalar then it follows that $\phi(y)$ is a  $\C$ scalar, $\psi_{\a}$ is a Weyl spinor, and $F(y)$ is a $\C$ scalar. In fact $F(y)$ will turn out to be a so-called \textit{auxillary} field, and it corresponds exactly the additional off-shell Bosonic d.o.f. we needed to add as per (iii) at the start of this chapter. This will be more clear soon. If we count d.o.f. we then have $2$ $\R$ from both $\phi$ and $F$ and $4$ $\R$ from $\psi$, which gives us a total of $8$. This is great because it's exactly what we wanted!
\er 

So we have an expression for our chiral superfield in $(y,\theta)$, but really we want to go back to $(x,\theta,\bar{\theta})$. This is easily achieved by putting the definition of $y$ back in and Taylor expanding around $y$. This is a rather tedious calculation, but we can show that the result is 
\mybox{
\be 
    \Phi(y,\theta) = \phi(x) + i\theta\sig^{\mu}\bar{\theta} \p_{\mu}\phi(x) - \frac{1}{4} \theta \theta \bar{\theta}\bar{\theta} \square \phi(x) + \sqrt{2}\theta\psi(x) - \frac{i}{\sqrt{2}} \theta\theta \p_{\mu}\psi(x) \sig^{\mu} \bar{\theta} - \theta\theta F(x),
\ee 
}
\noindent In order to obtain this result, the following relations are used\footnote{The proof that these hold is part of the worksheet questions for the course.}
\bse 
    \begin{split}
        \theta^{\a}\theta^{\beta} & = -\frac{1}{2}\epsilon^{\a\beta} \theta\theta \\
        \bar{\theta}^{\dot{\a}}\bar{\theta}^{\dot{\beta}} & = -\frac{1}{2}\epsilon^{\dot{\a}\dot{\beta}} \bar{\theta}\bar{\theta} \\
        (\theta\psi)(\theta\chi) & = -\frac{1}{2} (\theta\theta)(\psi\chi) \\
        (\theta\sig^{\mu}\bar{\theta})(\theta\sig^{\nu}\bar{\theta}) & = \frac{1}{2}(\theta\theta)(\bar{\theta}\bar{\theta}) \eta^{\mu\nu}
    \end{split}
\ese 
\bse 
    (\theta\psi)(\theta\chi) = -\frac{1}{2} (\theta\theta)(\psi\chi)
\ese 

Next we want to find the SUSY variation of the chiral superfield, again it's useful to switch from $(x,\theta,\bar{\theta})$ to $(y,\theta,\bar{\theta})$. In the same way that we found $\bar{D}_{\dot{\a}}$ in $(y,\theta,\bar{\theta})$ above, we obtain 
\bse 
    Q_{\a} = -i\p_{\a}, \qand \bar{Q}_{\dot{\a}} = i\bar{\p}_{\dot{\a}} + 2(\theta\sig^{\mu})_{\dot{\a}} \frac{\p}{\p y^{\mu}}.
\ese
From here we have 
\bse 
    \begin{split}
        \del_{\epsilon,\bar{\epsilon}}\Phi(y,\theta) & = i(\epsilon Q + \bar{\epsilon}\bar{Q}) \Phi(y,\theta) \\
        & = \bigg(\epsilon^{\a}\p_{\a} + 2i\theta\sig^{\mu}\bar{\epsilon}\frac{\p}{\p y^{\mu}}\bigg) \Phi(y,\theta) \\
        & = \sqrt{2}\epsilon \psi(y) - 2 \epsilon \theta F(y) + 2i\theta\sig^{\mu} \bar{\epsilon} \big( \p_{\mu}\phi(y) + \sqrt{2}\theta\p_{\mu}\psi(y) \big) \\
        & = \sqrt{2}\epsilon\psi(y) + \sqrt{2}\theta\big[ -\sqrt{2}\epsilon F(y) + \sqrt{2}i \sig^{\mu} \bar{\epsilon}\p_{\mu} \phi(y)\big] - \theta\theta i \sqrt{2}\p_{\mu}\psi(y) \sig^{\mu}\bar{\epsilon},
    \end{split}
\ese
where for the last term we have used the identity above again. From here we can read off the variations of the components:\footnote{This result is exactly what we were talking about in footnote 1 at the start of this chapter. Of course we didn't run into any problems with the algebra closing here because we had already introduced the $F$ so that everything worked. }
\mybox{
    \be 
    \label{eqn:VariationphipsiF}
        \begin{split}
            \del_{\epsilon,\bar{\epsilon}} \phi & = \sqrt{2}\epsilon\psi \\
            \del_{\epsilon,\bar{\epsilon}} \psi & = \sqrt{2} i \sig^{\mu} \bar{\epsilon} \p_{\mu}\phi - \sqrt{2}\epsilon F \\
            \del_{\epsilon,\bar{\epsilon}} F & = i \sqrt{2}\p_{\mu}\psi(y) \sig^{\mu}\bar{\epsilon}
        \end{split}
    \ee
}

\section{Superspace Integrals \& Supersymmetric Actions}

If we are going to write down a SUSY action, we obviously need to be able to integrate over our superspace $(x,\theta,\bar{\theta})$. We already know how to do the $x$ integrals, obviously, so it's just the Grassman integrals we need to define.

\subsection{Grassman Integral (or Berezin integrals) In 1 Variable $\theta$}

The Grassman integral is essentially defined via the two following conditions 
\mybox{
    \be
    \label{eqn:GrassmanIntegralDefinition}
        \int d\theta \, 1 = 0 \qand \int d\theta \, \theta = 1.
    \ee 
}
\noindent The second of these two results seems rather strange given our understanding of normal $x$ integrals. Why do we have it? The answer is that basically we want to be able to manipulate the Grassman integrals in the same ways we manipulate our $x$ integrals, and the three main properties we want to maintain are 
\ben[label=(\roman*)]
    \item Translation invariance 
    \bse 
        \int d\theta \, (\theta + \epsilon) = \int d\theta \, \theta.
    \ese 
    \item The integral over $\theta\del(\theta)$, with $\del$ being the delta function, vanishes
    \bse 
        \int d\theta \, \theta \del(\theta) =0.
    \ese 
    If we compare this to the fact that we know $\theta\theta=0$\footnote{The $\theta$ here is a single thing, it's not $\theta_{\a}$ as above so the square vanishes.} we see that essentially inside an integral $\theta = \del(\theta)$. 
    \item The integral of a total derivative vanishes 
    \bse 
        \int d\theta \frac{\p}{\p \theta}X = 0.
    \ese
    If we put this together with $\p_{\theta}\p_{\theta} = 0$, as we saw before, we see that the integral behaves like a derivative, i.e. 
    \bse 
        \int d\theta \sim \p_{\theta}.
    \ese 
\een
Putting these conditions together will give you exactly \Cref{eqn:GrassmanIntegralDefinition}.

\subsection{$\cN=1$ Superspace Integrals}

Ok so that was just the general discussion of how to integrate w.r.t. Grassman numbers, we now want to go back to SUSY and in particular $\cN=1$ SUSY. In this case we have $4$ Grassman numbers $\{\theta^1,\theta^2,\bar{\theta}^{\dot{1}}, \bar{\theta}^{\dot{2}}\}$. We need to adapt \Cref{eqn:GrassmanIntegralDefinition} to higher dimensional integrals so that we can integrate over all these Grassman numbers. We therefore define
\be 
\label{eqn:d2Theta}
    d^2\theta := \frac{1}{2}d\theta^1 d\theta^2 \qand d^2\bar{\theta} := \frac{1}{2}d\bar{\theta}^{\dot{2}} d\bar{\theta}^{\dot{1}}
\ee
where we note that for the barred version the $\dot{2}$ index comes first.\footnote{This comes from the fact that the Hermitian conjugate of two Grassman numbers swaps their order. So if we want $(d^2\theta)^{\dagger} = d^2\bar{\theta}$ then using $(\theta^i)^{\dagger} = \bar{\theta}^{\dot{i}}$ we get exactly the result above. } 
We also take the convention 
\bse 
    \int d\theta^1 d\theta^2 \, \theta^2 \theta^1 = 1,
\ese 
i.e. we "do the inner integral first". Note that this implies 
\bse 
    \int d\theta^1 d\theta^2 \theta^1 \theta^2 = - \int d\theta^1d\theta^2 \theta^2\theta^1 = -1,
\ese 
where we have used $\theta^1\theta^2 = - \theta^2\theta^1$. 

Collectively, then, we have 
\bse 
    \int d^2\theta (\theta\theta) = \int d^2\bar{\theta} (\bar{\theta}\bar{\theta}) = 1 
\ese 
and 
\bse 
    \int d^2\theta d^2 \bar{\theta} (\theta\theta)(\bar{\theta}\bar{\theta}) = 1.
\ese 

\bcl 
    We can use the above results to rewrite our integrals as 
    \be 
    \label{eqn:d2ThetaAsDerivatives}
        \int d^2 \theta = \frac{1}{4} \epsilon^{\a\beta} \p_{\a}\p_{\beta} \qand \int d^2 \bar{\theta} = -\frac{1}{4} \epsilon^{\dot{\a}\dot{\beta}} \bar{\p}_{\dot{\a}}\bar{\p}_{\dot{\beta}}.
    \ee 
\ecl 

\bq 
    Consider the $d^2\theta$ case:
    \bse 
        \begin{split}
            \int d^2\theta \, \theta\theta & =: \int d^2 \theta \, \theta^{\a}\theta_{\a} \\
            & = \int d^2\theta \, \epsilon_{\a\beta}\theta^{\a}\theta^{\beta} \\
            & = \int d^2\theta \, \big[-\theta^1\theta^2 + \theta^2\theta^1\big] \\
            & = \int d^2\theta \, 2\theta^2\theta^1,
        \end{split}
    \ese 
    where we have used 
    \bse 
        \epsilon_{\a\beta} = \begin{pmatrix}
            0 & -1 \\
            1 & 0
        \end{pmatrix}.
    \ese 
    which we want to be $1$ by our convention above. So combining this with the fact that our integrals act like derivatives, we see we need\footnote{Note we could have essentially guessed this from \Cref{eqn:d2Theta}.} 
    \bse
        \begin{split}
            \int d^2 \theta & = \frac{1}{2}\p_1\p_2 \\
            & = \frac{1}{4} \big(\p_1\p_2 + \p_1\p_2 \big) \\
            & = \frac{1}{4} \big(\p_1\p_2 - \p_2\p_1\big) \\
            & = \frac{1}{4} \epsilon^{\a\beta} \p_{\a}\p_{\beta},
        \end{split}
    \ese 
    where we have used $\epsilon^{\a\beta}= (\epsilon_{\a\beta})^{-1}$, so the minus sign swaps. This is the result we wanted. The barred version follows trivially from here we have the $\dot{2}$ integral to the left and so get a minus sign difference. 
\eq 

\br 
\label{rem:GrassmanIntegralSurvive}
    The important point to note about integrals over Grassman variables is that the only terms in the integrand that survive the integral are the ones with the matching Grassman structure. By which we mean if we do the integral over 
    \bse 
        \int d^2 \theta \, \Phi(\theta,\bar{\theta}),
    \ese
    only the term in $\Phi$ that contains two $\theta$s and no $\bar{\theta}$s will survive. The fact that we need two $\theta$s is clear from our "integrals act like derivatives" argument. The reason we don't want any $\bar{\theta}$s is that $\bar{\theta}$ is just a constant w.r.t. $d\theta$, and so by the first condition in \Cref{eqn:GrassmanIntegralDefinition} this vanishes. This idea will prove very useful to us going forward. 
\er 

\subsection{Manifestly Supersymmetric Integrals}

Ok great, so now we know how to integrate over the Grassman part of our superspace so we can begin to try and construct manifestily SUSY integrals. There are indeed two types:
\ben 
    \item Integral over \textit{all of superspace}:
    \bse 
        \int d^4x  d^2\theta d^2\bar{\theta} \, Y(x,\theta,\bar{\theta})
    \ese 
    with $Y(x,\theta,\bar{\theta})$ being a general superfield. This is is manifestly SUSY because the SUSY variation of $Y$ takes the form
    \bse 
        \del_{\epsilon,\bar{\epsilon}} Y = i(\epsilon Q + \bar{\epsilon}\bar{Q})Y = \p_{\a}( ... )^{\a} + \bar{\p}_{\dot{\a}} (...)^{\dot{\a}} + \p_{\mu}(...)^{\mu},
    \ese 
    but this is just a total derivative in superspace and so it must vanish. 
    \item Integral over \textit{chiral half} of superspace:
    \bse 
        \int d^4 y d^2\theta \,  W(y,\theta) = \int d^4 x d^2\theta \, W(x,\theta,\bar{\theta})
    \ese 
    with $W(y,\theta)$ being a chiral superfield. Again this is SUSY because 
    \bse 
        \del_{\epsilon,\bar{\epsilon}}W(y,\theta) = \p_{\a} (...)^{\a} + \frac{\p}{\p y^{\mu}} (...)^{\mu},
    \ese 
    which again is a total derivative in chiral half of superspace.\footnote{From now on we might simply write "chiral superspace" to mean "chiral half of superspace".} We obviously have a similar thing for the antichiral cases.  
\een

\bp  
    Any integral over superspace can be written as an integral over chiral superspace.
\ep 

\bq 
    At first this seems highly unintuitive and feels like it should be the other way. The proof comes from using \Cref{eqn:d2ThetaAsDerivatives}: we have 
    \bse 
        \int d^2\bar{\theta} = -\frac{1}{4}\epsilon^{\dot{\a}\dot{\beta}} \bar{\p}_{\dot{\a}} \bar{\p}_{\dot{\beta}} = -\frac{1}{4} \epsilon^{\dot{\a}\dot{\beta}} \bar{D}_{\dot{\a}}\bar{D}_{\dot{\beta}} + \text{total deriv. in } x,
    \ese
    and so we can write an integral over full superspace as 
    \bse 
        \int d^4 x d^2\theta d^2\bar{\theta} Y = - \frac{1}{4} \int d^4x d^2\theta \bar{D}^2 Y.
    \ese 
    This gives us the integral over the chiral half, but we still need to show that our integrand is a chiral superfield. This follows trivially from the fact that 
    \bse 
        \bar{D}_{\dot{\a}}\bar{D}^2 Y = 0
    \ese 
    for \textit{any} superfield $Y$ as there are too many $\theta$s. This tells us we can express any integral over all superspace as an integral over chiral superspace, and it is clear that the reverse is not true. This is just because if we have an integral over chiral superspace who's integrand is \textit{not} chirally exact we can't work backwards to obtain an integral over all superspace. Using the language that follows in a second, an integral over chiral superspace of a chiral superfield that is \textit{not} chirally exact can \textit{not} be expressed as the integral over full superspace of a general superfield. 
\eq 

\bd[Chirally Exact Superfield]
    We call a superfield of the form 
    \be 
    \label{eqn:ChirallyExact}
        \chi = \bar{D}^2 Y,
    \ee 
    with $Y$ being a general superfield, \textit{chirally exact}. 
\ed 

Now with \Cref{rem:GrassmanIntegralSurvive} in mind, we see that 
\bse 
    \int d^4 x d^2\theta d^2\bar{\theta} \, Y(x,\theta,\bar{\theta}) = \int d^4x \, D(x) \qand \int d^4 y d^2\theta \, W(y,\theta) = \int d^4x \, F_W(x)
\ese
where $D(x)$ is the top component of $Y$ and $F_W(x)$ is the top component of $W$. We therefore refer to 
\begin{itemize}
    \item F-terms: Integrals over chiral superspace with a non-chirally exact integral (i.e. of type 2. but not type 1.)
    \item D-terms: Integrals over full superspace (i.e. integrals of type 1.)
\end{itemize}

\subsection{Supersymmetric Actions}

Finally before to starting to study more specific examples of SQFTs, let's just make some comments on the conditions for us to have a SUSY action. 
\ben[label=(\roman*)]
    \item The action is meant to be real and so we require our $D$-terms to be real. Our $F$-terms can be complex as long as they are accompanied with their Hermitian conjugate term, i.e. the $\bar{F}$-terms arising from integrals over antichiral superspace.
    \item As they are top components they must be scalars (i.e. Bosonic). 
    \item The \textit{engineering dimension}\footnote{This is just the usual mass dimension in QFT. The reason we say "engineering" is to distinguish it from dimension in CFT, which is the dilatation weight. These two things need not agree for non-free fields} of $Y$ and $W$ are fixed by the fact that $[S]=0$. In particular we have 
    \bse 
        [P_{\mu}] = 1 \quad \implies \quad [Q] = [\bar{Q}] = 1/2,
    \ese 
    as the anticommutator of two $Q$s is $P$. From here (with $[x^{\mu}]=-1$) can obtain 
    \be 
    \label{eqn:ThetaDimension}
        [\theta] = [\bar{\theta}] = -1/2.
    \ee 
    Now we have to be careful and remember that the integral over Grassman variable is like a derivative and so we have 
    \bse 
        \bigg[\int dx\bigg] = -1, \qquad \text{but} \qquad \bigg[\int d\theta\bigg] = +1/2.
    \ese 
    Putting this together we concluse
    \mybox{
        \be 
        \label{eqn:YWDimensions}
            [Y] = 2 \qand [W] = 3.
        \ee 
    }
    Note we can also obtain this from $[F]=[D]=4$ (from the fact that they are integrated over $d^4x$ only) along with the fact that they appear as $\theta\theta\bar{\theta}\bar{\theta}D$ and $\theta\theta F$ and \Cref{eqn:ThetaDimension}. 
\een 
\chapter{SQFT Of Chiral Multiplets}

Now that we know how to construct a SUSY action, we can begin to actually study supersymmetric quantum field theories. Of course we will focus on our chiral superfields as we know these will form irreps. There are two main types
\ben 
    \item Wess-Zumino models: these are renormalisable and will be our main focus.
    \item Non-Linear Sigma models: these are non-renormalisable. We will discuss these a bit too.\footnote{\textcolor{red}{These were not part of the taught material but there's some stuff in Stefano's notes about it, so I shall add that later.}}
\een

\section{Content of SQFT}

Before moving on to discuss the above models, first let's break down the content of a SQFT. 

\subsubsection{Field Content}

The first obvious thing to as is what is the field content? As we have tried to stress above, if we want to get irreps, we need to restrict to chiral superfields, so our field content is simply a chiral superfield $\Phi^i$
\bse 
    \Phi^i(y,\theta) = \phi^i(y) + \sqrt{2} \theta\psi^i(y) - \theta\theta F^i(y) 
\ese 
and the antichiral superfield we get by taking the Hermitian conjugate $\bar{\Phi}^{\bar{i}} := (\Phi^i)^{\dagger}$:\footnote{We use a bar over the $i$ index instead of a dot just because $\dot{i}$ isn't very nice.}
\bse 
    \bar{\Phi}^{\bar{i}}(\bar{y},\bar{\theta}) = \bar{\phi}^{\bar{i}}(\bar{y}) + \sqrt{2} \bar{\theta} \bar{\psi}^{\bar{i}}(\bar{y}) - \bar{\theta}\bar{\theta} \bar{F}^{\bar{i}}(\bar{y}). 
\ese 

\subsubsection{$R$-Symmetry}

Next, what is our $R$ symmetry? Well we are considering $\cN=1$, $4D$ theories and so we simply have $U(1)_R$. We will denote the generator simply by $R$. It obeys
\bse
    [R,Q_{\a}] = - Q_{\a} \qand [R,\bar{Q}_{\dot{\a}}] = + Q_{\a}
\ese 
where the $\pm$ are the \textit{$R$-charges} of $Q/\bar{Q}$. From here we can we can work out the $R$-charge for $\theta/\bar{\theta}$ simply as
\bse 
    R[\theta^{\a}] = +1   \qand R[\bar{\theta}^{\dot{\a}}] =  -1
\ese 
where we note the sign flips compared to $Q/\bar{Q}$.

How does our $U(1)_R$ act on the fields? Well we express the result in terms of the $R$-charge of the chiral superfield $R[\Phi]$, i.e. 
\bse 
    \Phi \mapsto e^{iR[\Phi]\a}\Phi.
\ese 
We then just use you component decomposition above along with 
\bse 
    \theta  \mapsto e^{i\a} \theta
\ese 
to obtain
\mybox{
    \be 
    \label{eqn:RChargeOnFields}
        \begin{split}
            \phi & \mapsto e^{iR[\Phi]\a} \phi  \\
            \psi & \mapsto e^{i(R[\Phi]-1)\a} \psi\\
            F & \mapsto e^{i(R[\Phi]-2)\a} F.
        \end{split}
    \ee 
}
\noindent Then the Hermitian conjugates simply transform with the opposite sign (as $R[\theta]=-R[\bar{\theta}]$). 

\subsubsection{Global Flavour Symetries}

Recall that our global flavour symmetries are defined as those automorphisms which commute with the central charges. If we denote the generators by $F_I$, we have
\bse 
    [F_I, Q_{\a}] = [F_I, \bar{Q}_{\dot{\a}}] = 0 \qquad \implies \qquad F_I[\theta] = F_I[\bar{\theta}],
\ese 
and so \textit{all the components} have the \textit{same} $F_I$-charges.

\section{Most General SUSY Action}
\label{sec:MostGeneralSUSYAction}

We are now ready to write down the most general SUSY action for a Chiral superfield \textit{with at most 2 derivatives}. The answer is simply
\mybox{
    \be 
    \label{eqn:MostGeneralSUSYAction}
        S = \int d^4 d^2 \theta d^2\bar{\theta} \, K(\Phi,\bar{\Phi}) + \int d^4 y d^2 \theta \, W(\Phi) + \int d^4 y d^2 \bar{\theta} \, \bar{W}(\bar{\Phi}),
    \ee 
}
\noindent which is just a combination of a $D$-term, a $F$-term and $\bar{F}$-term. Let's make some comments.

\ben[label=(\roman*)]
    \item $K(\Phi,\bar{\Phi})$ is the \textit{K\"{a}hler potential}. We know this must be a real function of $\Phi^i$ and $\bar{\Phi}^{\bar{i}}$ of engineering dimension $2$. We call it a composite \textit{real} superfield. We can consider \textit{K\"{a}hler transformations} which shift the K\"{a}hler potential:
    \bse 
        K(\Phi,\bar{\Phi}) \to K(\Phi,\bar{\Phi}) + \Lambda(\Phi) + \bar{\Lambda}(\bar{\Phi})
    \ese 
    where $\Lambda/\bar{\Lambda}$ are holomorphic/antiholomorphic that are also (anti)chiral.\footnote{We shall prove this in a moment.} This leaves the action invariant because 
    \bse 
        \int d^4 d^2\theta d^2\bar{\theta} \Lambda(\Phi) = -\frac{1}{4} \int d^4 y d^2\theta \, \bar{D}^2 \Lambda(\Phi) = 0. 
    \ese
    The physical (invariant) quantity is 
    \bse 
        g_{i\bar{j}}(\Phi,\bar{\Phi}) = \p_i \bar{\p}_{\bar{j}} K(\Phi,\bar{\Phi})
    \ese 
    where 
    \bse 
        \p_i := \frac{\p}{\p \Phi^i} \qand \bar{\p}_{\bar{j}} := \frac{\p}{\p \bar{\Phi}^{\bar{j}}}.
    \ese 
    $g_{i\bar{j}}$ is known as the \textit{K\"{a}hler metric}. It is the metric on the complex manifold parameterised by the complex coordinates $\Phi^i$. 
    \item $W(\Phi)$ is called a \textit{superpotential}. This has to be a holomorphic function of $\Phi^i$ with $[W]=3$. It is called a composite \textit{chiral} superfield. 
\een 

Twice above we used the fact that a holomorphic function of a chiral superfield is a chiral superfield. Let's now actually prove this. 
\bse 
    \begin{split}
        \bar{D}_{\dot{\a}} W(\Phi) & = (\bar{D}_{\dot{\a}} \Phi^i) \p_i W + (\bar{D}_{\dot{\a}} \bar{\Phi}^{\bar{i}}) \bar{\p}_{\bar{i}} W \\
        & = 0 + 0,
    \end{split}
\ese
where the second term on the first line appears from the idea that 
\bse 
    \frac{\p}{\p \bar{z}}\frac{1}{z} \neq 0.
\ese 
The first 0 comes from it being Chiral and the second comes from $W$ being holomorphic, i.e. derivative w.r.t to conjugate variable vanishes.

%%% HERE %%%% 

\section{Global Symmetries of $S$}

The first thing to consider is our $R$-symmetries. Our action must have $R[S]=0$, and so using $R[d^2\theta] = -2$ and $R[d^2\bar{\theta}]= +2$ we can easily prove that 
\mybox{
    \be
    \label{eqn:RChargeKW}
        R[K] = 0 \qand R[W] = 2.
    \ee 
}
\noindent An analogous calculation then let's us also conclude 
\mybox{
    \be 
    \label{eqn:FChargeKW}
        F_I[W] = F_I[K] .
    \ee 
}
\noindent This means in particular if you have an abelian flavour symmetry that $F_I[K] = F_I[W]=0$, but for non-abelian its just that they are singlets. 

Now recall that a symmetry is \textit{explicitly broken} if we include some term in the action by hand that doesn't obey the symmetry. For example 
\bse 
    S = \int d^4 x \, (\p\phi)^2 + \l_4 \phi^4
\ese 
has a global $\Z_2$ symmetry given by $\phi\to-\phi$. This is explicitly broken if we introduce a $\l_3\phi^3$ term. 

We can reword this condition for us by saying that our $R$/$F_I$ symmetries are explictly broken if \Cref{eqn:RChargeKW}/(\ref{eqn:FChargeKW}) can only be met by assigning \textit{non-zero charges to the parameters}, e.g. the $\l_4/\l_3$ in the example above. 

\bex
    Let's consider 
    \bse 
        K(\Phi,\bar{\Phi}) = \bar{\Phi}\Phi \qand W(\Phi) = \frac{m}{2}\Phi^2 + \frac{\l}{3} \Phi^3.
    \ese 
    The K\"{a}hler potential condition tells us we need $R[\Phi]=1$. Given that $R[m]=0$, we see that our $U(1)_R$ symmetry then is \textit{only} preserved if $\l=0$. However if $\l\neq 0$ it is explicitly broken, because then $R[W] \overset{!}{=} 2$ requires $R[\l]=-1 \neq 0$. 
\eex 

\subsection{Spurion Analysis}

This seems like a very strange idea. In other words, why should we be assigning charges to the parameters? The answer is an old idea\footnote{Well at least as old as the chiral Lagrangian.} in QFT: we view parameters that explicitly break a symmetry as fixed \textit{background} values of external (or non-dynamical) fields that are charged under the broken symmetry. They are called \textit{spurions}. If the parameters did indeed correspond to dynamical fields, the symmetry would be broken spontaneously. The idea is that a low energy observer can't know if a parameter actually \textit{is} the vev of a field which is dynamical at higher energies or not. For this reason, we must treat parameters (non-dynamical fields) on the same footing as dynamical fields. That is we promote parameters to background values of non-dynamical fields, when we do that we restore the symmetry, which leads to selection rules that constrain the symmetry.

In SQFT, superpotential parameters will be treated on the same footing as background chiral superfields. This idea is at the root of one of the most powerful results about SQFTs: the \textit{non-renormalisation theorem for the superpotential}. We will discuss this later. 

\section{Wess-Zumino Models}

Ok let's now specialise to our Wess-Zumino (WZ) models and actually calculate some stuff. As we said before, these are renormalisable theories of chiral superfields. The components will have canonical dimensions, 
\bse 
    [\phi]= 1, \quad [\psi] = \frac{3}{2} \qand [F] = 2.
\ese 
We then require that $[\cL]=4$, which implies that $[K]=2$ and $[W]=3$, which restricts\footnote{Technically we have only restricted $K$ to be quadratic and terms like $\Phi\Phi + \bar{\Phi}\bar{\Phi}$ would be allowed. However it turns out that we can remove such terms using a K\"{a}hler transformation and diagonalise the $\bar{\Phi}\Phi$ terms, as in the equation here.}
\mybox{
    \be
    \label{eqn:KahlerPotentialWZ}
        K = \sum_i \bar{\Phi}^{\bar{i}} \Phi^i \equiv \bar{\Phi}_i \Phi^i
    \ee
}
\noindent this is the \textit{canonical K\"{a}hler potential}.

\br 
    The index above has been lowered using the K\"{a}hler metric, $g_{i\bar{j}}$ introduced above. 
\er 

Similarly, we also have 
\mybox{
    \be 
    \label{eqn:SuperpotentialWZ}
        W = \frac{1}{2}m_{ij} \Phi^i\Phi^j + \frac{1}{3}\l_{ijk} \Phi^i\Phi^j\Phi^k
    \ee 
}
\noindent which is the cubic \textit{superpotential}. 

\br 
    Note we don't include linear terms because we could always absorb them into the quadratic term using a K\"{a}hler transformation. Also constant terms vanish when we integrate over $\theta/\bar{\theta}$. 
\er 

We can expand these in components by Taylor expanding around $\phi(y)$ and truncating the result at $\theta\theta$. As $\psi(y)$ comes with a single $\theta$ we will get two terms in the Taylor expansion, whereas we will only get one from the $F(y)$ expansion. The result is 
\bse 
    \begin{split}
        W\big(\Phi(y,\theta)\big) & = W\big( \phi(y) + \sqrt{2}\theta \psi(y) - \theta\theta F(y)\big) \\
        & = W\big(\phi(y)\big) + \sqrt{2} \p_i W\big(\phi(y)\big) \theta \psi^i(y) - \theta\theta \big[ \p_i W\big(\phi(y)\big) F^i(y) + \frac{1}{2}\p_i\p_j W\big(\phi(y)\big) \psi^i(y)\psi^j(y)\big]
    \end{split}
\ese
where again we have used the identity 
\bse 
    (\theta\psi)(\theta\chi) = -\frac{1}{2}(\theta\theta)(\psi\chi)
\ese 
to get the $\p_i\p_j$ term. We have the bottom, middle and top components
\bse 
    W(\phi(y)), \qquad \p_i W(\phi(y))\psi^i(y) \qand \big[ \p_i W\big(\phi(y)\big) F^i(y) + \frac{1}{2}\p_i\p_j W\big(\phi(y)\big) \psi^i(y)\psi^j(y)\big]
\ese 
respectively. The most important one for us is the top component 
\bse 
    F_W =  \big(\p_i W(\phi)\big) F^i + \frac{1}{2}\big(\p_i\p_j W(\phi)\big) \psi^i\psi^j,
\ese 
which obeys
\bse 
    \int d^4x d^2 \theta \, W(\Phi) = - \int d^4 x \, F_W,
\ese
with the minus sign coming from our expansion above. 

This is all general, so we can now use our specific potential:\footnote{Note that it is little $\phi$ here as we have done the expansion around the bottom component $\phi$ so this is the argument.} 
\bse 
    \p_iW(\phi(y)) = m_{ij}\phi^i + \l_{ijk}\phi^j\phi^k \qand \p_i\p_j W = m_{ij} + 2\l_{ijk}\phi^k.
\ese 
Similarly we have 
\bse 
    \int d^4x d^2\theta d^2\bar{\theta} K = \int d^4x D_K
\ese
where $D_K$ is the top component of $K$ (i.e. the one with $\theta\theta\bar{\theta}\bar{\theta}$).

\bbox 
    By plugging in the component expansions into $K=\bar{\Phi}\Phi$, show that 
    \be 
    \label{eqn:DPhiPhiExercise}
        D_{\bar{\Phi}\Phi} = \p_{\mu}\bar{\phi}\p^{\mu}\phi - i\psi \sig^{\mu}\p_{\mu}\bar{\psi} + \bar{F}F + \text{total derivative} 
    \ee 
\ebox 

The middle term in the exercise above can be written as
\bse 
    -i\bar{\psi}\bar{\sig}^{\mu}\p_{\mu}\psi + \text{total deriv},
\ese 
which we shall do in the following. Collecting these results:
\bse 
    \begin{split}
        S & = \int d^4 x d^2\theta d^2\bar{\theta} \bar{\Phi}_i \Phi^i + \int d^4 d^2\theta W(\Phi) + h.c. \\
        & = \int d^4 x\bigg[ \p_{\mu}\bar{\phi}_i \p^{\mu}\phi^i - i\bar{\psi}_i \bar{\sig}^{\mu}\p_{\mu}\psi^i + \bar{F}_iF^i - (\p_i W) F^i - (\bar{\p}^i \bar{W}) \bar{F}_i - \frac{1}{2}(\p_i\p_j W) \psi^i\psi^j - \frac{1}{2}(\bar{\p}_i\bar{\p}_j \bar{W}) \bar{\psi}^i\bar{\psi}^j\bigg]
    \end{split}
\ese 
We can then group the terms as follows:
\begin{itemize}
    \item Kinetic terms: $\p_{\mu}\bar{\phi}_i \p^{\mu}\phi^i - i\bar{\psi}_i \bar{\sig}^{\mu}\p_{\mu}\psi^i$ 
    \item Scalar potential: $\bar{F}_iF^i - (\p_i W) F^i - (\bar{\p}^i \bar{W}) \bar{F}_i$. 
    \item Yakawa terms: $- \frac{1}{2}(\p_i\p_j W) \psi^i\psi^j - \frac{1}{2}(\bar{\p}_i\bar{\p}_j \bar{W}) \bar{\psi}^i\bar{\psi}^j$. We can see that these are indeed Yakawa terms by plugging in our $\p_i\p_j W$ term, which gives us a $\phi\psi\psi$ type term. Note they also give us a mass term for the $\psi$s. 
\end{itemize}

\br 
    We were careful to define this already, but just as a reminder: be careful not to mix $\p_{\mu}$ with $\p_i$. The former is a spacetime derivative, second is w.r.t $\phi^i$.
\er 

\subsection{Scalar Potential}

Note that $[F]=2$ which tells us that they will only appear quadratically and with out kinetic terms in the action. This is why they are called auxillary fields:\footnote{They are exactly those fields we mentioned in the footnotes last chapter needed to equate the off-shell Fermion and Boson degrees of freedom.} they can be integrated out exactly in the path integral, as they are just Gaussian. This has the same effect as replacing them by the solution to the equations of motion 
\mybox{
    \be
    \label{eqn:FEoM}
        F^i = \bar{\p}^i \bar{W}^i \qand \bar{F}_i = \p_i W.
    \ee 
}
\noindent When this is done we obtain the scalar potential 
\mybox{
    \be 
    \label{eqn:ScalarPotential}
        V(\phi,\bar{\phi}) = \sum_i |F^i|^2 = \sum_i |\p_i W|^2 
    \ee 
}
\noindent where the $F^i$s are \textit{on-shell} (as we have imposed the EoM). This will lead to mass terms, cubic terms and quartic terms, which go as 
\bse 
    m \bar{m} \bar{\phi}\phi \qquad m \bar{\l} \phi \bar{\phi}^2 + h.c. \qand \l\bar{\l} \phi^2\bar{\phi}^2.
\ese 

\subsection{Supersymmetric Vacua}

We now want to construct our SUSY vacua. 

\bcl 
    A Lorentz invariant vacuum requires the vev of a non-trivial Lorentz field to be vanishing, whereas we only require Lorentz scalars to have constant vev.  
\ecl 

\bq 
    The idea is based around our fields gaining a vev, e.g. 
    \bse 
        \phi \to \la \phi \ra + \phi \qand \psi \to \la \psi \ra + \psi.
    \ese 
    The first thing to note is that the vev itself is Lorentz invariant (it's a number). We then substitute these into the Lagrangian and then insist that we maintain Lorentz invariance. For example let's consider a mass term (any term in the Lagrangian will do, of course), and consider the scalar first:
    \bse 
        m^2\phi^2 \to m^2 (\la \phi\ra + \phi)^2 = m^2 \big(  \la\phi\ra^2 + 2 \la \phi \ra \phi + \phi^2\big),
    \ese 
    and every term here is Lorentz invariant. So we can have any constant number for the vev. However for something with non-trivial Lorentz transformation, e.g. a Fermion, we have 
    \bse 
        m \bar{\psi}\psi \to m \big(\la \bar{\psi}\ra + \bar{\psi}\big)\big(\la\psi\ra + \psi) = m \big( ... + \la \bar{\psi}\ra \psi + \la\psi\ra \bar{\psi} + ... \big),
    \ese 
    where all the terms we've dropped are Lorentz invariant. However, as we've said, the vev itself is Lorentz invariant and so the products written above are \textit{not} Lorentz invariant. So the only way we can obtain a Lorentz invariant theory is to make the vevs vanish. 
\eq 

Using the above idea, and the fact we've seen that we can have multiple SUSY vacua, we define the set of SUSY vacua as
\mybox{
    \be
    \label{eqn:SUSYVacua}
        \cM := \Big\{ \la \phi^i\ra = \text{const} \, \Big| \, \p_i W \big|_{\la\phi^i\ra} = 0 \quad  \forall i \Big\}
    \ee 
}
\noindent where the condition $\p_iW=0$ comes from wanting to minimise the scalar potential. Note that this condition gives us $\la F^i \ra =0$\footnote{So the constant for these scalars is $0$.} when we take the $F$ on-shell. 

\Cref{eqn:SUSYVacua} are indeed SUSY invariant by the following results:
\bse 
    \begin{split}
        \del_{\epsilon,\bar{\epsilon}} \la \phi^i \ra & \sim \la \psi^i \ra = 0  \\
        \del_{\epsilon,\bar{\epsilon}} \la \psi^i \ra & \sim \la \p_{\mu} \phi^i \ra + \la F^i \ra = 0  \\
        \del_{\epsilon,\bar{\epsilon}} \la F^i \ra & \sim \la \p_{\mu} \psi^i \ra = 0 
    \end{split}
\ese 

\br 
    Note that the supersymmetry vacua are zeros of the energy, which is in agreement with the fact SUSY is unbroken if and only if the vacuum energy vanishes. 
\er 

Often SUSY theories have exactly \textit{flat directions} of the scalar potential. The massless fields which parameterise the flat directions are called \textit{moduli} and the set of supersymmetric vacua $\cM$ is called \textit{moduli space of SUSY vacua}. We will see more about this later.

\subsection{EOM For Chiral Superfields}

We now want to find the equation of motion for a chiral superfield.\footnote{A more detailed analysis can be found in Chapter 9 of Wess \& Bagger} We'd like to vary the action
\bse 
    S = \int d^4 x d^2 \theta d^2\bar{\theta} \, \bar{\Phi}_i \Phi^i + \int d^4 xd^2\theta \, W(\Phi) + h.c.
\ese
w.r.t $\Phi^i$. Naively we would write 
\bse 
    \bar{\Phi}_i + \p_i W(\phi) = 0,
\ese 
however this must be wrong because if we considered the free theory (i.e. $W(\Phi)=0$), we would obtain $\Phi=0$, but this is not the equation of motion for a free Boson theory. 

Hmm... so what did we do wrong? The problem is that our $F$-term integral, $\int d^4xd^2\theta W(\Phi)$, is restricted to being a chiral superfield whereas the $D$-term integral, $\int d^4xd^2\theta d^2\bar{\theta}\bar{\Phi}_i\Phi^i$, is a general superfield. In other words, we have ignored the fact that one term in our action is integrated over \textit{full superspace} while the other was only over the \textit{chiral half}. We therefore want to put them on the same footing. The good thing is we already know how to do this, we simply write the full superspace integral as an chirally exact expression. We then get 
\bse 
    S = \int d^4 x d^2\theta \bigg(-\frac{1}{4}\bar{D}^2\big(\bar{\Phi}_i\Phi^i\big) + W(\Phi) \bigg) + \int d^4 x d^2\bar{\theta} \, \bar{W}(\bar{\Phi}).
\ese
Now use Liebniz\footnote{Note that this is a second order differential equation so its not simply $\bar{D}^2(\bar{\Phi})\Phi + \bar{\Phi}\bar{D}^2(\Phi)$, but we will also have $2(\bar{D}\bar{\Phi})(\bar{D}\Phi)$, but this term still vanishes by our chiral condition.} and the fact that we have chiral fields (so $\bar{D}_{\a}\Phi=0$) to obtain 
\bse 
    S = \int d^4 x d^2\theta \bigg(-\frac{1}{4}\bar{D}^2\big(\bar{\Phi}_i\big) \Phi^i + W(\Phi) \bigg) + \int d^4 x d^2\bar{\theta} \, \bar{W}(\bar{\Phi}) 
\ese 
\textit{Now} vary w.t.t. $\Phi^i$ to get 
\mybox{
    \be 
    \label{eqn:ChiralSuperfieldEOM}
        \frac{1}{4}\bar{D}^2 \bar{\Phi}_i = \p_i W \qand \frac{1}{4}D^2 \Phi^i = \bar{\p}^i \bar{W}.
    \ee 
}
\noindent where the second expression comes from if we had done the same thing for the antichiral case. 

\br 
    Note that if we have unbroken SUSY, i.e. \Cref{eqn:SUSYVacua} is not the empty set, the \Cref{eqn:ChiralSuperfieldEOM} let's us conclude 
    \bse 
        \la \bar{D}^2 \bar{\Phi}_i \ra = 0.
    \ese 
    In other words, \textit{any} chirally exact super field has vanishing vev, provided you have SUSY. 
\er 

\bbox 
    Consider the WZ model for a single chiral superfield $\Phi$ with 
    \bse 
        K(\bar{\Phi},\Phi) = \bar{\Phi}\Phi \qand W(\Phi) = \frac{m}{2}\Phi^2 + \frac{\l}{3}\Phi^3.
    \ese 
    \ben 
        \item Argue that the $W(\Phi)$ given above is the most general renormalisable superpotential.
        \item Find the SUSY vacua of this theory.
        \item Write down the Lagrangian in components, before and after integrating out the auxiliary fields. Check that $\phi$ and $\psi$ have the same bare mass $m$, and the same effective complex mass $m^{\C}_{\text{eff}}(\la\phi\ra) = m + 2\l\la\phi\ra$\footnote{Note the mass might change depending on the vacua, but they will always have the same mass as each other.} when the Lagrangian is expanded about a vacuum where $\phi$ takes a vev. How is the quartic coupling in the scalar potential related to the Yukawa coupling?
        \item Derive the EoM for the component fields $\phi,\psi$ and $F$ from the Lagrangian written in components.
        \item Expand the super-EoM for the superfield, $\Phi$, in components and rederive the EoM for the component fields derived above.
    \een 
    The last bit is meant to make you appreciate that working with superfields rather then components. 
\ebox  

\section{Non-Linear Sigma Models*}

\textcolor{red}{This material was not lectured and the material in the notes is very bare-bones (essentially just the final result). I will try come back and include some material for this section at a later date. For the time being the interested reader is directed to Section 7 of Bilal. }

\section{Non-Renormalisation Theorem (Seiberg 1993)}

We now want to prove an important theorem which states that the superpotential of a theory of chiral superfields does not flow under RG.\footnote{For more details on what this means, see the Renormalisation Group course.} That is, we'll use SUSY and holomorphy of $W$ to show that the superpotential is \textit{not} renormalised.

\br 
    There is a perturbative proof using so-called \textit{supergraphs}, but here we will use the more modern version due to Sieberg, which is a fully non-perturbative statement. 
\er 

Let's quickly recap what RG flow is:\footnote{Again, for more details see the RG course.} it is a statement about the \textit{Wilsonian effective action} (WEA), that is obtained by integrating our modes with Euclidean momenta $|p|>\mu$, where $\mu$ is our cut-off scale. Given the WEA at the cut-off scale $\mu$, we can obtain the WEA at a scale $\mu -d\mu$ by integrating out modes with $\mu-d\mu < |p| < \mu$. In the process, the values of the couplings/wave-functions could change, leading to a renormalisation of the couplings/wave-functions. This is one iteration step in our RG flow, and we repeatedly do this to obtain the IR behaviour. This is the basic idea behind \textit{Wilsonian renormalisation group}. 

For a SUSY field theory of chiral multiplets, the WEA takes the general form 
\bse 
    S_{\text{eff},\mu} = \int d^4 x d^2\theta d^2\bar{\theta} \, K_{\text{eff},\mu} + \int d^4 x d^2 \theta \, W_{\text{eff},\mu} + \int d^4 x d^2\bar{\theta} \, \bar{W}_{\text{eff},\mu},
\ese
up to higher order derivative terms. As always, we start from some max UV value $\mu =\Lambda_{\text{UV}}$ where the theory is defined. We then define
\bse 
    K_{\text{micro}} := K_{\text{eff},\Lambda_{\text{UV}}}, \qquad  W_{\text{micro}} := W_{\text{eff},\Lambda_{\text{UV}}}, \qand \bar{W}_{\text{micro}} := \bar{W}_{\text{eff},\Lambda_{\text{UV}}}.
\ese 
From now we will just write the $W$ expression and assume the $\bar{W}$ expression is implied. 

Due to quantum corrections, we would expect that when we do an RG iteration to scale $\mu$, that 
\bse 
    K_{\text{eff},\mu} \neq K_{\textbf{micro}} \qand W_{\text{eff},\mu} \neq W_{\text{micro}}. 
\ese 
This is the statement that we expect the fields to \textit{flow} under RG, in other words the values of $K$ and $W$ depend on the scale $\mu$. As we said at the start of this section, this is \textit{not} the case for the superpotential, $W$, and is the content of the next theorem.  

\bt[Non-Renormalisation Theorm]
    The superpotential of a theory of chiral superfields does \textit{not} flow under RG. That is 
    \be 
    \label{eqn:NonRenormThrm}
        W_{\text{eff},\mu} = W_{\text{micro}},
    \ee 
    for all $\mu$.
\et 

\bq 
    The key ideas needed to prove this theorem are 
    \ben[label=(\roman*)] 
        \item Holomorphy in the microscopic coupling constant
        \item Selection rules from symmetries under which the microscopic coupling constants may transform. 
        \item Smoothness of physics in various weak coupling limits, where we know how the theory should behave. 
    \een
    The first two follow from spurion analysis, i.e. from viewing all superpotential coupling constants as background chiral superfields.
    
    For this proof, we will consider the simplest case\footnote{Of course this itself does not prove the theorem, but for us it'll do.} of the WZ model of a single field
    \be 
    \label{eqn:WMicro}
        W_{\text{micro}} = \frac{1}{2}m\Phi^2 + \frac{1}{3}\l \Phi^3.
    \ee 
    \ben[label=(\roman*)]
        \item Holomorphy tells us $W_{\text{eff}}= f(\Phi,m,\l)$, which is fully holomorphic (i.e. no barred dependence). We leave the $\mu$ dependence of $f$ implicit. 
        \item From our spurion analysis conversation before, we allow $m$ and $\l$ to be charged under our symmetries. Then from $F[W]=0$ and $R[W]=2$, a quick calculation gives 
        \begin{center}
        	\begin{tabular}{@{} C{2cm} C{2cm} C{2cm}  @{}}
        		\toprule
        		 & $U(1)_F$ & $U(1)_R$ \\
        		\midrule 
        		$\Phi$ & $1$ & $1$ \\
        		$m$ & $-2$ & $0$ \\
        		$\l$ & $-3$ & $-1$ \\
        		$(\mu$ & $0$ & $0)$ \\
        		\bottomrule
        	\end{tabular}
        \end{center}
        where we have also included the charges for $\mu$. Now clearly our effective $W_{\text{eff},\mu}$ must also have $R$-charge $2$ and vanishing $F$-charge, so it takes the general form 
        \bse
            W_{\text{eff},\mu} = m\Phi^2 \cdot g\bigg(\frac{\l \Phi}{m}\bigg),
        \ese
        where we have used that
        \bse 
            R[m\Phi^2] = 2, \qquad  F[m\Phi^2] = 0 \qand R\bigg[\frac{\l\Phi}{m}\bigg] = F\bigg[\frac{\l\Phi}{m}\bigg] = 0.
        \ese
        If we then write $g$ as a power series, we get 
        \bse 
            W_{\text{eff},\mu} = \sum_n a_n \l^n \Phi^{2+n}  m^{1-n},
        \ese 
        where the $a_n$ coefficients are potentially $\mu$ dependent. The idea is to now write this in two forms 
        \be 
        \label{eqn:WeffmuSum}
            W_{\text{eff},\mu} = m\Phi^2 \sum_n \bigg(\frac{\l\Phi}{m}\bigg)^n = \l\Phi^3 \sum_n \bigg(\frac{m}{\l\Phi}\bigg)^{1-n},
        \ee 
        where the second equality follows from simply grouping stuff.\footnote{Note that the prefactor is the other term that appears in $W$ and so has $R$-charge 2 and vanishing $F$-charge. Indeed we could have started from this one and then obtained the $m\Phi^2$ term from it.} 
        \item We now impose weak coupling limits, namely 
            \ben
                \item $\l\to 0$: This corresponds to the free theory, i.e. we require $W_{\text{eff},\mu} \to \frac{1}{2}m\Phi^2$. From the first equality in \Cref{eqn:WeffmuSum} we see this forces us to set 
                \be
                \label{eqn:ana0}
                    a_n = 0 \qquad \forall n <0, \qand a_0 = \frac{1}{2}.
                \ee 
                The $a_n$ condition is a smoothness condition and the $a_0$ the leading order normalisation.
                \item $\l\to0$ and $\frac{m}{\l}\to 0$: This corresponds to just keeping the interaction term $W_{\text{eff},\mu}\to \frac{1}{3}\l \Phi^3$, and so the second equality in \Cref{eqn:WeffmuSum} forces us to set 
                \be 
                \label{eqn:ana1}
                    a_n = 0 \qquad \forall n >1 \qand a_1 = \frac{1}{3},
                \ee
                where again the first term is a smoothness condition and the second the leading order normalisation.
            \een 
        Putting \Cref{eqn:ana0,eqn:ana1} together, we obtain 
        \bse 
            W_{\text{eff},\mu} = \frac{1}{2}m\Phi^2 + \frac{1}{3}\l \Phi^3,
        \ese 
        for \textit{arbitrary} $\mu \leq \Lambda_{\text{UV}}$, but this is exactly our UV theory, \Cref{eqn:WMicro}, and so we have \Cref{eqn:NonRenormThrm}. In other words, we have shown that there is no $\mu$ dependence in the superpotential and so it doesn't flow under RG. 
    \een
\eq 

\br 
    Note we could have just used the fact that we are considering a strict holomorphic function to restrict the two sums $n\geq 0$ and $n \leq 1$ and then fix the coefficients from there. 
\er 

\br 
    It is important to note that just because $W_{\text{eff}}$ itself is not renormalised, it does \textit{not} mean that the couplings, $m$ and $\l$, are also not renormalised. Indeed they actually \textit{are}. The reason this is consistent is because both couplings and the wavefunction all renormalise, but their renormalisations exactly cancel in $W$. That is
    \bse
        \phi \mapsto Z_{\phi} \phi, \qquad m\mapsto Z_m m \qand \l \mapsto Z_{\l}\l 
    \ese 
    but we also have a condition
    \bse 
        Z_m Z_{\phi}^2 = 1 = Z_{\l} Z_{\phi}^3 = 1,
    \ese 
    which gives 
    \bse 
        W \mapsto W. 
    \ese
\er 

\bbox 
    Show that 
    \bse 
        W = \sum_{n=1}^N \l_n \Phi^n \qand W = \frac{1}{2}m_{ij} \Phi^i \Phi^j + \frac{1}{3}\l_{ijk}\Phi^i\Phi^j\Phi^k
    \ese 
    where $i,j,k$ are indices, \textit{not} powers, are not renormalised.
\ebox 
\chapter{SUSY Gauge Theories}

We are yet to discuss the SUSY version of gauge theories. These are obviously something we want to study if we want to do SUSY versions of QED/QCD. We now want to introduce our SUSY gauge fields. 

Recall that when we introduced a general SUSY field we showed that it had too many degrees of freedom to be in an irrep of the super Poincar\'{e} algebra, so we had to put some constraints on the components. We then did this and developed our chiral superfield descriptions above by using our supercovariant derivatives. However also recall that we could reduce the number of degrees of freedom by imposing a reality condition. We now do just this and we define a real superfield $V = V^{\dagger}$. The component expansion of such a field is the following mess\footnote{The strange $\square$ terms etc will become clearer soon.} 
\bse 
    \begin{split}
        V(x,\theta,\Bar{\theta}) & = C(x) + \theta\chi(x) + \Bar{\theta}\Bar{\chi}(x) + \theta \sig^{\mu} \Bar{\theta} A_{\mu}(x) + \theta\theta M(x) + \Bar{\theta}\Bar{\theta} \Bar{M} + i\theta\theta\Bar{\theta}\big(\Bar{\l}(x) + \frac{1}{2}\Bar{\sig}^{\mu} \p_{\mu}\chi(x)\big) \\
        & \qquad - i\Bar{\theta}\Bar{\theta}\big(\l(x) - \frac{1}{2}\sig^{\mu}\p_{\mu}\Bar{\chi}(x)\big)+ \frac{1}{2}\theta\theta\Bar{\theta}\Bar{\theta} \big( D(x) - \frac{1}{2}\square C(x)\big)
    \end{split}
\ese
with all barred things related by Hermitian conjugation, e.g. $C^{\dagger} \equiv \Bar{C} = C$, $\Bar{\chi}=\chi^{\dagger}$, $A_{\mu}=A_{\mu}^{\dagger}$ etc. 

Let's look at the off-shell degrees of freedom. 
\ben 
    \item Bosons:
    \ben 
        \item C: 1
        \item A: 4
        \item M: 2
        \item D: 1 
    \een 
    \item Fermions
    \ben
        \item $\chi$: 4 
        \item $\l$: 4 
    \een
\een 
so we have $8$ d.o.f. for both Fermions and Bosons. We actually already knew this from before, where we showed that a general complex superfield had 32, so imposing a reality condition will give us 16. 

This is good, but, as we said before, the best we can do for an irrep in $\cN=1$, $4D$ SUSY (without gravity) is $8$ d.o.f. So what do we do? Well we note that our $A_{\mu}$ is a real vector with $4$ d.o.f. and we recall that gauge Bosons have less d.o.f. So we can try to make $A_{\mu}$ into a gauge field and see how that helps. Of course we can look at both abelian gauge theories and nonabelian theories, which we consider in turn. 

\section{Abelian SUSY Gauge Theories}

First let's consider the easier case of an abelian gauge theory. We make $V$ into a gauge superfield by imposing the gauge symmetry
\mybox{
    \be
    \label{eqn:VAbelianGaugeTransformation}
        V \mapsto V + \Phi + \Bar{\Phi},
    \ee
}
\noindent where the parameters of our gauge symmetry, $\Phi/\Bar{\Phi}$, are chiral/antichiral superfields. It is important to note that these are \textit{gauge} parameters, and so they are \textit{not} dynamical superfields.

How does this gauge symmetry effect our components. Well using a $\phi,\psi$ and $F$ component decomposition of the gauge parameter $\Phi$, we can easily check that
\bse 
    \begin{split}
        C & \to C + 2\Re(\phi), \\
        \chi & \to \chi + \sqrt{2}\psi, \\
        M & \to M - F, \\
        D & \to D, \\
        \l & \to \l, \\
        A_{\mu} & \to A_{\mu} - 2\p_{\mu} \Im\phi.
    \end{split}
\ese

\bbox 
    Check the last three transformation behaviours. That is prove
    \bse 
        D \to D, \qquad \l \to \l \qand A_{\mu}\to A_{\mu} - 2\p_{\mu}\Im\phi.
    \ese
    The fact that $D$ and $\l$ don't transform is why we included some weird derivative terms in the expansion of $V(x,\theta,\bar{\theta})$ above.
\ebox 

Now what do we do? Well we emphasise again that $\Phi$ is a gauge parameter and so we can \textit{choose it}. In other words we can pick the values of $\phi, \psi$ and $F$. We can therefore pick $\Re\phi, \psi$ and $F$ to gauge away (i.e. set to zero) $C, \chi$ and $M$. This leaves us with $A_{\mu}, D$ and $\l$, with an ordinary gauge symmetry for $A_{\mu}$, i.e. $A_{\mu} \to A_{\mu} - \p_{\mu}$(gauge parameter).

This is a \textit{partial}\footnote{As not fixing $\Im\phi$} \textit{gauge choice}, and it goes by the name \textit{Wess-Zumino gauge}. In this gauge the vector superfield takes a simpler form, namely 
\mybox{
    \be
    \label{eqn:VWZ}
        V_{WZ} = \theta\sig^{\mu}\Bar{\theta} A_{\mu} + i\theta\theta\Bar{\theta} \Bar{\l} - i\Bar{\theta}\Bar{\theta}\theta \l + \frac{1}{2}\theta\theta\Bar{\theta}\Bar{\theta}D.
    \ee 
}
Now taking into account gauge symmetry, we can check the degrees of freedom again. 
\ben 
    \item Bosons 
    \ben 
        \item $A_{\mu}$: $4-1=3$ 
        \item $D$: 1 
    \een 
    \item Fermions:
    \ben 
        \item $\l$: 4,
    \een 
\een    
so in total we have $4$ of each. This gives us a total of $8$ d.o.f., which is \textit{exactly} the number we need to be able to produce an irrep. This is really a SUSY version of a gauge symmetry with $A_{\mu}$ being the gauge boson, $\l$ being the \textit{gaugino} (i.e. it is the superpartner of $A_{\mu}$), and $D$ is a real auxiliary field. 

Let's make a couple comments: 
\ben[label=(\roman*)]
    \item In the WZ gauge, many computations are easier because any term that has 3 $V_{WZ}$ terms must vanish, i.e.
    \bse 
        (V_{WZ})^3 = (V_{WZ})^2D_{\a}V_{WZ} = (V_{WZ})^2\bar{D}_{\a}V_{WZ} = ... = 0.
    \ese 
    This is easily justified by a power of $\theta/\bar{\theta}$ argument. So we just need to consider \Cref{eqn:VWZ} and
    \be
    \label{eqn:VWZSquared}
        V_{WZ}^2 = \theta\sig^{\mu} \Bar{\theta} \theta\sig^{\nu}\Bar{\theta} A_{\mu}A_{\nu} = \frac{1}{2}\theta\theta\Bar{\theta}\Bar{\theta} A_{\mu}A^{\mu},
    \ee 
    where we have made use of the identity
    \bse 
        (\theta\sig^{\mu}\bar{\theta})(\theta\sig^{\nu} \bar{\theta}) = \frac{1}{2}(\theta\theta)(\bar{\theta}\bar{\theta})\eta^{\mu\nu}.
    \ese 
    So whenever we have a Taylor expansion, in WZ gauge it will truncate at quadratic order. 
    \item The WZ gauge is \textit{not} supersymmetric. That is a SUSY transformation brings us out of the WZ gauge. This is because we only partially fixed our gauge, i.e. because we haven't fixed the imaginary part of $\phi$. It then follows trivially from \Cref{eqn:VariationphipsiF} that if we start in the WZ gauge and then perform a SUSY transformation we will leave WZ gauge (we'll end up with a $\psi$ term). We therefore  need to follow this SUSY transformation up with a \textit{compensating gauge transformation}, $V\to V + \Phi + \Bar{\Phi}$ with an appropriate $\Phi$ to go back to the WZ gauge. 
\een 

\bbox 
    Compute the SUSY variation of a vector superfield $\del_{\epsilon,\Bar{\epsilon}} V$ in WZ gauge and find the compensating gauge transformation that brings you back to WZ gauge. 
\ebox  

Ok great so we have a SUSY version of a gauge field. The next thing we want to do is find the SUSY version of the field strength $F_{\mu\nu}$.  We do this by defining something called the \textit{gaugino superefilds}\footnote{Note that these are automatically chiral/antichiral by the $D^2$.}
\mybox{
    \be 
    \label{eqn:GauginoSuperfield}
        W_{\a} := -\frac{1}{4} \Bar{D}^2 D_{\a} V \qand \bar{W}_{\dot{\a}} := -\frac{1}{4} D^2 \bar{D}_{\dot{\a}} V.
    \ee 
}

\bcl 
    Our gaugino superfields are gauge invariant.\footnote{Note that this is an \textit{abelian} thing. That is we know from QCD that the field strength is only gauge \textit{covariant}, but it is \textit{not} gauge invariant itself. That is it's only $\Tr[F_{\mu\nu}F^{\mu\nu}]$ that is gauge invariant.}
\ecl 

\bq 
    A gauge transformation acts as 
    \bse
        \begin{split}
            W_{\a} & \mapsto W_{\a} - \frac{1}{4}\bar{D}^2 D_{\a} \Phi - \frac{1}{4}\bar{D}^2 D_{\a}\bar{\Phi} \\
            & = W_{\a} + \frac{1}{4}\bar{D}^{\dot{\a}}\bar{D}_{\dot{\a}} D_{\a} \Phi \\
            & - W_{\a} + \frac{1}{4}\bar{D}^{\dot{\a}}\big\{ \bar{D}_{\dot{\a}}, D_{\a}\big\} \Phi \\
            & = W_{\a} + \frac{i}{2} \sig^{\mu}_{\a\dot{\a}} \bar{D}^{\dot{\a}} \p_{\mu}\Phi \\
            & = W_{\a},
        \end{split}
    \ese 
    where we have repeatedly used that we have chiral/antichiral $\Phi/\bar{\Phi}$ and used the fact that $\bar{D}$ doesn't depend on $x$ at all so we can move it inside the $\p_{\mu}$ on the penultimate line. The proof for $\bar{W}_{\dot{\a}}$ is analogous. 
\eq 

Now, since $W_{\a}$ is gauge invariant, we can compute it in WZ gauge
\bse 
    W_{\a} = -\frac{1}{4}\bar{D}^2 D_{\a} V_{WZ}.
\ese 
In $(y,\theta,\bar{\theta})$ coordinates, we have 
\bse 
    V_{WZ} = \theta\sig^{\mu} \bar{\theta} A_{\mu}(y) + i\theta\theta\bar{\theta} \bar{\l}(y) - i \bar{\theta}\bar{\theta}\theta\l(y) + \frac{1}{2}\theta\theta\bar{\theta}\bar{\theta}\big( D(y) - i \p_{\mu} A^{\mu}(y)\big).
\ese 
Next using the relation
\bse 
    \sig^{\nu}\bar{\sig}^{\mu} = 2\sig^{\nu\mu} + 2\eta^{\nu\mu}
\ese 
it is easy to show that  
\be 
\label{eqn:DVWZExercise}
    D_{\a} V_{WZ} = \big(\sig^{\mu}\bar{\theta})_{\a} A_{\mu} + 2i\theta_{\a}\bar{\theta}\bar{\l} - i\bar{\theta}\bar{\theta}\l_{\a} + \theta_{\a} \bar{\theta}\bar{\theta}D + 2i(\sig^{\mu\nu}\theta)_{\a} \bar{\theta}\bar{\theta} \p_{\mu}A_{\nu} + \theta\theta\bar{\theta}\bar{\theta}(\sig^{\mu}\p_{\mu}\l)_{\a}.
\ee 
Then finally using
\bse 
    -\frac{1}{4}\bar{D}^2 \bar{\theta}\bar{\theta} = 1 \qand \bar{D}^2 \theta = \bar{D}^2 \bar{\theta} = 0,
\ese 
we have 
\mybox{
    \be
    \label{eqn:GauginoSuperfieldComponents}
        W_{\a}(y,\theta) = -i\l_{\a}(y) + \theta_{\a}D(y) + i(\sig^{\mu\nu}\theta)_{\a} F_{\mu\nu}(y) + \theta\theta\big(\sig^{\mu}\p_{\mu}\bar{\l}\big)_{\a}(y)
    \ee 
    with $F_{\mu\nu} = \p_{\mu}A_{\nu}-\p_{\nu}A_{\mu}$.
}

We now see where $W_{\a}$ gets the name "gaugino superfield": its bottom component is the gaugino. It is also sometimes known as the \textit{superfield strength} as it contains $F_{\mu\nu}$. 

\subsection{Supersymmetric Actions For Abelian Vector SFs}

We now want to construct the SUSY action for our abelian vector superfield.

\subsubsection{Kinetic Term}

The first thing we want to do is make our fields dynamical, i.e. we want a kinetic term. Our experience from non-SUSY QFTs tells us that we essentially want something of the form $F_{\mu\nu}F^{\mu\nu}$, which here simply becomes $W_{\a}W^{\a}$. Now as this is a chiral superfield, we integrate it over chiral superspace, giving us 
\bse 
    \int d^4 x d^2\theta W^{\a} W_{\a} + h.c.
\ese 

As we know, the integral will kill everything but the $\theta\theta$ term,
\bse 
    W^{\a}W_{\a}\big|_{\theta\theta} = -2i\l \sig^{\mu} \p_{\mu}\bar{\l} + D^2 - \frac{1}{2}(\sig^{\mu\nu})^{\a\beta}(\sig^{\rho\tau})_{\a\beta} F_{\mu\nu}F_{\rho\tau},
\ese
where we have used that the cross terms between $D$ and $F_{\mu\nu}$ turn out to vanish. Now using the relation
\bse 
    (\sig^{\mu\nu})^{\a\beta}(\sig^{\rho\tau})_{\a\beta} = \frac{1}{2}\big( \eta^{\mu\rho}\eta^{\nu\tau} - \eta^{\nu\rho} \eta^{\mu\tau}\big) - \frac{i}{2}\epsilon^{\mu\nu\rho\nu},
\ese
we get 
\bse 
    \int d^2 \theta W^{\a}W_{\a} = -\frac{1}{2} F_{\mu\nu}F^{\mu\nu} - 2i\l\sig^{\mu}\p_{\mu}\bar{\l} + D^2 + \frac{i}{4}\epsilon^{\mu\nu\rho\sig}F_{\mu\nu}F_{\rho\sig}. 
\ese 
The first three terms are real (Hermitian) while the last term is imaginary (antihermitian). We therefore introduce the \textit{complexified gauge coupling} $\tau$
\be 
\label{eqn:Tau}
    \tau = \frac{\theta}{2\pi} + \frac{4\pi i}{g^2}
\ee 
where $g$ is the gauge coupling and $\theta \sim \theta + 2\pi$ is the \textit{theta angle}.\footnote{This is the same theta angle we have mentioned a couple times in other courses, but are yet to study. It corresponds to a topological contribution to the different fields.} This then give us the Maxwell-type SUSY action
\mybox{
    \be 
    \label{eqn:SMaxwellAbelian}
        \begin{split}
            S_{\text{Maxwell}} & = \Im \bigg( \int d^4 x d^2 \theta \frac{\tau}{8\pi} W^{\a}W_{\a}\bigg) \\
            & = \int d^4 x \bigg[ \frac{1}{g^2}\bigg(-\frac{1}{4} F_{\mu\nu}F^{\mu\nu} - i\l \sig^{\mu}\p_{\mu} \bar{\l} + \frac{1}{2}D^2\bigg) + \frac{\theta}{32\pi^2}F_{\mu\nu} \tilde{F}^{\mu\nu} \bigg]
        \end{split}
    \ee 
}
\noindent where 
\bse 
    \Tilde{F}^{\mu\nu} = \frac{1}{2}\epsilon^{\mu\nu\rho\sig}F_{\rho\sig}
\ese
is that dual field strength. Let's make a couple comments:

\ben[label=(\roman*)]
    \item We have written the Maxwell SUSY action over chiral superspace so we are tempted to say that we have an $F$ term. However we have to remember that the integrand is actually chirally exact because 
    \bse
        W^{\a}W_{\a} = -\frac{1}{4}\bar{D}^2\big((D^{\a}V)W_{\a}\big),
    \ese 
    where we have used $\bar{D}^2W_{\a} = 0$ to take the $\bar{D}^2$ derivative over the whole expression. Therefore we really have a $D$-term here. 
    \item As we have presented it $\tau$ is just a parameter and so we could take it outside the integral. However with our comments on spurion analysis in mind, we might want to promote it to a background superfield, so we put it inside the integral. We then notice that it appears inside a chiral superspace integral and so the background field it corresponds to must be a chiral superfield. 
\een 

\subsubsection{Fayet-Ilioupulos Term}

We now add a new looking term, known as a \textit{Fayet-Ilioupoulos} (FI) term. It is simply given by 
\mybox{
    \be 
    \label{eqn:SFI}
        S_{FI} = -2\xi \int d^4 x d^2\theta d^2\bar{\theta} \, V = - \xi \int d^4 x \, D
    \ee 
}
\noindent where $\xi$ is a \textit{FI parameter}. The $-2$ is included here for later convenience. 

The claim this is that this is gauge invariant. The argument is essentially the same as when we introduced a Kahler potential, see \Cref{sec:MostGeneralSUSYAction}. This is a peculiarity of abelian theories. That is for non-abelian theories we will \textit{not} be able to write it down (unless we have some abelian terms). 

Note just as we argued that $\tau$ would become a chiral superfield, we see that when we promote $\xi$ to a background superfield it will be a general superfield, we will also want it to be real.

\subsection{Matter Fields}

Consider Chiral superfield $\Phi$ with charge $Q[\Phi] = q$ under a $U(1)$ gauge symmetry. This $\Phi$ is obviously not the same $\Phi$ that appears in the gauge transformation. We therefore relabel the chiral superfield in the gauge parameter as $-i\Lambda$. 

We have 
\be 
\label{eqn:AbelianMatterFieldsGaugeTransformation}
    \begin{split}
        \Phi & \to e^{iq\Lambda} \Phi \\
        \bar{\Phi} & \to  \bar{\Phi} e^{-iq\bar{\Lambda}} \\
        V & \to V + \Im\Lambda = V - \frac{i}{2}\Lambda + \frac{i}{2}\bar{\Lambda}.
    \end{split}
\ee
The bottom component of $\Lambda$ will give us a `normal' U(1). Then it's easy to write down a gauge invariant kinetic term for $\Phi$ 
\bse 
    \int d^4 x d^2\theta d^2\bar{\theta} \, \bar{\Phi} e^{2qV}\Phi,
\ese 
where the $e^{2qV}$ cancels the transformation terms. Setting $q=1$ for simplicity (can be recovered by $V\to qV$), we have in WZ gauge
\bse 
    \begin{split}
        e^{2V_{WZ}} & = 1 + 2V_{WZ} + 2V_{ZW}^2  \\
        & = 1 + 2\theta \sig^{\mu}\bar{\theta} A_{\mu} + 2i\theta\theta\bar{\theta}\bar{\l} - 2i\bar{\theta}\bar{\theta}\theta\l + \theta\theta\bar{\theta}\bar{\theta} \big(D + A_{\mu}A^{\mu}\big).
    \end{split}
\ese
Plugging this in the integrand above, we have (only keeping the top component as the rest will vanish in the integral)
\bse 
    \bar{\Phi}e^{2V}\Phi\big|_{\theta\theta\bar{\theta}\bar{\theta}} = |D_{\mu} \phi|^2 - i \bar{\psi}\bar{\sig}^{\mu}D_{\mu} \psi + |F|^2 + i\sqrt{2}\bar{\phi}\l\psi - i\sqrt{2}\phi \bar{\l} \bar{\psi} + D|\phi|^2 + \text{total deriv.}
\ese
where we have defined the familiar 
\bse
    D_{\mu} := \p_{\mu} - i A_{\mu}.
\ese 
Note the partial derivative terms come simply from the $\bar{\Phi}\Phi$ term, as we showed in \Cref{eqn:DPhiPhiExercise} earlier.

Restoring the charge $q$, we have the gauge invariant kinetic term for matter fields take the following form 
\mybox{
    \be 
    \label{eqn:SMatterAbelian}
        \begin{split}
            S_{\text{matter}} & = \int d^4x d^2\theta d^2\bar{\theta} \, \bar{\Phi}e^{2qV}\Phi \\
            & = \int d^4x \Big[ |D_{\mu} \phi|^2 - i \bar{\psi}\bar{\sig}^{\mu}D_{\mu} \psi + |F|^2 + i\sqrt{2} q \bar{\phi}\l\psi - i\sqrt{2} q \phi \bar{\l} \bar{\psi} +  q D|\phi|^2 \Big]
        \end{split}
    \ee 
}
\noindent where, by putting the $q$ back in, we now have 
\be 
\label{eqn:DmuAbelian}
    D_{\mu} := \p_{\mu} -iqA_{\mu}.
\ee 

\subsection{Abelian SUSY Gauge Theory}

We have just constructed the separate parts of our action for a single abelian vector superfield $V$. That is the most general abelian, renormalisable, gauge invariant superaction is 
\mybox{
    \be 
    \label{eqn:AbelianAction}
        S = S_{\text{Maxwell}} + S_{\text{matter}} + S_{FI} + S_W.
    \ee 
}

We can easily extend this result to an abelian vector multiplet $\{V^a \, | \, a=1,...,r\}$ with gauge group $U(1)^r$, and a chiral multiplet $\{\Phi^i\, | \, i=1,...,N\}$ with charges $Q_a[\Phi^i]=q_a^i$.  We sill have \Cref{eqn:AbelianAction} but now with 
\mybox{
    \be 
    \label{eqn:AbelianActionsPlural}
        \begin{split}
            S_{\text{Maxwell}} & = \sum_{a=1}^r \Im\bigg( \int d^4x d^2 \theta \, \frac{\tau_a}{8\pi} W^{a\a}W^a_{\a}\bigg) \\
            S_{FI} & = -2 \sum_{a=1}^r \xi_a \int d^4x d^2\theta d^2\bar{\theta} \, V^a \\
            S_{\text{matter}} & = \sum_{i=1}^N \int d^4xd^2\theta d^2\bar{\theta} \, \bar{\Phi}^i e^{2\sum_{a=1}^r q_a^i V^a} \Phi^i \\
            S_W & = \int d^4x d^2\theta \, W(\Phi^i) + h.c.
        \end{split}
    \ee 
}
\noindent where 
\bse 
    \tau_a := \frac{\theta_a}{2\pi} + \frac{4\pi i}{g_a^2} \qand W_{\a}^a = -\frac{1}{4}\bar{D}^2D_{\a} V^a,
\ese
and where $W(\Phi)$ is a gauge invariant, holomorphic function of $\Phi$, i.e. 
\bse 
    Q_a[W] = 0 \qquad \forall \,  a \in \{1,...,r\}.
\ese 

\br 
    Note that we \textit{don't} have a sum for $S_W$. This is because $W(\Phi^i)$ is a polynomial of all the $\Phi^i$s already. For example, as we will see shortly, for the SUSY version of QED we have 
    \bse 
        W = m \Phi^1 \Phi^2,
    \ese 
    with $\Phi^1$ and $\Phi^2$ being two different chiral superfields.\footnote{We will use the notation $\Phi^1=\widetilde{Q}$ and $\Phi^2 = Q$ for this later.}
\er 

Let's focus on the terms that involve the auxiliary fields, $F^i$ and $D^a$. We have 
\bse 
    \int d^4x \Bigg( \sum_i \big(\textcolor{blue}{|F_i|^2} \textcolor{red}{- \p_i W F^i -\Bar{\p}^i\Bar{W}\Bar{F}_i} \big) + \sum_a \bigg[ \textcolor{orange}{\frac{1}{2g_a^2}(D^a)^2} - \textcolor{purple}{\xi_aD^a} + \textcolor{orange}{\sum_{i=1}^N q_a^i |\phi^i|^2 D^a}  \bigg] \Bigg)
\ese 
where the colour coding tells us where the terms come from, namely: \textcolor{blue}{$S_{\text{matter}}$}, \textcolor{red}{$S_W$}, \textcolor{orange}{$S_{\text{Maxwell}}$} and \textcolor{purple}{$S_{FI}$}. 

We can then use these to find the equations of motion for our auxiliary fields simply as 
\mybox{
    \be
    \label{eqn:FDEoM}
        \Bar{F}_i = \p_i W(\phi), \qquad F_i = \bar{\p}_i \bar{W}(\bar{\phi}),  \qand D_a = - g_a^2 \bigg(\sum_{i=1}^N q_a^i |\phi^i|^2 - \xi_a\bigg)
    \ee 
}
\noindent We can rewrite this in terms of a so-called \textit{moment map} of the $a$th $U(1)$ gauge group
\be
\label{eqn:MomentMapAbelian}
    \mu_a(\phi,\bar{\phi}) := \sum_{i=1}^N q_a^i |\phi^i|^2,
\ee 
so that we simply have 
\bse
    D^a = - g_a^2 \Big(\mu_a(\phi,\Bar{\phi}) - \xi_a\Big).
\ese 

We can then obtain the scalar potential simply as
\mybox{
    \be 
    \label{eqn:ScalarPotentialAbelian}
        \begin{split}
            V(\phi,\Bar{\phi}) & = \sum_i |F^i|^2 + \sum_a \frac{1}{2g_a^2}(D^a)^2 \\
            & = \sum_i |\p_i W(\phi)|^2 + \sum_a \frac{g_a^2}{2} \big(\mu_a(\phi,\Bar{\phi})-\xi_a\big)^2
        \end{split}
    \ee
}
\noindent where the second line is understood as on-shell values (i.e. we used the EoM to get there). We now recall that if we want to have unbroken SUSY, we require that the vevs of these terms, i.e. $\la F^i \ra $ and $\la D^a \ra $, to vanish. This is just the statement that we want the lowest energy of our system to be vanishing, as otherwise we have broken SUSY. This motivates the next subsection. 

\subsection{Moduli Space Of Supersymmetric Vacua}

\bnn 
    To keep our notation short, we will omite all angular brackets, e.g. $F_i$ instead of $\la F_i \ra$. However obviously it's important that we remember they are there, as a vanishing vev is not at all the same as $F_i$ itself vanishing. 
\enn 

Just as before, we have the moduli space of SUSY vacua given by $V=0$, which is simply\footnote{Note that is the bottom component, little $\phi$, not the full chiral superfield $\Phi$.} 
\mybox{
    \be 
    \label{eqn:ModuliSpaceSUSYVacuaAbelian}
        \begin{split}
            \cM & = \{(\phi,\Bar{\phi}) \, | \, \Bar{F}_i = 0 = F_i \, \forall i \, \text{ and } D^a =  0 \, \forall a\}\big/U(1)^r \\
            & = \{(\phi,\Bar{\phi}) \, | \, \p_iW(\phi) = 0 = \bar{\p}_i\bar{W}(\bar{\phi}) \, \forall i \, \text{ and } \mu_a(\phi,\bar{\phi}) = \xi_a \, \forall a\}\big/U(1)^r
        \end{split}
    \ee 
}
\noindent where the quotient is taken to account for over counting of fields related by gauge transformations. 

\br
    Imposing the condition $D^a=0$ (i.e. $\mu_a(\phi,\Bar{\phi})=\xi_a$) and taking the $U(1)^r$ gauge symmetry quotient $(\phi^i \sim e^{i\sum_a q_a^i \a^a} \phi^i$) is called a \textit{K\"{a}hler quotient}. The name comes from the fact that the result is a K\"{a}hler manifold. 
\er 

Now, the nice thing about \Cref{eqn:ModuliSpaceSUSYVacuaAbelian} is that  the $F$-terms are holomorphic. However the the $D$ terms are real so a bit tricker to deal with. Luckily, we have the following nice theorem.

\bt 
    The space $\cM$ has a complex algebraic description as follows 
    \be 
    \label{eqn:ModuliSpaceComplexVersion}
        \cM = \{ (\phi^i) \in \C^N \, | \, \p_i W(\phi) = 0 \, \forall i\}\big/(\C^*)^r,
    \ee
    where $\C^* := \C\setminus\{0\}$ is the \textit{complexified gauge group} giving the equivalence relation 
    \bse 
        \phi^i \sim \Big(\prod_{a=1}^N \l_a^{q_a^i}\Big) \phi^i \qquad \text{with} \qquad \l_a \in \C^*.
    \ese 
\et 

This theorem tells us that instead of setting $D^a=0$ and dividing by the $U(1)^r$ gauge group, we can simply divide by the complexified gauge group $(\C^*)^r$. We can think of $\l_a$ as the bottom component of the chiral superfield gauge parameter.
\bse 
    e^{i\Lambda^a} = e^{-\Im\Lambda^a} e^{i\Re\Lambda^a}\in\C^*.
\ese 
This might seem like we made the problem more complicated, but it has the advantage that we can just work with complexified fields, namely the $F$s. 
\bex 
    SQED (SUSY QED) with $N_f$ flavours. It has $G=U(1)$ and matter fields $Q^i/\widetilde{Q}^i$\footnote{This $Q$ is not to be confused with the supercharges $Q_{\a}$.} of charges $\pm 1$ which are chiral superfields, with $i=1,...,N_f$. $Q^i$ is in the antifundamental representation of $SU(N_f)_L$ and $\widetilde{Q}^i$ in the fundamental representation of $SU(N_f)_R$.
\eex 

\bbox 
    Consider SQED with only one flavour, $N_f=1$, and $W=m\widetilde{Q}Q$.
    \ben 
        \item Write down the scalar potential 
        \item Determine the moduli space of supersymmetric vacua for 
            \ben 
                \item $m=0=\xi$ 
                \item $m=0$ but $\xi\neq0$
                \item $m\neq 0$ and $\xi=0$
                \item $m\neq0 \neq \xi$.
            \een 
        \item Determine the allowed vevs of the gauge invariant operator  $M:= \widetilde{Q}Q$ for the cases a), b) and c) above, and  show that for d) SUSY is broken.  
    \een
\ebox  

\section{Non-abelian SUSY Gauge Theories}

Ok great, we have discussed the SUSY version of abelian gauge theories, the next obvious thing to discuss is the SUSY version of \textit{non-abelian} gauge theories. As is normally the case when discussing non-ableian gauge theories, we will power through quite a lot here as the general idea is similar to the abelian case, and just highlight where the differences arise.

\subsection{Non-Abelian gauge symmetry}

Our non-abelian gauge symmetry acts as
\bse 
    \begin{split}
        \Phi & \to e^{i\Lambda} \Phi \\
        \Phi^{\dagger} = \bar{\Phi} & \to  \bar{\Phi}e^{-i\bar{\Lambda}}\\
        e^{2V} & \to e^{i\bar{\Lambda}} e^{2V} e^{-i\Lambda}\\
    \end{split}
\ese 
where $\Lambda$ is chiral superfield gauge parameter and $\bar{\Lambda}=\Lambda^{\dagger}$ is antichiral. The transformation for $e^{2V}$ comes from wanting $\bar{\Phi}e^{2V}\Phi$ to be gauge invariant.

As we have a non-abelian gauge theory, we have multiple generators, and so we can decompose our gauge parameters via 
\bse
    \Lambda = \Lambda_a T^a_R
\ese 
where $T^a_R$ are the generators in representation $R$. We will work with $R$ being the representation in which the chiral superfield $\Phi$ transforms. Just as we normally decompose $A_{\mu}$ in terms of the generators, here we decompose $V = V_aT^a_R$ so that
\bse 
    e^{2V}\Phi = e^{2V_aT^a_R} \Phi. 
\ese
In particular, if the group is $U(1)$ then all the representations are irreps, and labelled by the charges $q$. So here we can think of the generators simple as $T_q = q\b1$. 

\subsection{Gaugino Superfield}

Next we want to talk about the gaugino superfield. Now recall that it is a peculiarity of abelian theories that the field strength $F_{\mu\nu}$ be itself gauge \textit{invariant}, and that for a non-abelian theory it will only be gauge \textit{covariant}. This will obviously carry over to our gaugino superfield, which we now show. 

We have
\mybox{
    \be 
    \label{eqn:GauginoSuperfildNonAbelian}
        W_{\a} = -\frac{1}{8} \bar{D}^2\big( e^{-2V} D_{\a} e^{2V}\big) \qand \bar{W}_{\dot{\a}} = +\frac{1}{8} D^2 \big( e^{2V} \bar{D}_{\bar{\a}} e^{-2V}\big),
    \ee 
}
\noindent which is the extension of the abelian case, \Cref{eqn:GauginoSuperfield}.\footnote{Bonus exercise, check they agree to leading order. Note that $\bar{W}_{\dot{\a}}$ comes with a positive sign. The minus sign seen in \Cref{eqn:GauginoSuperfield} comes from the $-2V$ in the exponential.} 

Under a gauge transformation, we have 
\bse 
    W_{\a} \mapsto -\frac{1}{8}\bar{D}^2 \big( e^{i\Lambda} e^{-2V} e^{-i\bar{\Lambda}} D_{\a} e^{i\bar{\Lambda}} e^{2V} e^{-i\Lambda}\big)
\ese    
where the $D_{\a}$ acts on everything to it's right. We now note that $\bar{\Lambda}$ is antichiral do $D_{\a}e^{i\bar{\Lambda}}=0$ so can move it across. Similarly the $\bar{D}^2$ can be moved. This then gives us 
\bse 
    \begin{split}
        -\frac{1}{8} e^{i\Lambda} \bar{D}^2 \big( e^{-2V} D_{\a} \big(e^{2V}e^{-i\Lambda}\big)\big) & = -\frac{1}{8} e^{i\Lambda} \bar{D}^2 \Big( e^{-2V} \big[ D_{\a} \big(e^{2V}\big)e^{-i\Lambda} + e^{2V} D_{\a} e^{-i\Lambda}\big]\Big) \\
        & = -\frac{1}{8} e^{i\Lambda} \bar{D}^2 \Big( e^{-2V}D_{\a} \big(e^{2V}\big)e^{-i\Lambda} + D_{\a} e^{-i\Lambda}\Big).
    \end{split}
\ese 
Next use that $D_{\a}e^{-i\Lambda}$ is a chiral superfield, so $\bar{D}^2$ on it vanishes. We are then left with  
\bse 
    -\frac{1}{8}e^{i\Lambda}\bar{D}^2 \big(e^{2V}(D_{\a}e^{2V}) e^{-i\Lambda}\big) 
\ese 
Finally recall the Liebniz rule
\bse 
    (fg)'' = f''g + 2f'g' + fg'' 
\ese 
and again use $\bar{D}e^{-i\Lambda}=0$ so that only the first term survives, leaving us with
\bse 
    -\frac{1}{8} e^{i\Lambda} \bar{D}^2 \big(e^{-2V}D_{\a} e^{2V}\big) e^{-i\Lambda} = e^{i\Lambda} W_{\a} e^{-i\Lambda},
\ese
which is a gauge covariant result; it is the adjoint transformation.\footnote{Note that we only have $\Lambda$s no $\bar{\Lambda}$, so we really should say it transforms in the \textit{chiral} adjoint representation.}

To compute $W_{\a}$ in WZ gauge, we Taylor expand the exponential and use the properties of the WZ gauge, i.e. $V_{WZ}^3 = 0$, $V^2_{WZ}D_{\a}V_{WZ}= 0$ etc. We then have\footnote{Note that the prefactor before the commutator is \textit{not} $2$, as you might expect at first. If you work through the calculation you will get $1$.} 
\bse 
    W_{\a,WZ} = -\frac{1}{4}\bar{D}^2\big( D_{\a} V_{WZ} + [D_{\a} V_{WZ}, V_{WZ}]\big)
\ese 
When we expand this in components we get 
\mybox{
    \be 
    \label{eqn:WAlphaNonAbelian}
        W_{\a, WZ} = -i\l_{\a}(y) + \theta_{\a} D(y) + i (\sig^{\mu\nu}\theta)_{\a} F_{\mu\nu}(y) + \theta\theta \big(\sig^{\mu}D_{\mu} \bar{\l}(y)\big)_{\a}. 
    \ee 
    with 
    \bse 
        F_{\mu\nu} := \p_{\mu}A_{\nu} - \p_{\nu}A_{\mu} - i[A_{\mu},A_{\nu}] \qand D_{\mu}\bar{\l} = \p_{\mu}\bar{\l} - i[A_{\mu}, \bar{\l}]. 
    \ese
}

\subsection{SUSY Actions}

\bnn 
    In all the actions that follow, the gauge indices are implicitly contracted in order to pull out the singlet. We obviously need this for our actions to be gauge invariant. 
\enn 

So we have shown that the gaugino superfield transforms gauge covariantly, and so if we want a gauge invariant action, we take the trace (just like we do for non-SUSY theories). We then have the SUSY Yang Mills action
\mybox{
    \be 
    \label{eqn:SYM}
        \begin{split}
            S_{YM} & = \Im\bigg(\int d^4x d^2\theta \, \frac{\tau}{4\pi} \Tr(W^{\a}W_{\a}) \bigg) \\
            & = \int d^4x \, \bigg[ \frac{1}{g^2} \Tr( -\frac{1}{2}F_{\mu\nu} F^{\mu\nu} -2i\l\sig^{\mu}D_{\mu}\bar{\l} + D^2) + \frac{\theta}{16\pi^2} \Tr(F_{\mu\nu}\widetilde{F}^{\mu\nu})\bigg],
        \end{split}
    \ee 
}
\noindent where again the last term is a topological theta term. We can, of course, rewrite this in terms of things like $F_{\mu\nu}^a$ using our decomposition in terms of the generators. 

We also have our matter action given by
\mybox{
    \be 
    \label{eqn:SMatterNonAbelian}
        \begin{split}
            S_{\text{matter}} & = \int d^4 x d^2\theta d^2\bar{\theta} \, \bar{\Phi} e^{2V} \Phi \\
            & = \int d^4x \, \Big[ (D_{\mu}\phi)^{\dagger} D^{\mu}\phi - i\bar{\psi} \bar{\sig}^{\mu} D_{\mu}\psi + F^{\dagger}F + i \sqrt{2}\phi^{\dagger}\l\psi - i\sqrt{2} \bar{\psi}\bar{\l} \phi + \phi^{\dagger}D \phi\Big],
        \end{split}
    \ee
}
\noindent where we have dropped the "$+$ total deriv." on the last line.

Lastly we have
\mybox{
    \be 
    \label{eqn:SWNonAbelian}
        \begin{split}
            S_W & = \int d^4x d^2\theta W(\Phi) + h.c. \\
            & = - \int d^4x \bigg[ \p_i W(\phi) F^i + \frac{1}{2} \p_i\p_j W(\phi)\psi^i \psi^j\bigg] + h.c.,
        \end{split}
    \ee 
}
\noindent where $W(\Phi)$ is a gauge invariant polynomial of $\Phi$. 

We said lastly above because, as we have said a few times, the Fayet-Ilioupoulos action $S_{FI}$ is \textit{not} gauge invariant when we only have non-abelian terms. Therefore we do not have a $S_{FI}$ term here.

\subsection{Non-Abelian SUSY Gauge Theory}

We can then put all these actions together to give us the full non-abelian gauge invariant SUSY action 
\bse 
    S = S_{YM} + S_{\text{matter}} + S_W. 
\ese
Again we then extend this to the case of a non-abelian vector multiplet $\{V_A \, A = 1,...,r\}$, with gauge group $G = \otimes_A G_A$, and multiple chiral multiplets $\{\Phi^i\, | \, i=1,..., N\}$, with representations $R_i$ of $G$. Just as before, we get
\ben[label=(\roman*)]
    \item One $S_{YM}$ term for each simple\footnote{Simple in the sense of a simple Lie group.} factor $G_A$ of $G$. 
    \item One $S_{\text{matter}}$ for each chiral multiplet $\Phi^i$
    \item A \textit{single} $S_W$ from the $G$-invariant superpotential. 
\een 

\br 
    Note that if $G$ contains an abelian factor, then we will get a Fayet-Ilioupoulos term. 
\er 

We can again find the auxiliary field equations of motion, which turn out to simply be 
\mybox{
    \be 
    \label{eqn:FDEoMNonAbelian}
        F_i = \bar{\p_i} W^{\dagger}(\bar{\phi}), \qquad F_i^{\dagger} = \p_i W(\phi), \qand D_A^a = -g_A^2 \sum_i \phi_i^{\dagger} T^a_{A,R_i} \phi^i,
    \ee 
}
\noindent where $T^a_{A,R_i}$ are the generators for the $G_A$ in representation $R_i$. Again we can rewrite using a moment map
\bse 
    \mu_A^a(\phi^{\dagger},\phi) := \sum_i \phi_i^{\dagger} T^a_{A,R_i} \phi^i. 
\ese 
From here we obtain the scalar potential 
\mybox{
    \be 
    \label{eqn:ScalarPotentialNonAbelian}
        \begin{split}
            V(\phi^{\dagger},\phi) & = \sum_i F^{\dagger}_i F^i + \sum_A \frac{1}{2g_A^2} \sum_{a=1}^{\dim G_A} \big(D^a_A)^2 \\
            & = \sum_i \big(\p_iW(\phi)\big)^{\dagger} \big(\p^i W(\phi)\big) + \sum_A \frac{g_A^2}{2} \sum_a \big(\mu_A^a(\phi^{\dagger},\phi)\big)^2
        \end{split}
    \ee 
}
\noindent which again is just of the form $(F$-terms$)^2 + (D$-terms$)^2$. 

\subsection{Moduli Space Of SUSY Vacua}

Next we construct the moduli space of SUSY vacua as before:
\mybox{
    \be 
        \begin{split}
            \cM & = \{(\phi^{\dagger},\phi) \, | \, F_i^{\dagger} = 0 = F_i \, \forall i \, \text{ and } D_A^a = 0 \, \forall A, a\}\big/G \\
            & = \{ \phi \, | \, \p_i W(\phi) = 0 \, \forall i \}\big/G_{\C},
        \end{split}
    \ee 
}
\noindent where the second line follows from our theorem before that we can replace the $D$ constraint at the expense of complexifying the gauge quotient. 

\bex 
    Just as we gave the SQED example above, we can discuss SQCD with $N_f$ flavors. The gauge group here is $SU(N_c)$. Again we have two chiral superfields, $Q$ and $\widetilde{Q}$. We have multiple types of symmetry, which we can group into three categories 
    \ben[label=(\roman*)]
        \item Gauge symmetry: $SU(N_C)$.
        \item Global non-$R$ symmetry: $SU(N_f)_L$, $SU(N_f)_R$, $U(1)_B$ and $U(1)_A$.
        \item $R$-symmetry: $U(1)_R$.
    \een
    We have three classes of gauge invariant operators:
    \ben 
        \item Mesons: ${M_{\widetilde{i}}}^j = \widetilde{Q}^a_{\widetilde{i}} Q^j_a$. 
        \item Baryons: $B^{j_1...j_{N_C}} = \epsilon^{a_1...a_{N_C}} Q^{j_1}_{a_1} ... Q^{j_{N_C}}_{a_{N_C}}$. 
        \item AntiBaryons: $\widetilde{B}_{\widetilde{j_1}...\widetilde{j_{N_C}}} = \epsilon_{a_1...a_{N_C}} Q_{\widetilde{j_1}}^{a_1} ... Q_{\widetilde{j_{N_C}}}^{a_{N_C}}$.
    \een 
    where $a = 1, ..., N_c^2-1$ is a gauge index and and $i,\widetilde{i}=1,...,N_f$ is a flavour index. The (anti)baryons only exist is $N_f\geq N_C$. 
\eex 

\section{Minimal Supersymmetric Standard Model}

Now that we have our SUSY gauge theories, we can briefly discuss the \textit{minimal supersymmetric standard model} (MSSM). As the name suggests, this is the minimal SUSY extension of the standard model. We do not discuss it in detail here but simply explain what it consists of. 

The field content is obtained by promoting the field content of the SM to superfields as follows\footnote{The right-handed Fermions are traded for charge conjugates of left-handed ones. This is what the $c$ on the right column in the middle row means.} 

\begin{center}
	\begin{tabular}{@{} C{7cm} C{7cm}  @{}}
		\toprule
		Standard Model & Minimal SUSY Standard Model \\
		\midrule 
		Gauge fields & Vector superfields \\ \\
		Left-handed Fermions & Chiral superfields: $Q, U^c, D^c, L$ and $E^c$ \\ \\
		Higgs & \textit{Two} chiral superfields: $H_u$ and $H_d$ \\
		\bottomrule
	\end{tabular}
\end{center}

Probably the most surprising is the fact that we need \textit{two} chiral superfields for the single SM Higgs. This is because we need one to cancel so-called \textit{gauge anomolies} between Fermions (Higgsinos), and another to write Yakawa couplings from a superpotential. 

Just as we did in the SM course, we can write down a table for the different gauge group charges for these fields as follows. 

\begin{center}
	\begin{tabular}{@{} C{3cm} C{2cm} C{2cm} C{2cm} @{}}
		\toprule
		Superfield & $SU(3)_C$ & $SU(2)_L$ & $U(1)_Y$ \\
		\midrule 
		$Q_i$ & $\mathbf{3}$ & $\mathbf{2}$ & $1/6$ \\ 
		$U^c_i$ & $\mathbf{\bar{3}}$ & $\mathbf{1}$ & $-2/3$ \\ 
		$D^c_i$ & $\mathbf{\bar{3}}$ & $\mathbf{1}$ & $1/3$ \\ 
		$L_i$ & $\mathbf{1}$ & $\mathbf{2}$ & $-1/2$ \\ 
		$E^c_i$ & $\mathbf{1}$ & $\mathbf{1}$ & $1$ \\ 
		$H_u$ & $\mathbf{1}$ & $\mathbf{2}$ & $1/2$ \\ 
		$H_d$ & $\mathbf{1}$ & $\mathbf{2}$ & $-1/2$ \\ 
		\bottomrule
	\end{tabular}
\end{center}

The superpotential is a gauge invariant, renormalisable expression that preservers $R$-parity, which is given by 
\bse 
    P_R := (-1)^{3(B-L)+2s} 
\ese 
which gives $P_R=+1$ for SM particles and $P_R = -1$ for superpartners. The superpotential is explictly given by 
\bse 
    W = \mu H_u H_d + y_uH_uQU^c + y_dH_dQD^c + y_{\ell} H_dLE^c
\ese 
where the Yakawa couplings have \textit{generation indices} $(y)_{ij}$. 
\chapter{Spontaneous Symmetry Breaking}

\textcolor{red}{This material was not lectured as we ran out of time. I will try return later and update the notes to include some stuff here. For now the interested reader is directed to Sections 7,8 and 11 of Bertolini's notes.}

% -------------------------------------------------------------------
% Bibliography/Further Readings
% -------------------------------------------------------------------

\chapter*{Useful Texts \& Further Readings}

\section*{Books}
\begin{itemize}
    \item J. Wess and J. Bagger, \textit{Supersymmetry \& Supergravity}, Princeton, USA: Univ. Pr. (1992) 259 p. 
    \item P. C. West, \textit{Introduction To Supersymmetry \& Supergravity}, Singapore: World Scientific (1990) 425 p.
    \item S. Weinberg, \textit{The Quantum Theory Of Fields. Vol. 3: Supersymmetry}, Cambridge, UK: Univ. Pr. (2000) 419 p.
    \item J. Terning \textit{Modern Supersymmetry: Dynamics \& Duality}, Oxford, UK: Clarendon (2006) 324 p.
\end{itemize}

\section*{Reviews \& Lecture Notes}
\begin{itemize}
    \item A. Bilal, \href{https://arxiv.org/pdf/hep-th/0101055.pdf}{Introduction to supersymmetry}.
    \item M. F. Sohnius, \href{https://www.sciencedirect.com/science/article/abs/pii/0370157385900237?via\%3Dihub}{Introducing Supersymmetry}.
    \item S. P. Martin, \href{https://arxiv.org/pdf/hep-ph/9709356.pdf}{A Supersymmetry Primer}.
    \item J. Figueroa-O’Farrill, \href{https://arxiv.org/pdf/hep-th/0109172.pdf}{BUSSTEPP Lectures on Supersymmetry}.
    \item P. Argyres, \href{https://homepages.uc.edu/~argyrepc/cu661-gr-SUSY/index.html}{An Introduction to Global Supersymmetry}.
    \item M. J. Strassler, \href{https://arxiv.org/pdf/hep-th/0309149.pdf}{An Unorthodox Introduction to Supersymmetric Gauge Theory}.
    \item K. A. Intriligator and N. Seiberg, \href{https://arxiv.org/pdf/hep-th/9509066.pdf}{Lectures On Supersymmetric Gauge Theories \& Electric-Magnetic Duality}, Nucl. Phys. Proc. Suppl. 45BC (1996) 28.
    \item M. Bertolini, \href{https://people.sissa.it/~bertmat/susycourse.pdf}{Lectures on Supersymmetry}.
    \item R. Argurio, \href{http://homepages.ulb.ac.be/~rargurio/susycourse.pdf}{Introduction to Supersymmetry}.
\end{itemize}


%\bibliographystyle{agsm} 
%\bibliography{mybibliography} 
%\printbibliography[heading=bibintoc]


% -------------------------------------------------------------------
% Appendices
% -------------------------------------------------------------------

%\begin{appendices}
%\chapter{An Appendix of Some Kind}

\lipsum[1-3]  % Replace with your text

%\end{appendices}

\end{document}
